\documentclass[a4paper,11pt]{article}
        \usepackage[top=15mm,bottom=18mm,left=16mm,right=16mm]{geometry}
        \usepackage{fontspec}
        \usepackage{xeCJK}
        \setmainfont{Times New Roman}
        \setCJKmainfont{Yu Gothic}
        \setmonofont{Consolas}
        \setCJKmonofont{Yu Gothic}
        \usepackage{booktabs}
        \usepackage{longtable}
        \usepackage{tabularx}
        \usepackage{array}
        \usepackage{enumitem}
        \usepackage{hyperref}
        \usepackage{listings}
        \usepackage{xcolor}
        \usepackage{amsmath}
        \usepackage{amssymb}
        \hypersetup{
          colorlinks=true,
          linkcolor=blue,
          urlcolor=blue,
          pdftitle={live / sim 共通 MPC 本体},
          pdfauthor={Codex}
        }
        \lstset{
          basicstyle=\ttfamily\small,
          breaklines=true,
          columns=fullflexible,
          keepspaces=true,
          frame=single,
          framerule=0.2pt,
          rulecolor=\color{black},
          showstringspaces=false
        }
        \setlength{\parskip}{0.42em}
        \setlength{\parindent}{1em}
        \setlength{\emergencystretch}{3em}
        \renewcommand{\arraystretch}{1.15}
        \title{11. live / sim 共通 MPC 本体}
        \author{solar\_ws0129-main}
        \date{2026-07-29}
        \begin{document}
        \maketitle
        \tableofcontents
        \clearpage

        \section{対象}
        \begin{tabularx}{\linewidth}{>{\raggedright\arraybackslash}p{0.18\linewidth} X}
        \toprule
        項目 & 内容 \\
        \midrule
        ファイル & \path|mpc_solarcar/mpc_node.py| \\
        ソースSHA-256 & \path|fd107b7606b4eab77abdeb47d7d6c0fc8a151a56dba13fab2a277a538f4ed56a| \\
        種別 & Python \\
        区分 & runtime core \\
        役割 & forecast、route、vehicle telemetry、maps を使って上位速度計画と下位追従指令を出す ROS2 ノード。 \\
        起動文脈 & sim/live/live\_wifi で中心に動く単一障害点に近いノード。 \\
        \bottomrule
        \end{tabularx}

        \section{このファイルを読む理由}
        forecast、route、vehicle telemetry、maps を使って上位速度計画と下位追従指令を出す ROS2 ノード。

        \begin{itemize}[leftmargin=1.6em]
        \item SolarCarModel を直接生成する。
\item 1 Hz upper timer と lower timer 群を並列 callback group で回す。
\item calibration topic で内部係数を上書きする。
        \end{itemize}

        \section{呼び出し関係}
        \subsection{どこから呼ばれるか}
        \begin{itemize}[leftmargin=1.6em]
        \item \path|live_launch.py|
\item \path|solarcar_sim.launch.py|
        \end{itemize}

        \subsection{次に読むと理解が進むファイル}
        \begin{itemize}[leftmargin=1.6em]
        \item \path|mpc_solarcar/model.py|
\item \path|mpc_solarcar/upper_cost.py|
\item \path|mpc_solarcar/estimator.py|
        \end{itemize}

        \subsection{同一リポジトリ内の主要依存}
        \begin{itemize}[leftmargin=1.6em]
        \item \path|mpc_solarcar/estimator.py|
\item \path|mpc_solarcar/forecast_grid.py|
\item \path|mpc_solarcar/model.py|
\item \path|mpc_solarcar/path_utils.py|
\item \path|mpc_solarcar/route_utils.py|
\item \path|mpc_solarcar/schedule_utils.py|
\item \path|mpc_solarcar/signal_utils.py|
\item \path|mpc_solarcar/upper_cost.py|
\item \path|mpc_solarcar/upper_horizon.py|
\item \path|mpc_solarcar/upper_policy.py|
\item \path|mpc_solarcar/upper_solver.py|
        \end{itemize}

        \section{ソース冒頭の把握}
        モジュール docstring または先頭意図の要約:

        モジュール docstring は明示されていない。

        \subsection{代表コード断片}
        \begin{lstlisting}[language=Python]
from .upper_solver import hybrid_bounded_minimize


class MPCNode(Node):
    """
    MPC node with two modes:
      - Default: solarcar MPC (forecast-driven)
      - Passo mode: fuel-minimizing advisory MPC
    """

    def __init__(self):
        super().__init__('mpc_node')
        self.params_cfg = {}
        # A full-race upper solve can take longer than the 1 Hz control period.
        # Keep telemetry and lower control responsive while that solve runs.
        self.telemetry_callback_group = MutuallyExclusiveCallbackGroup()
        self.upper_callback_group = MutuallyExclusiveCallbackGroup()
        self.lower_callback_group = MutuallyExclusiveCallbackGroup()
        self.command_callback_group = MutuallyExclusiveCallbackGroup()
        self.declare_parameter('passo_mode', False)
        self.passo_mode = bool(self.get_parameter('passo_mode').value)
        if self.passo_mode:
\end{lstlisting}


        \section{主要構造の読み取り}
        主要クラスは MPCNode。 主要関数は z\_next\_for, quad\_penalty, cost, expand\_ctrl, build\_balance\_seed, integrate\_stationary\_duration, step\_wait, apply\_control\_stop\_at。 ROS パラメータ宣言は 151 件。 ROS I/O は publisher 14 件、subscription 19 件。

        \subsection{主要クラス・関数}
            \begin{longtable}{p{0.07\linewidth} p{0.16\linewidth} p{0.25\linewidth} p{0.40\linewidth}}
            \toprule
            行 & 種別 & 名前 & 補足 \\
            \midrule
            42 & class & \path|MPCNode| & bases: Node \\
49 & function & \path|__init__| & args: self \\
66 & function & \path|_load_stops| & args: self, stop\_yaml \\
76 & function & \path|_load_forecast_file| & args: self, path \\
102 & function & \path|_apply_params_yaml| & args: self, path \\
172 & function & \path|_maybe_reload_forecast| & args: self \\
192 & function & \path|_on_s_km_solar| & args: self, msg \\
208 & function & \path|_on_speed_solar| & args: self, msg \\
216 & function & \path|_on_soc_solar| & args: self, msg \\
225 & function & \path|_on_tb_solar| & args: self, msg \\
234 & function & \path|_on_i_solar| & args: self, msg \\
242 & function & \path|_on_v_solar| & args: self, msg \\
250 & function & \path|_measured_distance_km| & args: self \\
256 & function & \path|_sync_measured_state| & args: self \\
270 & function & \path|_on_calib_solar_gain| & args: self, msg \\
276 & function & \path|_on_calib_drive_power_gain| & args: self, msg \\
284 & function & \path|_on_calib_aux_power| & args: self, msg \\
291 & function & \path|_init_solar| & args: self \\
780 & function & \path|_current_bin_index| & args: self \\
830 & function & \path|_horizon_data| & args: self, k0 \\
872 & function & \path|_forecast_at_time| & args: self, t\_utc, s\_km, drive, control\_stop \\
948 & function & \path|_sample_plan_segments| & args: self, dt\_sample \\
959 & function & \path|_route_value| & args: self, s\_km, field, default \\
970 & function & \path|_route_average| & args: self, s0\_km, s1\_km, field, default \\
979 & function & \path|_speed_limit_at| & args: self, s\_km, default\_kmh \\
987 & function & \path|_soc_guard_speed| & args: self, v\_kmh, s\_km, d0 \\
999 & function & \path|z_next_for| & args: v\_kmh\_local \\
1045 & function & \path|_control_stop_catalog| & args: self \\
1065 & function & \path|_update_live_control_stop_hold| & args: self, s\_km, v\_kmh \\
1115 & function & \path|_dwell_penalty| & args: self, s\_km, v\_ms \\
1127 & function & \path|_mpc_solve_solar| & args: self, data \\
1164 & function & \path|quad_penalty| & args: x, cap \\
1171 & function & \path|cost| & args: v \\
1291 & function & \path|_mpc_solve_solar_distance| & args: self, t0\_utc, s0\_km, x0 \\
1379 & function & \path|expand_ctrl| & args: u\_vec \\
1382 & function & \path|build_balance_seed| &  \\
1483 & function & \path|integrate_stationary_duration| & args: t\_utc, z, Tb, s\_km, dt\_wait \\
1538 & function & \path|step_wait| & args: t\_utc, z, Tb, s\_km \\
1547 & function & \path|apply_control_stop_at| & args: t\_utc, z, Tb, s\_km \\
1565 & function & \path|cost| & args: u\_vec \\
1816 & function & \path|_publish_upper_plan| & args: self \\
1823 & function & \path|_publish_lower_plan| & args: self, v\_seq\_ms \\
1830 & function & \path|_publish_plan_path| & args: self, data \\
1858 & function & \path|_publish_metrics| & args: self, d0, v\_exec\_kmh, s\_for\_profile \\
1897 & function & \path|_publish_summary| & args: self, \_v\_exec\_kmh \\
1929 & function & \path|_interp_upper_speed| & args: self, t\_sec \\
1952 & function & \path|_distance_plan_speed| & args: self, s\_km \\
1958 & function & \path|_tau_max_for_mode| & args: self, maps, mode \\
1965 & function & \path|_tau_limits| & args: self \\
1981 & function & \path|_traction_force| & args: self, tau\_nm \\
1988 & function & \path|_pack_from_tau| & args: self, v\_ms, tau\_nm, z, Tb, env \\
2012 & function & \path|_build_lower_ref| & args: self, base\_time\_utc, s\_km, d0 \\
2039 & function & \path|_lower_rollout| & args: self, v0\_ms, u\_seq, env, z0, Tb0, tau\_drive, tau\_regen \\
2072 & function & \path|_lower_mpc_solve| & args: self, v0\_ms, s0\_km, z0, Tb0, env, v\_ref\_seq \\
2185 & function & \path|cost| & args: u\_vec \\
2273 & function & \path|_step_solar| & args: self \\
2549 & function & \path|_step_lower| & args: self \\
2600 & function & \path|_publish_lower_command_cycle| & args: self \\
2618 & function & \path|_init_passo| & args: self \\
2709 & function & \path|_on_s_km| & args: self, msg \\
2712 & function & \path|_on_speed| & args: self, msg \\
            \bottomrule
            \end{longtable}

        \subsection{ファイルを上から読んだときの定義順}
            \begin{longtable}{p{0.07\linewidth} p{0.84\linewidth}}
            \toprule
            行 & 内容 \\
            \midrule
            42 & クラス MPCNode を定義する。 \\
2965 & 関数 main を定義する。 \\
            \bottomrule
            \end{longtable}

        \subsection{import 群の詳細}
            \begin{longtable}{p{0.07\linewidth} p{0.33\linewidth} p{0.50\linewidth}}
            \toprule
            行 & import 文 & このファイルで読む理由 \\
            \midrule
            2 & \path|import math| & Python標準のスカラー数値関数と定数を使うため。実際に参照する名前と行は使用位置欄で確認する。 このファイル内での主な使用位置は L196, L199, L254, L259, L261, L264, L267, L616, ...。 \\
3 & \path|import os| & パス、環境変数、プロセス外部状態を扱うため。 このファイル内での主な使用位置は L180, L434。 \\
4 & \path|import time| & wall/monotonic time に基づく周期制御や freshness 判定を行うため。 このファイル内での主な使用位置は L175, L437, L1070, L1915, L2017, L2374, L2398, L2401, ...。 \\
5 & \path|from collections import deque| & 固定長の時系列や遅延キューを効率よく保持するため。 このファイル内での主な使用位置は L2676。 \\
6 & \path|from datetime import datetime, timezone, timedelta| & UTC 時刻や相対時間を扱うため。 このファイル内での主な使用位置は L713, L717, L718, L786, L810, L815, L824, L846, ...。 \\
8 & \path|import rclpy| & ROS 2 Python ノードとして起動・spin するため。 このファイル内での主な使用位置は L2966, L2975。 \\
9 & \path|from rclpy.callback_groups import MutuallyExclusiveCallbackGroup| & 同一ノード内 callback の排他・並行関係をCallback Groupとして指定するため。 このファイル内での主な使用位置は L54, L55, L56, L57。 \\
10 & \path|from rclpy.executors import MultiThreadedExecutor| & 複数 callback group を並列実行する executor を使うため。 このファイル内での主な使用位置は L2968。 \\
11 & \path|from rclpy.node import Node| & ROS 2 ノード本体の基底クラスとして使うため。 このファイル内での主な使用位置は L42。 \\
12 & \path|from rclpy.parameter import Parameter| & rclpy.parameter から Parameter を読み込み、このファイルの処理を組み立てるため。 このファイル内での主な使用位置は L166, L2750, L2752, L2754, L2756。 \\
14 & \path|import yaml| & profile、stop、schedule、summary YAML を読み書きするため。 このファイル内での主な使用位置は L70, L105, L2740。 \\
15 & \path|import pandas as pd| & CSV や時刻列の読込・整形・再サンプリングに使うため。 このファイル内での主な使用位置は L77, L79, L98, L100, L451, L460, L1904。 \\
16 & \path|import numpy as np| & 数値配列、clip、補間、統計計算を行うため。 このファイル内での主な使用位置は L221, L230, L265, L268, L272, L278, L542, L545, ...。 \\
18 & \path|from std_msgs.msg import Bool, Float32, Float32MultiArray, String| & Float32、Bool、Stringなどの標準ROS messageで値をpublishまたはsubscribeするため。 このファイル内での主な使用位置は L192, L208, L216, L225, L234, L242, L270, L276, ...。 \\
19 & \path|from nav_msgs.msg import Path| & 将来軌跡を dashboard や logger へ出すため。 このファイル内での主な使用位置は L741, L1833, L2701, L2880。 \\
20 & \path|from geometry_msgs.msg import PoseStamped| & Path を構成する waypoint pose を組み立てるため。 このファイル内での主な使用位置は L1842, L1851, L2886。 \\
22 & \path|from scipy.optimize import minimize| & 目的関数と制約・boundsに基づく連続数値最適化を解くため。 このファイル内での主な使用位置は L1278, L2251, L2874。 \\
24 & \path|from .model import SolarCarModel, Params| & 車体物理・電気モデル本体 から SolarCarModel, Params を読み込み、このファイルの内部処理を分担させるため。 実体ファイルは mpc\_solarcar/model.py。 このファイル内での主な使用位置は L509, L513。 \\
25 & \path|from .path_utils import resolve_path| & ROS share / 相対パス解決 から resolve\_path を読み込み、このファイルの内部処理を分担させるため。 実体ファイルは mpc\_solarcar/path\_utils.py。 このファイル内での主な使用位置は L135, L423, L424, L442, L450, L459, L470, L483, ...。 \\
26 & \path|from .route_utils import average_profile, interpolate_profile| & route\_utils.py から average\_profile, interpolate\_profile を読み込み、このファイルの内部処理を分担させるため。 実体ファイルは mpc\_solarcar/route\_utils.py。 このファイル内での主な使用位置は L963, L974, L983。 \\
27 & \path|from .schedule_utils import DriveSchedule| & schedule\_utils.py から DriveSchedule を読み込み、このファイルの内部処理を分担させるため。 実体ファイルは mpc\_solarcar/schedule\_utils.py。 このファイル内での主な使用位置は L484。 \\
28 & \path|from .estimator import BatteryMHE, MheInput, MheMeas| & Battery MHE から BatteryMHE, MheInput, MheMeas を読み込み、このファイルの内部処理を分担させるため。 実体ファイルは mpc\_solarcar/estimator.py。 このファイル内での主な使用位置は L724, L2493, L2504。 \\
29 & \path|from .forecast_grid import build_forecast_grid_payload, interp_forecast_grid| & forecast\_grid.py から build\_forecast\_grid\_payload, interp\_forecast\_grid を読み込み、このファイルの内部処理を分担させるため。 実体ファイルは mpc\_solarcar/forecast\_grid.py。 このファイル内での主な使用位置は L96, L885, L886, L887, L890, L891, L896, L898, ...。 \\
30 & \path|from .signal_utils import RobustScalarFilter, finite_float, fresh_enough, slew_limit| & signal\_utils.py から RobustScalarFilter, finite\_float, fresh\_enough, slew\_limit を読み込み、このファイルの内部処理を分担させるため。 実体ファイルは mpc\_solarcar/signal\_utils.py。 このファイル内での主な使用位置は L195, L252, L253, L258, L263, L266, L636, L643, ...。 \\
31 & \path|from .upper_cost import load_upper_cost_config, upper_stage_cost, upper_terminal_cost, quad_penalty| & 上位MPC 目的関数 から load\_upper\_cost\_config, upper\_stage\_cost, upper\_terminal\_cost, quad\_penalty を読み込み、このファイルの内部処理を分担させるため。 実体ファイルは mpc\_solarcar/upper\_cost.py。 このファイル内での主な使用位置は L590, L1227, L1231, L1236, L1238, L1242, L1250, L1251, ...。 \\
32 & \path|from .upper_horizon import build_upper_distance_horizon, plan_segment_index| & 上位MPC 距離メッシュ生成 から build\_upper\_distance\_horizon, plan\_segment\_index を読み込み、このファイルの内部処理を分担させるため。 実体ファイルは mpc\_solarcar/upper\_horizon.py。 このファイル内での主な使用位置は L1293, L1953。 \\
33 & \path|from .upper_policy import absolute_control_distances, interpolate_upper_policy, load_upper_policy_csv, shift_upper_policy_warm_start| & 上位速度計画の補間と warm start から absolute\_control\_distances, interpolate\_upper\_policy, load\_upper\_policy\_csv, shift\_upper\_policy\_warm\_start を読み込み、このファイルの内部処理を分担させるため。 実体ファイルは mpc\_solarcar/upper\_policy.py。 このファイル内での主な使用位置は L471, L1335, L1345, L1355。 \\
39 & \path|from .upper_solver import hybrid_bounded_minimize| & 上位探索ソルバ から hybrid\_bounded\_minimize を読み込み、このファイルの内部処理を分担させるため。 実体ファイルは mpc\_solarcar/upper\_solver.py。 このファイル内での主な使用位置は L1719。 \\
            \bottomrule
            \end{longtable}

        \section{このファイルを読む前に必要な基礎知識}
        次の章は、構文やROS用語を既知と仮定しないための説明である。

        \subsection{プログラム、プロセス、メモリ、オブジェクトを区別する}
ソースファイルはディスク上の文字列であり、それ自体は走っていない。Python実行可能プログラムがソースを読み、OSがその実行を一つのプロセスとして管理する。プロセスには仮想メモリ、開いているファイル、スレッド、終了コードなどが対応する。


\begin{lstlisting}
ソースファイル -> Pythonインタプリタが読み込む -> OS上のプロセス -> Pythonオブジェクトをメモリに生成 -> 関数やcallbackを実行
\end{lstlisting}


「メモリ上に生成する」とは、実行中プロセスが使う記憶領域に、その値の型、属性、参照関係を表す実体を用意することである。変数は実体そのものというより、そのオブジェクトを指す名前として理解するとPythonの代入が読みやすい。


\begin{lstlisting}[language=python]
node = MPCNode()
alias = node
# nodeとaliasは同じオブジェクトを参照する。
\end{lstlisting}


一つのプロセスには複数スレッドを持たせられる。スレッドは同じプロセスのメモリを共有するため、受信callbackと最適化callbackが同じself.zを書き換える場合は実行順と排他を検討する必要がある。

\paragraph{根拠資料}
\begin{itemize}[leftmargin=1.6em]
\item \href{https://docs.python.org/3/tutorial/classes.html}{Python公式チュートリアル: Classes}
\end{itemize}

\subsection{名前、代入、型、参照}
Pythonの代入文は、右辺を先に評価し、その結果のオブジェクトへ左辺の名前を結び付ける。`x = f()`では、まずfを呼び、その戻り値が得られてからxが更新される。


`float(...)`、`int(...)`、`bool(...)`は型名を呼び出して値を変換する。外部CSV、YAML、ROS parameterから来た値は型が期待どおりとは限らないため、このリポジトリでは明示変換が多い。


`None`は値が無いことを示す単一のオブジェクト、`math.nan`は浮動小数点値ではあるが有効な数値ではないことを示す。両者は用途が異なるため、`is None`と`math.isfinite`を使い分ける。

\paragraph{根拠資料}
\begin{itemize}[leftmargin=1.6em]
\item \href{https://docs.python.org/3/tutorial/classes.html}{Python公式チュートリアル: Classes}
\end{itemize}

\subsection{関数、引数、戻り値、スコープ}
`def`は関数本体をその場で実行する文ではない。関数オブジェクトを作り、その名前を現在の名前空間へ登録する。`f`は関数そのもの、`f()`はその関数を呼んだ結果である。


仮引数は関数定義側の受け取り名、実引数は呼出側から渡す値である。位置引数、キーワード引数、既定値付き引数、`*args`、`**kwargs`は渡し方が異なる。


\begin{lstlisting}[language=python]
def clip_speed(value, minimum=0.0, maximum=110.0):
    return min(maximum, max(minimum, value))

v = clip_speed(120.0, maximum=90.0)
\end{lstlisting}


関数内で作った通常の名前はローカルスコープに属する。メソッドが`self.z`を更新すればオブジェクトの状態が残るが、単に`z`を更新しただけならその呼出しのローカル値である。

\paragraph{根拠資料}
\begin{itemize}[leftmargin=1.6em]
\item \href{https://docs.python.org/3/tutorial/classes.html}{Python公式チュートリアル: Classes}
\end{itemize}

\subsection{module、package、import、console\_scripts}
Pythonファイルはmoduleとして読み込める。複数moduleをディレクトリにまとめたものがpackageである。`from .model import SolarCarModel`の先頭の点は、現在と同じpackage内のmodel moduleを指す相対importである。


import時にはファイルのトップレベル文が上から一度実行される。`def`や`class`は関数・クラスを登録するが、その本体の通常処理は呼び出すまで走らない。


setuptoolsの`console\_scripts`は、端末で使う実行可能名と`package.module:function`を対応付けるインストール時メタデータである。生成される実行用ラッパーはmoduleをimportし、指定関数を呼び、戻り値を終了コードとして扱う小さな入口である。


\begin{lstlisting}
ROS launchのexecutable名 -> install済みconsole script -> mpc_solarcar.mpc_nodeをimport -> main()を呼ぶ
\end{lstlisting}

\paragraph{根拠資料}
\begin{itemize}[leftmargin=1.6em]
\item \href{https://setuptools.pypa.io/en/latest/userguide/entry_point.html}{setuptools公式: Entry Points / Console Scripts}
\item \href{https://docs.python.org/3/tutorial/classes.html}{Python公式チュートリアル: Classes}
\end{itemize}

\subsection{例外、try/finally、with、資源解放}
例外は通常の戻り値とは別経路で異常を呼出元へ伝える。`try/except`は想定した異常を処理し、`finally`は成功・失敗にかかわらず後始末を行う。


`with open(...) as f:`はcontext managerを使い、ブロックを出るとファイルを閉じる。CSVやログの破損を避けるため、開いた資源の所有者と閉じる場所を明確にする。


`except Exception: pass`はノードを止めない利点がある一方、入力異常を隠して原因追跡を難しくする。安全に関係する値では、少なくとも頻度制限付き警告、異常カウンタ、fallback状態のpublishを検討する。



\subsection{class、独自クラス、インスタンス、self、継承}
クラスは、保持するデータと、そのデータを扱う関数を一つの型としてまとめる設計単位である。「独自クラス」とはPythonやROS 2が最初から用意した型ではなく、このプロジェクトが目的に合わせて定義した新しい型を指す。


`class MPCNode(Node)`は、rclpyのNodeを基底クラスとして継承し、Nodeが持つpublisher、subscription、timer、parameterなどの機能にMPC固有機能を追加する。丸括弧内は関数引数ではなく基底クラス指定である。


`MPCNode()`はクラスオブジェクトを呼び出して新しいインスタンスを作る。Pythonは新しい実体を用意し、その実体を第1引数selfとして`\_\_init\_\_`へ渡す。`self`は予約語ではなく慣習的な名前だが、この慣習を守ることでコードの意味が共有される。


\begin{lstlisting}[language=python]
class Counter:
    def __init__(self):
        self.value = 0

    def increment(self):
        self.value += 1

counter = Counter()
counter.increment()
\end{lstlisting}


`super().\_\_init\_\_(...)`は継承元の初期化処理を呼ぶ。MPCNodeではこれを省くとROSノードとして必要な内部実体が作られず、`create\_publisher`などを正しく使えない。

\paragraph{根拠資料}
\begin{itemize}[leftmargin=1.6em]
\item \href{https://docs.python.org/3/tutorial/classes.html}{Python公式チュートリアル: Classes}
\item \href{https://docs.ros.org/en/humble/Tutorials/Beginner-CLI-Tools/Understanding-ROS2-Nodes/Understanding-ROS2-Nodes.html}{ROS 2 Humble公式: Understanding nodes}
\end{itemize}

\subsection{先頭アンダースコア、\_\_init\_\_、名前付けの作法}
`self.\_step\_solar`の先頭1個のアンダースコアは「クラス内部で使う実装詳細」という慣習であり、アクセスを強制禁止する機構ではない。外部コードから呼べるが、公開APIとして依存しない意思表示である。


`\_\_init\_\_`の前後2個のアンダースコアはPythonが定めた特殊メソッド名である。クラス生成時に自動で呼ばれる。任意の名前に前後2個を付けて独自仕様を作ることは避ける。


`\_`で始まるローカル変数は、未使用または外部へ見せない意図を示す場合がある。例えば`\_base\_csv`は戻り値として受け取る必要はあるが、その後の処理では使わないことを読者へ示す。

\paragraph{根拠資料}
\begin{itemize}[leftmargin=1.6em]
\item \href{https://docs.python.org/3/tutorial/classes.html}{Python公式チュートリアル: Classes}
\end{itemize}

\subsection{list、dict、deque、NumPy配列、pandas DataFrame}
listは順序付き可変列、tupleは順序付きで通常変更しない列、dictはキーから値を引く対応表、dequeは両端追加・削除が効率的なキューである。どの構造を選ぶかはアクセス方法と更新方法で決まる。


NumPy配列は同種数値を連続的に扱い、要素ごとの演算、clip、補間、線形代数を簡潔に書く。shapeは各次元の要素数であり、速度系列なら通常`(N,)`、候補集団なら`(population, N)`となる。


pandas DataFrameは列名を持つ表である。CSV読込後は列型、欠損、単位、timezone、並び順、重複を明示的に処理しなければ、数値計算が動いても意味が誤る。



\subsection{rclpy、rcl、rmw、DDS/RTPSの層}
rclpyはPython利用者向けROS 2 client libraryである。利用者のNode、publisher、subscriptionなどをPython APIとして提供し、その下で共通C層のrclを利用する。


\begin{lstlisting}
本プロジェクトのPythonコード -> rclpy -> rcl -> rmw -> DDS/RTPS実装 -> ネットワークまたは同一PC内通信
\end{lstlisting}


rclは言語に依存しない共通ROS機能を提供するC API、rmwはROS 2と具体的middleware実装の境界である。DDS/RTPS側が探索、serialize、publish/subscribe、request/replyなどを担う。executorの実行モデルはrclだけで完結せずclient library側にも実装される。

\paragraph{根拠資料}
\begin{itemize}[leftmargin=1.6em]
\item \href{https://docs.ros.org/en/humble/Concepts/About-Internal-Interfaces.html}{ROS 2 Humble公式: Internal ROS 2 interfaces}
\end{itemize}

\subsection{ROS graph、Node、topic、service、action、parameter}
ROS graphは、実行中Nodeと、それらが持つpublisher、subscription、service、actionなどの接続関係である。Pythonクラス、プロセス、ROS graph上のNode名は関連するが同一物ではない。


topicは継続データ向けの非同期一方向publish/subscribe、serviceは短いrequest/response、actionは時間のかかる目標へfeedback、cancel、resultを持たせる。parameterはNodeごとの設定値である。


publisherが送った時点とsubscription callbackが実行される時点は同じとは限らない。通信遅延、QoS queue、executor待ちを挟むため、センサ値にはtimestampとfreshness判定が必要になる。

\paragraph{根拠資料}
\begin{itemize}[leftmargin=1.6em]
\item \href{https://docs.ros.org/en/humble/Tutorials/Beginner-CLI-Tools/Understanding-ROS2-Nodes/Understanding-ROS2-Nodes.html}{ROS 2 Humble公式: Understanding nodes}
\item \href{https://docs.ros.org/en/humble/Concepts/Basic/Interfaces-Topics-Services-Actions.html}{ROS 2 Humble公式: Topics, Services, Actions}
\item \href{https://docs.ros.org/en/humble/Concepts/Basic/About-Parameters.html}{ROS 2 Humble公式: Parameters}
\end{itemize}

\subsection{callback、timer、Executor、spin、Callback Group}
callbackは、メッセージ受信、timer満了、service requestなどのイベントが成立した後でExecutorから呼ばれる関数である。登録時に`self.\_on\_speed`と括弧なしで渡すのは、今実行せず後で呼ぶ関数オブジェクトを渡すためである。


Executorは実行可能になったcallbackを見つけ、Callback Groupの条件を確認して実行する。`spin()`は終了要求まで待機・dispatchを続けるため、通常運転中はmainの次の行へ戻らない。


MutuallyExclusiveCallbackGroupは同じgroup内のcallbackを同時実行させない。別group間はMultiThreadedExecutorで並行実行し得る。groupを分けただけでは共有属性への完全な排他にはならないため、複数groupが同じ状態を読む・書く場合は設計確認が必要である。


\begin{lstlisting}
DDS受信またはtimer満了 -> wait setでready -> Executorが選択 -> Callback Groupが許可 -> worker threadがcallback実行 -> 終了後Executorへ戻る
\end{lstlisting}

\paragraph{根拠資料}
\begin{itemize}[leftmargin=1.6em]
\item \href{https://docs.ros.org/en/rolling/Concepts/Intermediate/About-Executors.html}{ROS 2公式: Executors}
\item \href{https://docs.ros.org/en/humble/Concepts/About-Internal-Interfaces.html}{ROS 2 Humble公式: Internal ROS 2 interfaces}
\end{itemize}

\subsection{rqt\_graph、ros2 CLI、rosbag2をいつ使うか}
rqt\_graphは接続関係を見る道具であり、数値の正しさや更新周期までは保証しない。起動直後、topic名変更後、publisherが複数存在する疑いがある時に使う。


\begin{lstlisting}[language=bash]
ros2 node list
ros2 node info /mpc_node
ros2 topic list -t
ros2 topic info -v /planner/speed_cmd
ros2 topic hz /planner/speed_cmd
ros2 topic echo /planner/status
rqt_graph
\end{lstlisting}


rosbag2はtopicメッセージを時系列のまま記録・再生する。通信不具合、freshness、再計画trigger、実車とSILSの差を再現可能にするため、本番前試験では制御入力だけでなく原因となる全telemetry、status、parameter情報を記録する。


\begin{lstlisting}[language=bash]
ros2 bag record -o outputs/bags/preflight \
  /vehicle/s_km /vehicle/speed_kmh /vehicle/batt_soc \
  /vehicle/batt_temp_c /vehicle/batt_current_a /vehicle/batt_voltage_v \
  /planner/upper_speed_cmd /planner/speed_cmd /planner/status

ros2 bag info outputs/bags/preflight
ros2 bag play outputs/bags/preflight --clock
\end{lstlisting}


bag再生時はQoS互換性、simulation time、外部publisherとの二重入力に注意する。実車Nodeを同時に動かす場合はnamespaceまたはremappingで入力源を明確に分離する。

\paragraph{根拠資料}
\begin{itemize}[leftmargin=1.6em]
\item \href{https://docs.ros.org/en/ros2_packages/humble/api/rqt_graph/index.html}{ROS 2 Humble公式: rqt\_graph}
\item \href{https://docs.ros.org/en/humble/Tutorials/Beginner-CLI-Tools/Recording-And-Playing-Back-Data/Recording-And-Playing-Back-Data.html}{ROS 2 Humble公式: Recording and playing back data}
\item \href{https://docs.ros.org/en/humble/Tutorials/Beginner-CLI-Tools/Understanding-ROS2-Nodes/Understanding-ROS2-Nodes.html}{ROS 2 Humble公式: Understanding nodes}
\end{itemize}

\subsection{目的関数、制約、L-BFGS-B、SHGO、有限grid証明}
数値最適化器は、利用者が与えた目的関数を複数の候補点で評価し、より小さい値を持つ候補を探す。solverが物理を理解するのではなく、物理と運用価値はcost関数へ書かれる。


L-BFGS-Bは変数ごとの上下限を扱える局所最適化法である。初期値の近くの谷へ収束し得るため、非凸問題では複数seedや大域探索と組み合わせる。successがFalseでも有限な候補が返る場合があるため、採用条件をコード側で決める。


SHGOは定めたsamplingと局所最適化を組み合わせる大域最適化法である。有限Cartesian gridの全列挙は、そのgrid上の最良を証明できるが、連続領域全体の最良を自動的に証明しない。資料ではこの証明範囲を区別する。

\paragraph{根拠資料}
\begin{itemize}[leftmargin=1.6em]
\item \href{https://docs.scipy.org/doc/scipy/reference/generated/scipy.optimize.minimize.html}{SciPy公式: scipy.optimize.minimize}
\item \href{https://docs.scipy.org/doc/scipy/reference/generated/scipy.optimize.shgo.html}{SciPy公式: scipy.optimize.shgo}
\end{itemize}

\subsection{Cross-Entropy Methodを式と実装で理解する}
CEMは候補を生成する確率分布を持ち、良かったelite候補から分布を更新する反復的な確率最適化である。このリポジトリのupper\_solver.pyは各制御点速度を独立正規分布で生成し、上下限へclipする。


\[
u_i^{(j)} \sim \mathcal{N}(\mu_i^{(g)},(\sigma_i^{(g)})^2),\qquad u_i^{(j)}\leftarrow\operatorname{clip}(u_i^{(j)},l_i,h_i)
\]

\[
\mu_i^{(g+1)}=\frac{1}{K}\sum_{j\in\mathcal{E}_g}u_i^{(j)},\qquad \sigma_i^{(g+1)}=\max\left(\operatorname{Std}_{j\in\mathcal{E}_g}u_i^{(j)},0.05(h_i-l_i)\right)
\]

ここでE\_gはcostが小さい上位K候補である。平均は良い領域へ移り、標準偏差は探索幅を表す。標準偏差の下限は探索が完全に潰れることを避ける。


現行hybrid\_bounded\_minimizeは、deterministic seedを評価し、上位候補をL-BFGS-Bで局所refineし、設定とseed間不一致に応じてCEMを実行し、最後に再度局所refineする。したがってCEM単独ではなくhybrid solverである。


CEMで落とした候補を永久保存しないこと自体は通常の最適化として自然だが、off-nominal状態からの再利用には別のpolicy library設計が必要である。状態を無制限に全組合せ保存する代わりに、SoC、進捗、時刻、予報誤差、停止状態などのscenarioを設計し、近傍policyを検索してMPCで再最適化する。

\paragraph{根拠資料}
\begin{itemize}[leftmargin=1.6em]
\item \href{https://link.springer.com/book/10.1007/978-1-4757-4321-0}{Rubinstein and Kroese: The Cross-Entropy Method}
\item \href{https://docs.scipy.org/doc/scipy/reference/generated/scipy.optimize.minimize.html}{SciPy公式: scipy.optimize.minimize}
\end{itemize}

\subsection{warm startは何を保存し、どう効くか}
warm startは前回またはoffline探索で得た入力系列を、次の最適化の初期候補として渡すことである。答えを固定するのではなく、探索開始点と候補libraryの一員を与える。


\[
u^{0,\mathrm{new}}_i=\operatorname{interp}\left(s_i^{\mathrm{new}};\,s_j^{\mathrm{old}},u_j^{\ast,\mathrm{old}}\right)
\]

距離基準計画では、現在より後ろの旧制御点を捨て、新しい絶対距離制御点へ補間してshiftする。初期状態や天候が変わればcost評価は新条件で行われるので、warm startが不適切でも他seed、CEM、安全fallbackが補う設計が必要である。


warm startの効き方は、局所法なら収束先と反復数、CEMなら初期meanまたは候補pool、receding horizonなら前回解の時間・距離shiftとして現れる。どの位置に渡しているかを呼出引数まで追う。

\paragraph{根拠資料}
\begin{itemize}[leftmargin=1.6em]
\item \href{https://docs.scipy.org/doc/scipy/reference/generated/scipy.optimize.minimize.html}{SciPy公式: scipy.optimize.minimize}
\item \href{https://link.springer.com/book/10.1007/978-1-4757-4321-0}{Rubinstein and Kroese: The Cross-Entropy Method}
\end{itemize}

\subsection{車両力学、電力収支、効率map}
\[
F_{\mathrm{aero}}=\frac{1}{2}\rho C_dA(v+v_{\mathrm{wind}})^2,\quad F_{\mathrm{roll}}\approx mgC_{rr}\cos\theta,\quad F_{\mathrm{grade}}=mg\sin\theta
\]

車輪機械powerは概ね`P\_mech = F\_total * v + P\_inertia`で、駆動時と回生時で異なる効率mapを通してDC側powerへ変換する。mapは速度・torqueなどを軸にした実測または同定tableであり、範囲外の補間・clip規則もモデルの一部である。


\[
P_{\mathrm{pack}}=P_{\mathrm{drive,dc}}-P_{\mathrm{regen,dc}}+P_{\mathrm{aux}}-P_{\mathrm{pv}}
\]

符号規約は必ずコードで確認する。本プロジェクトでは正のpack powerを放電側としてSoCを減らす処理が中心であり、発電と回生はpack負荷を下げる方向に働く。



\subsection{SoC、内部抵抗、端子電圧、温度、MHE}
\[
V=V_{\mathrm{oc}}(z,T)-IR_{\mathrm{total}}(z,T,I),\qquad z_{k+1}=z_k-\frac{\eta(I,T)I\Delta t}{Q_{\mathrm{eff}}}
\]

SoCは直接完全には観測できないため、電流積算、OCV、端子電圧、温度、容量、効率を組み合わせる。内部抵抗はSoC、温度、電流方向で変わり、発熱と電圧制約の両方へ影響する。


MHEは有限時間窓の状態と観測誤差をまとめて最適化し、SoCや温度を推定する。古い測定や欠損値を同じ重みで使うと推定が壊れるため、timestamp freshnessと観測可用性を入口で確認する。

\paragraph{根拠資料}
\begin{itemize}[leftmargin=1.6em]
\item \href{https://docs.scipy.org/doc/scipy/reference/generated/scipy.optimize.minimize.html}{SciPy公式: scipy.optimize.minimize}
\end{itemize}

\subsection{天候、route、補間、時刻、単位}
予報は時刻だけの系列か、時刻とroute距離の2次元gridになり得る。現在時刻tと距離sがgrid点の間なら、周囲の値から補間して日射、温度、風を求める。


UTC、現地timezone、naive datetimeを混ぜると数時間のずれが発生する。入力でtimezoneを確定し、内部比較はUTC、表示は現地時刻という役割分担が安全である。


route距離のkm、速度のkm/hとm/s、時間の秒、energyのWhとJを跨ぐため、変換係数3.6、3600、1000の意味を各式で確認する。



\subsection{freshness、filter、guard、fallback、fail-safe}
分散システムでは最後に受け取った値が現在も有効とは限らない。受信時刻とtimeoutからfreshnessを判定し、stale値を計画状態へ無条件に同期しない。


filterはnoiseと一時的な飛び値を抑えるが、遅れを生む。slew limitは指令変化率を制限する。安全guardはsolverのcost罰則とは別に、現在出力へ強制制約を適用する最後の防波堤である。


fallbackは失敗時の代替動作を事前に決める設計である。前回計画保持、物理に基づく決定論的入力、停止、低速制限などから、故障modeごとに選ぶ。fallback発生はstatusとlogへ残し、正常解と区別する。



\subsection{制御、状態、入力、モデル予測制御MPC}
制御対象の内部を表す状態をx、操作入力をu、外乱・予報をwとすると、離散モデルは`x[k+1] = f(x[k], u[k], w[k])`と書ける。ソーラーカーではSoC、電池温度、距離、速度などが状態候補、速度目標や駆動トルクが入力候補になる。


\[
\min_{u_0,\ldots,u_{N-1}} \sum_{k=0}^{N-1}\ell(x_k,u_k,w_k)+V_f(x_N)
\]

MPCは現在状態からNステップ先まで予測し、目的関数と制約を満たす入力系列を求める。ただし実際に適用するのは通常先頭入力だけで、次回は新しい実測状態から再び解く。これがreceding horizonである。


予測モデル、目的関数、制約、ホライズン、solver、初期値のどれかが変わると答えも変わる。「MPCを使う」だけでは仕様は決まらず、これらを単位付きで追う必要がある。

\paragraph{根拠資料}
\begin{itemize}[leftmargin=1.6em]
\item \href{https://docs.scipy.org/doc/scipy/reference/generated/scipy.optimize.minimize.html}{SciPy公式: scipy.optimize.minimize}
\end{itemize}

\subsection{上位MPCと下位MPCの役割}
上位層は長い距離または時間を見て、エネルギー、到着、停止、速度制限、終端SoCを考えた速度計画を出す。下位層は短い周期で実測速度を見て、上位速度へ追従する駆動・回生入力を求める。


\begin{lstlisting}
予報・route・SoC -> 上位MPC -> 将来速度列 -> 下位MPC -> throttle/回生/drive mode -> driverまたはvehicle -> 新しいtelemetry
\end{lstlisting}


上位解が遅い間も下位出力を止めないこと、古い計画を安全に保持すること、上位と下位で単位と時刻基準を一致させることが実装上重要である。



\subsection{launch Action、Node Action、実行可能名、remapping}
`launch\_ros.actions.Node(...)`はrclpyのNode基底クラスではなく、指定したpackageのexecutableをプロセスとして起動するlaunch Actionである。


`DeclareLaunchArgument`はlaunch実行時に受け取る入力欄を宣言し、`LaunchConfiguration`はその値を後で解決するsubstitutionを表す。`perform(context)`は実行時contextから確定文字列を取り出す。


launchの`name`はNode名override、`parameters`は起動時parameter、`remappings`はNodeやtopicの既定名を別名へ対応付ける。launchはPythonファイルからmainという名前を推測せず、executableとしてインストールされたconsole scriptを起動する。

\paragraph{根拠資料}
\begin{itemize}[leftmargin=1.6em]
\item \href{https://docs.ros.org/en/humble/Tutorials/Beginner-CLI-Tools/Understanding-ROS2-Nodes/Understanding-ROS2-Nodes.html}{ROS 2 Humble公式: Understanding nodes}
\item \href{https://setuptools.pypa.io/en/latest/userguide/entry_point.html}{setuptools公式: Entry Points / Console Scripts}
\end{itemize}

\subsection{単体試験、SILS、replay、観測可能性}
単体試験は関数やモデル項の局所契約、SILSは複数Nodeと時間進行を含む閉ループ、historical replayは実測入力への再現性、本番preflightは実機通信と停止動作を確認する。目的が異なるため一つで代用しない。


再現可能な診断には、入力、出力、内部状態、solver成否、候補数、計算時間、fallback、source revision、profile hashを同じrun IDで記録する。


非同期不具合は最終値だけでは追えない。topic周期、message age、callback所要時間、queue遅延、Node生存、publisher数を同時に観測する。

\paragraph{根拠資料}
\begin{itemize}[leftmargin=1.6em]
\item \href{https://docs.ros.org/en/ros2_packages/humble/api/rqt_graph/index.html}{ROS 2 Humble公式: rqt\_graph}
\item \href{https://docs.ros.org/en/humble/Tutorials/Beginner-CLI-Tools/Recording-And-Playing-Back-Data/Recording-And-Playing-Back-Data.html}{ROS 2 Humble公式: Recording and playing back data}
\item \href{https://docs.ros.org/en/rolling/Concepts/Intermediate/About-Executors.html}{ROS 2公式: Executors}
\end{itemize}

        \section{mpc\_node.py専用の統合解説}
起動機構とMPC内部計算を一つの時間軸で接続する。

\subsection{起動経路を大本から一本につなぐ}
\begin{lstlisting}
SolarSim.ps1 -> scripts/solar_control.sh -> ROS 2 launch -> launch_ros.actions.Node -> setup.pyのconsole_scripts -> mpc_solarcar.mpc_node:main -> rclpy.init -> MPCNode() -> __init__ -> _init_solar -> Executor.add_node -> spin
\end{lstlisting}


この対応は推測ではない。setup.pyのconsole\_scriptsには`mpc\_node = mpc\_solarcar.mpc\_node:main`が登録され、live\_launch.pyとsolarcar\_sim.launch.pyはexecutableとしてmpc\_nodeを指定する。launchがmainという名前を探索するのではない。


launch側のNodeはプロセス起動指示、rclpy.node.Nodeは基底クラス、MPCNodeはこのリポジトリが定義した派生クラス、`node = MPCNode()`のnodeは実行中メモリに作られたインスタンス、`/mpc\_node`はROS graph上の名前である。



\subsection{mainとspinの前後で実行方式が変わる}
\begin{lstlisting}[language=python]
def main():
    rclpy.init()
    node = MPCNode()
    executor = MultiThreadedExecutor(num_threads=4)
    executor.add_node(node)
    try:
        executor.spin()
    finally:
        executor.shutdown()
        node.destroy_node()
        rclpy.shutdown()
\end{lstlisting}


spinまでは通常のPythonとして上から一度だけ進む。spin後はExecutorのevent loopへ入り、timerと受信がreadyになった時にcallbackが呼ばれる。callback終了後はmainの次行ではなくExecutorへ戻る。


finallyはCtrl+C、launch停止、例外などでspinを抜けた後の後始末である。Executor、Node、rclpy contextの順に明示終了する。



\subsection{MPCNode.\_\_init\_\_とselfの正確な意味}
`MPCNode()`を評価すると、新しいインスタンスが作られ、その参照がselfとして`MPCNode.\_\_init\_\_`へ入る。`super().\_\_init\_\_('mpc\_node')`が基底Nodeを初期化して初めてROS Node機能を持つ。


`self.z`、`self.Tb`、`self.model`は同じMPCNodeインスタンスに保存される状態である。callbackが別時刻に呼ばれても同じ属性を参照する。`z`のような局所変数はその関数呼出し内だけである。


先頭アンダースコア付きメソッドは内部実装という慣習、`\_\_init\_\_`はPython特殊メソッドである。`@Class`という一般構文はなく、このファイルのMPCNode自体にはdataclass decoratorは使われていない。



\subsection{四つのCallback Groupと共有状態}
telemetry、upper、lower、commandの四つを別MutuallyExclusiveCallbackGroupへ置く。同じgroup内は同時実行しないが、別groupは四workerで並行実行し得る。長い上位solve中にもtelemetry受信と下位出力を続ける意図である。


別groupはself.z、self.Tb、self.v\_now、self.last\_dataなどを共有する。コードは上位solve後に`\_sync\_measured\_state()`を再実行して長時間solve中に受けた最新値を反映するが、明示lockですべてを原子的snapshot化しているわけではない。


PythonのGILは複数文にまたがる状態整合性を保証しない。診断時はcallback開始終了時刻、plan ID、telemetry timestampを併記し、どの状態snapshotで解いたか確認する。



\subsection{時間基準上位MPC}
`\_mpc\_solve\_solar(data)`は将来速度列を最適化変数とし、各予測stepで慣性power、electrical balance、SoC、温度、距離、制約penaltyを順に計算する。


\[
P_{\mathrm{inertia},k}=\frac{\frac12m(v_k^2-v_{k-1}^2)}{\Delta t},\quad s_{k+1}=s_k+\frac{v_k\Delta t}{1000}
\]

`model.electrical\_balance`が電流、電圧、pack power、損失を返し、`model.soc\_step`がSoCを進める。温度更新とcost蓄積は呼出側にもある。最後にSciPy L-BFGS-Bへcost、初期値、boundsを渡す。



\subsection{距離基準上位MPCとCEM}
`\_mpc\_solve\_solar\_distance`はrouteを距離区間へ分け、制御点速度uを区間速度へ線形補間する。停止座標をmesh edgeへ追加し、到着、dwell発電、再出発を同じ座標で評価する。


初期候補は、前回planの距離shift、offline initial policyの補間、現在速度一定の順で選ぶ。別途balance seedを作り、`hybrid\_bounded\_minimize`へ渡す。


solverはseed評価、上位seedのL-BFGS-B、条件付きCEM、再L-BFGS-Bを行う。CEMの各世代ではmean候補を1個必ず含め、残りを正規乱数で生成し、elite平均と標準偏差で次世代分布を更新する。


finite grid全列挙の証明は宣言grid上に限る。CEMや局所候補を含む連続領域の大域最適性を意味しないため、solve\_infoの`discrete\_global\_proof`、`finite\_library\_global\_proof`、`certificate\_scope`を区別する。



\subsection{warm startとoff-nominal状態}
前回速度列を新しい絶対距離制御点へ補間するため、予定より遅い・速い場合でも過去区間を除いた残り計画を初期値にできる。これはplanを固定せず、現在SoC、温度、天候でcostを再評価する。


offline CEMで理想軌道だけを残すと大きな逸脱時のseed品質は落ちる。対策は全連続状態の全組合せ保存ではなく、SoC偏差、進捗偏差、時刻、予報scale、停止継続、温度などのscenario libraryを設計し、近傍policyとphysics seedを併用してlive MPCで修正することである。


この現行ノードは前回plan、initial\_upper\_policy、balance seed、generic seed、CEMを併用するが、多次元scenario library検索を完全実装したものではない。この境界を資料上で保証と将来拡張に分ける。



\subsection{下位MPCと指令継続}
`\_build\_lower\_ref`が上位planから短期参照速度列を作り、`\_lower\_mpc\_solve`が駆動・回生入力uを求める。逆動力学seedを必ず作り、設定した場合だけL-BFGS-Bでrefineする。


上位solve中はoptionにより下位refineを省略し、決定論的入力を使う。`\_publish\_lower\_command\_cycle`は独立timerから保存済み指令を出すため、optimizer完了を待たず出力を継続する。


ただしdocstringは`1 Hz output path`と書く一方、実際のtimer周期は`1/lower\_rate\_hz`で、既定lower\_rate\_hzが5なら5 Hzである。説明では実装値を正とし、この不一致を既知の文書上問題として記す。



\subsection{実測同期、freshness、MHE、fallback}
各telemetry callbackは受信時刻とfilter済み値を保存する。`\_sync\_measured\_state`はtimeout内の値だけを状態へ反映し、古い速度はNaNに戻す。距離には大きな後退値を捨てるguardがある。


MHE有効時は観測可能なSoC、温度、電流、電圧を窓へpushして状態を推定する。無効時は車両モデルで状態を進めるが、新鮮なSoCまたは温度実測がある項目はモデル上書きを避ける。


solver失敗時はwarm-start planまたは決定論的入力へfallbackする。さらにschedule、速度制限、SoC guard、control stop holdをoptimizer後段で強制するため、soft penaltyだけに安全を依存しない。



\subsection{実機で使う診断順}
\begin{lstlisting}[language=bash]
ros2 node list
ros2 node info /mpc_node
ros2 topic info -v /vehicle/speed_kmh
ros2 topic info -v /planner/speed_cmd
ros2 topic hz /vehicle/speed_kmh
ros2 topic hz /planner/speed_cmd
ros2 topic echo /planner/status
rqt_graph
\end{lstlisting}


Nodeが無い場合はlaunch・console script・process、topicが無い場合はpublisher、周期が遅い場合は通信またはExecutor、値が古い場合はtimestamp/freshness、指令が0ならschedule/stop/SoC guard、solver失敗ならstatusとlog、という順で範囲を狭める。


空転試験では距離を進めない入力sourceを使い、実GPSとreplay publisherを同時に接続しない。bagへtelemetry、upper/lower指令、statusを同時記録し、停止後に同じprofileとsource revisionで再生する。



\subsection{現行改修の設計意図と残る注意点}
現行コードのコメントと差分から、長時間full-race solve中の応答維持、2次元forecast補間、実測freshness、runtime profile override、offline policy warm start、control stopの正確なmesh分割、決定論的lower fallback、solve時間logを目的とした改修を確認できる。


一方、異なるCallback Group間の共有状態snapshot、広い`except Exception`、commandの二経路publish、1 Hzというdocstringと実timerの不一致は、利用者が挙動を理解するうえで注意が必要である。これらは機能説明と保証範囲を分けて記載する。


Git index上ではmpc\_node.pyが競合未解決状態であるため、本資料は2026-07-29時点のワークツリー内容を根拠とし、merge完了後はmanifestのsource hashを更新して再生成する必要がある。



        \section{関数・クラスを上から順に解説}
        \subsection{L42 クラス MPCNode}
            \begin{itemize}[leftmargin=1.6em]
              \item 行範囲: L42--L2962
              \item 所属: module直下
              \item 定義: \path|MPCNode(bases=Node)|
              \item docstring: MPC node with two modes:
  - Default: solarcar MPC (forecast-driven)
  - Passo mode: fuel-minimizing advisory MPC
              \item このブロックが直接呼ぶ主な関数・メソッド: 特に明示されていない。
              \item 戻り値の要点: 戻り値が無いか、return 文が明示されていない。
              \item 主なローカル代入名: 特になし。
              \item 読み取る主なインスタンス属性: 特になし。
              \item 更新する主なインスタンス属性: 特になし。
              \item 明示的に送出する例外: 明示的raiseはない。
              \item 制御構造の規模: 条件分岐 0、ループ 0、try 0
            \end{itemize}
            \paragraph{この定義を読むためのPython構文}
            \begin{itemize}[leftmargin=1.6em]
            \item class文は新しい型を定義する。丸括弧内は継承する基底クラスである。
\item lambdaは名前を付けずに短い関数オブジェクトを作る構文である。
\item 内包表記は、反復・条件・値生成を一つの式で記述してcollectionを作る。
\item with文はcontext managerに開始・終了処理を任せ、ファイルなどの資源を確実に解放する。
\item try文は例外が起き得る処理と、異常時または終了時の経路を分ける。
            \end{itemize}
            \paragraph{上から順の処理}
            \begin{enumerate}[leftmargin=1.8em]
            \item docstring または説明文字列を置く。
\item 関数 \_\_init\_\_ を定義する。
\item 関数 \_load\_stops を定義する。
\item 関数 \_load\_forecast\_file を定義する。
\item 関数 \_apply\_params\_yaml を定義する。
\item 関数 \_maybe\_reload\_forecast を定義する。
\item 関数 \_on\_s\_km\_solar を定義する。
\item 関数 \_on\_speed\_solar を定義する。
\item 関数 \_on\_soc\_solar を定義する。
\item 関数 \_on\_tb\_solar を定義する。
\item 関数 \_on\_i\_solar を定義する。
\item 関数 \_on\_v\_solar を定義する。
            \end{enumerate}
            \paragraph{代表コード断片}
            \begin{lstlisting}[language=Python]
class MPCNode(Node):
    """
    MPC node with two modes:
      - Default: solarcar MPC (forecast-driven)
      - Passo mode: fuel-minimizing advisory MPC
    """

    def __init__(self):
        super().__init__('mpc_node')
        self.params_cfg = {}
        # A full-race upper solve can take longer than the 1 Hz control period.
        # Keep telemetry and lower control responsive while that solve runs.
        self.telemetry_callback_group = MutuallyExclusiveCallbackGroup()
        self.upper_callback_group = MutuallyExclusiveCallbackGroup()
        self.lower_callback_group = MutuallyExclusiveCallbackGroup()
        self.command_callback_group = MutuallyExclusiveCallbackGroup()
        self.declare_parameter('passo_mode', False)
        self.passo_mode = bool(self.get_parameter('passo_mode').value)
        if self.passo_mode:
            self._init_passo()
        else:
            self._init_solar()

    # -------------------- common helpers --------------------
    def _load_stops(self, stop_yaml: str):
        self.stops = []
        try:
            with open(stop_yaml, 'r', encoding='utf-8') as f:
                y = yaml.safe_load(f) or {}
                self.stops = y.get('stops', [])
                self.get_logger().info(f'Loaded {len(self.stops)} stop points from {stop_yaml}')
        except Exception:
            self.get_logger().info('No stop_points.yaml provided. Running without dwell constraints.')

    def _load_forecast_file(self, path: str):
...
\end{lstlisting}

\subsection{L49 関数 MPCNode.\_\_init\_\_}
            \begin{itemize}[leftmargin=1.6em]
              \item 行範囲: L49--L63
              \item 所属: \path|MPCNode|
              \item 定義: \path|__init__(self)|
              \item docstring: 明示されていない。
              \item このブロックが直接呼ぶ主な関数・メソッド: MutuallyExclusiveCallbackGroup, \_\_init\_\_, \_init\_passo, \_init\_solar, bool, declare\_parameter, get\_parameter, super
              \item 戻り値の要点: 戻り値が無いか、return 文が明示されていない。
              \item 主なローカル代入名: 特になし。
              \item 読み取る主なインスタンス属性: self.\_init\_passo, self.\_init\_solar, self.declare\_parameter, self.get\_parameter, self.passo\_mode
              \item 更新する主なインスタンス属性: self.command\_callback\_group, self.lower\_callback\_group, self.params\_cfg, self.passo\_mode, self.telemetry\_callback\_group, self.upper\_callback\_group
              \item 明示的に送出する例外: 明示的raiseはない。
              \item 制御構造の規模: 条件分岐 1、ループ 0、try 0
            \end{itemize}
            \paragraph{この定義を読むためのPython構文}
            \begin{itemize}[leftmargin=1.6em]
            \item 第1引数selfは、このメソッドを呼んだインスタンス自身を受け取る慣習名である。
            \end{itemize}
            \paragraph{上から順の処理}
            \begin{enumerate}[leftmargin=1.8em]
            \item super().\_\_init\_\_(...) を実行する。
\item self.params\_cfg に \{\} の結果を代入する。
\item self.telemetry\_callback\_group に MutuallyExclusiveCallbackGroup() の結果を代入する。
\item self.upper\_callback\_group に MutuallyExclusiveCallbackGroup() の結果を代入する。
\item self.lower\_callback\_group に MutuallyExclusiveCallbackGroup() の結果を代入する。
\item self.command\_callback\_group に MutuallyExclusiveCallbackGroup() の結果を代入する。
\item self.declare\_parameter(...) を実行する。
\item self.passo\_mode に bool(self.get\_parameter('passo\_mode').value) の結果を代入する。
\item 条件 self.passo\_mode を判定し、真なら内部処理を行う。
\item   self.\_init\_passo(...) を実行する。
\item   上の条件が偽の場合:
\item   self.\_init\_solar(...) を実行する。
            \end{enumerate}
            \paragraph{代表コード断片}
            \begin{lstlisting}[language=Python]
    def __init__(self):
        super().__init__('mpc_node')
        self.params_cfg = {}
        # A full-race upper solve can take longer than the 1 Hz control period.
        # Keep telemetry and lower control responsive while that solve runs.
        self.telemetry_callback_group = MutuallyExclusiveCallbackGroup()
        self.upper_callback_group = MutuallyExclusiveCallbackGroup()
        self.lower_callback_group = MutuallyExclusiveCallbackGroup()
        self.command_callback_group = MutuallyExclusiveCallbackGroup()
        self.declare_parameter('passo_mode', False)
        self.passo_mode = bool(self.get_parameter('passo_mode').value)
        if self.passo_mode:
            self._init_passo()
        else:
            self._init_solar()
\end{lstlisting}

\subsection{L66 関数 MPCNode.\_load\_stops}
            \begin{itemize}[leftmargin=1.6em]
              \item 行範囲: L66--L74
              \item 所属: \path|MPCNode|
              \item 定義: \path|_load_stops(self, stop_yaml: str)|
              \item docstring: 明示されていない。
              \item このブロックが直接呼ぶ主な関数・メソッド: get, get\_logger, info, len, open, safe\_load
              \item 戻り値の要点: 戻り値が無いか、return 文が明示されていない。
              \item 主なローカル代入名: f, y
              \item 読み取る主なインスタンス属性: self.get\_logger, self.stops
              \item 更新する主なインスタンス属性: self.stops
              \item 明示的に送出する例外: 明示的raiseはない。
              \item 制御構造の規模: 条件分岐 0、ループ 0、try 1
            \end{itemize}
            \paragraph{この定義を読むためのPython構文}
            \begin{itemize}[leftmargin=1.6em]
            \item 第1引数selfは、このメソッドを呼んだインスタンス自身を受け取る慣習名である。
\item with文はcontext managerに開始・終了処理を任せ、ファイルなどの資源を確実に解放する。
\item try文は例外が起き得る処理と、異常時または終了時の経路を分ける。
            \end{itemize}
            \paragraph{上から順の処理}
            \begin{enumerate}[leftmargin=1.8em]
            \item self.stops に [] の結果を代入する。
\item 例外処理を伴う try ブロックを実行する。
\item   with 文で open(stop\_yaml, 'r', encoding='utf-8') を管理しながら処理する。
\item     y に yaml.safe\_load(f) or \{\} の結果を代入する。
\item     self.stops に y.get('stops', []) の結果を代入する。
\item     self.get\_logger().info(...) を実行する。
\item   Exceptionを捕捉した場合:
\item   self.get\_logger().info(...) を実行する。
            \end{enumerate}
            \paragraph{代表コード断片}
            \begin{lstlisting}[language=Python]
    def _load_stops(self, stop_yaml: str):
        self.stops = []
        try:
            with open(stop_yaml, 'r', encoding='utf-8') as f:
                y = yaml.safe_load(f) or {}
                self.stops = y.get('stops', [])
                self.get_logger().info(f'Loaded {len(self.stops)} stop points from {stop_yaml}')
        except Exception:
            self.get_logger().info('No stop_points.yaml provided. Running without dwell constraints.')
\end{lstlisting}

\subsection{L76 関数 MPCNode.\_load\_forecast\_file}
            \begin{itemize}[leftmargin=1.6em]
              \item 行範囲: L76--L100
              \item 所属: \path|MPCNode|
              \item 定義: \path|_load_forecast_file(self, path: str)|
              \item docstring: 明示されていない。
              \item このブロックが直接呼ぶ主な関数・メソッド: Index, all, build\_forecast\_grid\_payload, drop\_duplicates, dropna, get\_logger, getattr, isna, read\_csv, sort\_values, str, to\_datetime
              \item 戻り値の要点: 戻り値が無いか、return 文が明示されていない。
              \item 主なローカル代入名: t, tzname
              \item 読み取る主なインスタンス属性: self.df, self.get\_logger
              \item 更新する主なインスタンス属性: self.df, self.forecast\_grid, self.forecast\_times
              \item 明示的に送出する例外: 明示的raiseはない。
              \item 制御構造の規模: 条件分岐 5、ループ 0、try 1
            \end{itemize}
            \paragraph{この定義を読むためのPython構文}
            \begin{itemize}[leftmargin=1.6em]
            \item 第1引数selfは、このメソッドを呼んだインスタンス自身を受け取る慣習名である。
\item try文は例外が起き得る処理と、異常時または終了時の経路を分ける。
            \end{itemize}
            \paragraph{上から順の処理}
            \begin{enumerate}[leftmargin=1.8em]
            \item self.df に pd.read\_csv(path) の結果を代入する。
\item 条件 'time' in self.df.columns を判定し、真なら内部処理を行う。
\item   t に pd.to\_datetime(self.df['time'], format='mixed', errors='coerce') の結果を代入する。
\item   tzname に str(getattr(self, 'forecast\_time\_tz', 'UTC') or 'UTC') の結果を代入する。
\item   条件 t.dt.tz is None を判定し、真なら内部処理を行う。
\item     条件 tzname.upper() == 'UTC' を判定し、真なら内部処理を行う。
\item       t に t.dt.tz\_localize('UTC') の結果を代入する。
\item       上の条件が偽の場合:
\item       例外処理を伴う try ブロックを実行する。
\item     上の条件が偽の場合:
\item     t に t.dt.tz\_convert('UTC') の結果を代入する。
\item   self.df['time'] に t の結果を代入する。
\item   条件 self.df['time'].isna().all() を判定し、真なら内部処理を行う。
\item     self.get\_logger().warn(...) を実行する。
\item   上の条件が偽の場合:
\item   self.get\_logger().warn(...) を実行する。
\item self.forecast\_grid に build\_forecast\_grid\_payload(self.df) の結果を代入する。
\item 条件 'time' in self.df.columns and (not self.df['time'].isna().all()) を判定し、真なら内部処理を行う。
\item   self.forecast\_times に pd.Index(self.df['time'].dropna().drop\_duplicates().sort\_values()) の結果を代入する。
\item   上の条件が偽の場合:
\item   self.forecast\_times に pd.Index([]) の結果を代入する。
            \end{enumerate}
            \paragraph{代表コード断片}
            \begin{lstlisting}[language=Python]
    def _load_forecast_file(self, path: str):
        self.df = pd.read_csv(path)
        if 'time' in self.df.columns:
            t = pd.to_datetime(self.df['time'], format='mixed', errors='coerce')
            tzname = str(getattr(self, 'forecast_time_tz', 'UTC') or 'UTC')
            if t.dt.tz is None:
                if tzname.upper() == 'UTC':
                    t = t.dt.tz_localize('UTC')
                else:
                    try:
                        t = t.dt.tz_localize(tzname, ambiguous='NaT', nonexistent='NaT').dt.tz_convert('UTC')
                    except Exception:
                        t = t.dt.tz_localize('UTC')
            else:
                t = t.dt.tz_convert('UTC')
            self.df['time'] = t
            if self.df['time'].isna().all():
                self.get_logger().warn("forecast 'time' column could not be parsed; falling back to index bins.")
        else:
            self.get_logger().warn("forecast CSV has no 'time' column; falling back to index bins.")
        self.forecast_grid = build_forecast_grid_payload(self.df)
        if 'time' in self.df.columns and not self.df['time'].isna().all():
            self.forecast_times = pd.Index(self.df['time'].dropna().drop_duplicates().sort_values())
        else:
            self.forecast_times = pd.Index([])
\end{lstlisting}

\subsection{L102 関数 MPCNode.\_apply\_params\_yaml}
            \begin{itemize}[leftmargin=1.6em]
              \item 行範囲: L102--L170
              \item 所属: \path|MPCNode|
              \item 定義: \path|_apply_params_yaml(self, path: str)|
              \item docstring: 明示されていない。
              \item このブロックが直接呼ぶ主な関数・メソッド: Parameter, append, dict, float, get, get\_logger, get\_parameter, has\_parameter, hasattr, info, isinstance, items
              \item 戻り値の要点: 戻り値が無いか、return 文が明示されていない。
              \item 主なローカル代入名: cfg, f, k, key, model\_cfg, motor\_type, mpc\_cfg, mt, params, runtime\_mode, runtime\_overrides, runtime\_section, v, v\_str, val
              \item 読み取る主なインスタンス属性: self.get\_logger, self.get\_parameter, self.has\_parameter, self.model, self.set\_parameters
              \item 更新する主なインスタンス属性: self.params\_cfg
              \item 明示的に送出する例外: 明示的raiseはない。
              \item 制御構造の規模: 条件分岐 15、ループ 2、try 6
            \end{itemize}
            \paragraph{この定義を読むためのPython構文}
            \begin{itemize}[leftmargin=1.6em]
            \item 第1引数selfは、このメソッドを呼んだインスタンス自身を受け取る慣習名である。
\item with文はcontext managerに開始・終了処理を任せ、ファイルなどの資源を確実に解放する。
\item try文は例外が起き得る処理と、異常時または終了時の経路を分ける。
            \end{itemize}
            \paragraph{上から順の処理}
            \begin{enumerate}[leftmargin=1.8em]
            \item 例外処理を伴う try ブロックを実行する。
\item   with 文で open(path, 'r', encoding='utf-8') を管理しながら処理する。
\item     cfg に yaml.safe\_load(f) or \{\} の結果を代入する。
\item   Exceptionを捕捉した場合:
\item   self.get\_logger().warn(...) を実行する。
\item    を返す。
\item self.params\_cfg に cfg if isinstance(cfg, dict) else \{\} の結果を代入する。
\item model\_cfg に cfg.get('model', cfg) if isinstance(cfg, dict) else \{\} の結果を代入する。
\item 条件 isinstance(model\_cfg, dict) を判定し、真なら内部処理を行う。
\item   motor\_type に None の結果を代入する。
\item   model\_cfg.items() を順に走査し、各要素を (k, v) に入れて処理する。
\item     条件 hasattr(self.model.p, k) を判定し、真なら内部処理を行う。
\item       例外処理を伴う try ブロックを実行する。
\item     条件 k == 'motor\_type' を判定し、真なら内部処理を行う。
\item       motor\_type に str(v) の結果を代入する。
\item     条件 k == 'drive\_mode' を判定し、真なら内部処理を行う。
\item       self.model.drive\_mode に str(v) の結果を代入する。
\item     条件 k == 'drive\_mode\_tau\_margin' を判定し、真なら内部処理を行う。
\item       例外処理を伴う try ブロックを実行する。
\item     条件 k == 'ocv\_soc\_map' を判定し、真なら内部処理を行う。
\item       例外処理を伴う try ブロックを実行する。
\item   条件 motor\_type を判定し、真なら内部処理を行う。
\item     mt に motor\_type.lower() の結果を代入する。
\item     条件 mt in ('inwheel', 'hub') を判定し、真なら内部処理を行う。
\item       条件 not 'gear\_ratio' in model\_cfg を判定し、真なら内部処理を行う。
\item       条件 not 'gear\_eta' in model\_cfg を判定し、真なら内部処理を行う。
\item mpc\_cfg に dict(cfg.get('mpc', \{\}) or \{\}) の結果を代入する。
\item runtime\_mode に str(self.get\_parameter('profile\_runtime\_mode').value or '').strip() の結果を代入する。
\item runtime\_section に cfg.get(runtime\_mode, \{\}) if runtime\_mode else \{\} の結果を代入する。
\item runtime\_overrides に runtime\_section.get('mpc\_overrides', \{\}) if isinstance(runtime\_section, dict) else \{\} の結果を代入する。
\item 条件 isinstance(runtime\_overrides, dict) and runtime\_overrides を判定し、真なら内部処理を行う。
\item   mpc\_cfg.update(...) を実行する。
\item   self.get\_logger().info(...) を実行する。
\item 条件 isinstance(mpc\_cfg, dict) を判定し、真なら内部処理を行う。
\item   params に [] の結果を代入する。
\item   mpc\_cfg.items() を順に走査し、各要素を (key, val) に入れて処理する。
\item     条件 isinstance(val, dict) or val is None or (not self.has\_parameter(key)) を判定し、真なら内部処理を行う。
\item       Continue 文を実行する。
\item     例外処理を伴う try ブロックを実行する。
\item       params.append(...) を実行する。
\item       (TypeError, ValueError)を捕捉した場合:
\item       self.get\_logger().warn(...) を実行する。
\item   条件 params を判定し、真なら内部処理を行う。
\item     self.set\_parameters(...) を実行する。
            \end{enumerate}
            \paragraph{代表コード断片}
            \begin{lstlisting}[language=Python]
    def _apply_params_yaml(self, path: str):
        try:
            with open(path, 'r', encoding='utf-8') as f:
                cfg = yaml.safe_load(f) or {}
        except Exception as exc:
            self.get_logger().warn(f'params_yaml load failed: {exc}')
            return
        self.params_cfg = cfg if isinstance(cfg, dict) else {}
        model_cfg = cfg.get('model', cfg) if isinstance(cfg, dict) else {}
        if isinstance(model_cfg, dict):
            motor_type = None
            for k, v in model_cfg.items():
                if hasattr(self.model.p, k):
                    try:
                        setattr(self.model.p, k, float(v))
                    except Exception:
                        try:
                            setattr(self.model.p, k, v)
                        except Exception:
                            pass
                if k == 'motor_type':
                    motor_type = str(v)
                if k == 'drive_mode':
                    self.model.drive_mode = str(v)
                if k == 'drive_mode_tau_margin':
                    try:
                        self.model.drive_mode_tau_margin = float(v)
                    except Exception:
                        pass
                if k == 'ocv_soc_map':
                    try:
                        v_str = str(v).strip()
                        if v_str:
                            self.model.load_ocv_map(resolve_path(v_str, 'maps'))
                    except Exception:
...
\end{lstlisting}

\subsection{L172 関数 MPCNode.\_maybe\_reload\_forecast}
            \begin{itemize}[leftmargin=1.6em]
              \item 行範囲: L172--L190
              \item 所属: \path|MPCNode|
              \item 定義: \path|_maybe_reload_forecast(self)|
              \item docstring: 明示されていない。
              \item このブロックが直接呼ぶ主な関数・メソッド: \_load\_forecast\_file, get\_logger, getmtime, info, monotonic, warn
              \item 戻り値の要点: 戻り値が無いか、return 文が明示されていない。
              \item 主なローカル代入名: mtime, now
              \item 読み取る主なインスタンス属性: self.\_load\_forecast\_file, self.forecast\_mtime, self.forecast\_path, self.forecast\_reload\_sec, self.get\_logger, self.last\_forecast\_check
              \item 更新する主なインスタンス属性: self.forecast\_mtime, self.forecast\_reloaded, self.last\_forecast\_check
              \item 明示的に送出する例外: 明示的raiseはない。
              \item 制御構造の規模: 条件分岐 3、ループ 0、try 2
            \end{itemize}
            \paragraph{この定義を読むためのPython構文}
            \begin{itemize}[leftmargin=1.6em]
            \item 第1引数selfは、このメソッドを呼んだインスタンス自身を受け取る慣習名である。
\item try文は例外が起き得る処理と、異常時または終了時の経路を分ける。
            \end{itemize}
            \paragraph{上から順の処理}
            \begin{enumerate}[leftmargin=1.8em]
            \item 条件 self.forecast\_reload\_sec <= 0 を判定し、真なら内部処理を行う。
\item    を返す。
\item now に time.monotonic() の結果を代入する。
\item 条件 now - self.last\_forecast\_check < self.forecast\_reload\_sec を判定し、真なら内部処理を行う。
\item    を返す。
\item self.last\_forecast\_check に now の結果を代入する。
\item 例外処理を伴う try ブロックを実行する。
\item   mtime に os.path.getmtime(self.forecast\_path) の結果を代入する。
\item   Exceptionを捕捉した場合:
\item    を返す。
\item 条件 self.forecast\_mtime is None or mtime > self.forecast\_mtime を判定し、真なら内部処理を行う。
\item   self.forecast\_mtime に mtime の結果を代入する。
\item   例外処理を伴う try ブロックを実行する。
\item     self.\_load\_forecast\_file(...) を実行する。
\item     self.forecast\_reloaded に True の結果を代入する。
\item     self.get\_logger().info(...) を実行する。
\item     Exceptionを捕捉した場合:
\item     self.get\_logger().warn(...) を実行する。
            \end{enumerate}
            \paragraph{代表コード断片}
            \begin{lstlisting}[language=Python]
    def _maybe_reload_forecast(self):
        if self.forecast_reload_sec <= 0:
            return
        now = time.monotonic()
        if (now - self.last_forecast_check) < self.forecast_reload_sec:
            return
        self.last_forecast_check = now
        try:
            mtime = os.path.getmtime(self.forecast_path)
        except Exception:
            return
        if self.forecast_mtime is None or mtime > self.forecast_mtime:
            self.forecast_mtime = mtime
            try:
                self._load_forecast_file(self.forecast_path)
                self.forecast_reloaded = True
                self.get_logger().info('Forecast CSV reloaded.')
            except Exception as exc:
                self.get_logger().warn(f'Forecast reload failed: {exc}')
\end{lstlisting}

\subsection{L192 関数 MPCNode.\_on\_s\_km\_solar}
            \begin{itemize}[leftmargin=1.6em]
              \item 行範囲: L192--L206
              \item 所属: \path|MPCNode|
              \item 定義: \path|_on_s_km_solar(self, msg: Float32)|
              \item docstring: 明示されていない。
              \item このブロックが直接呼ぶ主な関数・メソッド: bool, finite\_float, float, get\_clock, get\_parameter, isfinite, now, update
              \item 戻り値の要点: 戻り値が無いか、return 文が明示されていない。
              \item 主なローカル代入名: current, now\_sec, raw
              \item 読み取る主なインスタンス属性: self.distance\_meas\_filter, self.distance\_meas\_max\_backtrack\_km, self.get\_clock, self.get\_parameter, self.s\_meas
              \item 更新する主なインスタンス属性: self.s\_km, self.s\_meas, self.s\_meas\_time
              \item 明示的に送出する例外: 明示的raiseはない。
              \item 制御構造の規模: 条件分岐 3、ループ 0、try 1
            \end{itemize}
            \paragraph{この定義を読むためのPython構文}
            \begin{itemize}[leftmargin=1.6em]
            \item 第1引数selfは、このメソッドを呼んだインスタンス自身を受け取る慣習名である。
\item try文は例外が起き得る処理と、異常時または終了時の経路を分ける。
            \end{itemize}
            \paragraph{上から順の処理}
            \begin{enumerate}[leftmargin=1.8em]
            \item 例外処理を伴う try ブロックを実行する。
\item   now\_sec に self.get\_clock().now().nanoseconds / 1000000000.0 の結果を代入する。
\item   raw に finite\_float(msg.data) の結果を代入する。
\item   条件 not math.isfinite(raw) を判定し、真なら内部処理を行う。
\item      を返す。
\item   current に float(self.distance\_meas\_filter.value) の結果を代入する。
\item   条件 math.isfinite(current) and raw < current - self.distance\_meas\_max\_backtrack\_km を判定し、真なら内部処理を行う。
\item      を返す。
\item   self.s\_meas に float(self.distance\_meas\_filter.update(raw, now=now\_sec)) の結果を代入する。
\item   self.s\_meas\_time に now\_sec の結果を代入する。
\item   条件 bool(self.get\_parameter('use\_measured\_s').value) を判定し、真なら内部処理を行う。
\item     self.s\_km に float(self.s\_meas) の結果を代入する。
\item   Exceptionを捕捉した場合:
\item   Pass 文を実行する。
            \end{enumerate}
            \paragraph{代表コード断片}
            \begin{lstlisting}[language=Python]
    def _on_s_km_solar(self, msg: Float32):
        try:
            now_sec = self.get_clock().now().nanoseconds / 1e9
            raw = finite_float(msg.data)
            if not math.isfinite(raw):
                return
            current = float(self.distance_meas_filter.value)
            if math.isfinite(current) and raw < current - self.distance_meas_max_backtrack_km:
                return
            self.s_meas = float(self.distance_meas_filter.update(raw, now=now_sec))
            self.s_meas_time = now_sec
            if bool(self.get_parameter('use_measured_s').value):
                self.s_km = float(self.s_meas)
        except Exception:
            pass
\end{lstlisting}

\subsection{L208 関数 MPCNode.\_on\_speed\_solar}
            \begin{itemize}[leftmargin=1.6em]
              \item 行範囲: L208--L214
              \item 所属: \path|MPCNode|
              \item 定義: \path|_on_speed_solar(self, msg: Float32)|
              \item docstring: 明示されていない。
              \item このブロックが直接呼ぶ主な関数・メソッド: float, get\_clock, now, update
              \item 戻り値の要点: 戻り値が無いか、return 文が明示されていない。
              \item 主なローカル代入名: now\_sec
              \item 読み取る主なインスタンス属性: self.get\_clock, self.speed\_meas\_filter
              \item 更新する主なインスタンス属性: self.v\_meas\_time, self.v\_now
              \item 明示的に送出する例外: 明示的raiseはない。
              \item 制御構造の規模: 条件分岐 0、ループ 0、try 1
            \end{itemize}
            \paragraph{この定義を読むためのPython構文}
            \begin{itemize}[leftmargin=1.6em]
            \item 第1引数selfは、このメソッドを呼んだインスタンス自身を受け取る慣習名である。
\item try文は例外が起き得る処理と、異常時または終了時の経路を分ける。
            \end{itemize}
            \paragraph{上から順の処理}
            \begin{enumerate}[leftmargin=1.8em]
            \item 例外処理を伴う try ブロックを実行する。
\item   now\_sec に self.get\_clock().now().nanoseconds / 1000000000.0 の結果を代入する。
\item   self.v\_now に float(self.speed\_meas\_filter.update(msg.data, now=now\_sec)) の結果を代入する。
\item   self.v\_meas\_time に now\_sec の結果を代入する。
\item   Exceptionを捕捉した場合:
\item   Pass 文を実行する。
            \end{enumerate}
            \paragraph{代表コード断片}
            \begin{lstlisting}[language=Python]
    def _on_speed_solar(self, msg: Float32):
        try:
            now_sec = self.get_clock().now().nanoseconds / 1e9
            self.v_now = float(self.speed_meas_filter.update(msg.data, now=now_sec))
            self.v_meas_time = now_sec
        except Exception:
            pass
\end{lstlisting}

\subsection{L216 関数 MPCNode.\_on\_soc\_solar}
            \begin{itemize}[leftmargin=1.6em]
              \item 行範囲: L216--L223
              \item 所属: \path|MPCNode|
              \item 定義: \path|_on_soc_solar(self, msg: Float32)|
              \item docstring: 明示されていない。
              \item このブロックが直接呼ぶ主な関数・メソッド: clip, float, get\_clock, now, update
              \item 戻り値の要点: 戻り値が無いか、return 文が明示されていない。
              \item 主なローカル代入名: now\_sec
              \item 読み取る主なインスタンス属性: self.get\_clock, self.model, self.soc\_meas\_filter, self.solar\_soc\_meas
              \item 更新する主なインスタンス属性: self.soc\_meas\_time, self.solar\_soc\_meas, self.z
              \item 明示的に送出する例外: 明示的raiseはない。
              \item 制御構造の規模: 条件分岐 0、ループ 0、try 1
            \end{itemize}
            \paragraph{この定義を読むためのPython構文}
            \begin{itemize}[leftmargin=1.6em]
            \item 第1引数selfは、このメソッドを呼んだインスタンス自身を受け取る慣習名である。
\item try文は例外が起き得る処理と、異常時または終了時の経路を分ける。
            \end{itemize}
            \paragraph{上から順の処理}
            \begin{enumerate}[leftmargin=1.8em]
            \item 例外処理を伴う try ブロックを実行する。
\item   now\_sec に self.get\_clock().now().nanoseconds / 1000000000.0 の結果を代入する。
\item   self.solar\_soc\_meas に float(self.soc\_meas\_filter.update(msg.data, now=now\_sec)) の結果を代入する。
\item   self.soc\_meas\_time に now\_sec の結果を代入する。
\item   self.z に float(np.clip(self.solar\_soc\_meas, self.model.p.soc\_min, self.model.p.soc\_max)) の結果を代入する。
\item   Exceptionを捕捉した場合:
\item   Pass 文を実行する。
            \end{enumerate}
            \paragraph{代表コード断片}
            \begin{lstlisting}[language=Python]
    def _on_soc_solar(self, msg: Float32):
        try:
            now_sec = self.get_clock().now().nanoseconds / 1e9
            self.solar_soc_meas = float(self.soc_meas_filter.update(msg.data, now=now_sec))
            self.soc_meas_time = now_sec
            self.z = float(np.clip(self.solar_soc_meas, self.model.p.soc_min, self.model.p.soc_max))
        except Exception:
            pass
\end{lstlisting}

\subsection{L225 関数 MPCNode.\_on\_tb\_solar}
            \begin{itemize}[leftmargin=1.6em]
              \item 行範囲: L225--L232
              \item 所属: \path|MPCNode|
              \item 定義: \path|_on_tb_solar(self, msg: Float32)|
              \item docstring: 明示されていない。
              \item このブロックが直接呼ぶ主な関数・メソッド: clip, float, get\_clock, now, update
              \item 戻り値の要点: 戻り値が無いか、return 文が明示されていない。
              \item 主なローカル代入名: now\_sec
              \item 読み取る主なインスタンス属性: self.get\_clock, self.model, self.solar\_tb\_meas, self.tb\_meas\_filter
              \item 更新する主なインスタンス属性: self.Tb, self.solar\_tb\_meas, self.tb\_meas\_time
              \item 明示的に送出する例外: 明示的raiseはない。
              \item 制御構造の規模: 条件分岐 0、ループ 0、try 1
            \end{itemize}
            \paragraph{この定義を読むためのPython構文}
            \begin{itemize}[leftmargin=1.6em]
            \item 第1引数selfは、このメソッドを呼んだインスタンス自身を受け取る慣習名である。
\item try文は例外が起き得る処理と、異常時または終了時の経路を分ける。
            \end{itemize}
            \paragraph{上から順の処理}
            \begin{enumerate}[leftmargin=1.8em]
            \item 例外処理を伴う try ブロックを実行する。
\item   now\_sec に self.get\_clock().now().nanoseconds / 1000000000.0 の結果を代入する。
\item   self.solar\_tb\_meas に float(self.tb\_meas\_filter.update(msg.data, now=now\_sec)) の結果を代入する。
\item   self.tb\_meas\_time に now\_sec の結果を代入する。
\item   self.Tb に float(np.clip(self.solar\_tb\_meas, self.model.p.T\_min, self.model.p.T\_max)) の結果を代入する。
\item   Exceptionを捕捉した場合:
\item   Pass 文を実行する。
            \end{enumerate}
            \paragraph{代表コード断片}
            \begin{lstlisting}[language=Python]
    def _on_tb_solar(self, msg: Float32):
        try:
            now_sec = self.get_clock().now().nanoseconds / 1e9
            self.solar_tb_meas = float(self.tb_meas_filter.update(msg.data, now=now_sec))
            self.tb_meas_time = now_sec
            self.Tb = float(np.clip(self.solar_tb_meas, self.model.p.T_min, self.model.p.T_max))
        except Exception:
            pass
\end{lstlisting}

\subsection{L234 関数 MPCNode.\_on\_i\_solar}
            \begin{itemize}[leftmargin=1.6em]
              \item 行範囲: L234--L240
              \item 所属: \path|MPCNode|
              \item 定義: \path|_on_i_solar(self, msg: Float32)|
              \item docstring: 明示されていない。
              \item このブロックが直接呼ぶ主な関数・メソッド: float, get\_clock, now, update
              \item 戻り値の要点: 戻り値が無いか、return 文が明示されていない。
              \item 主なローカル代入名: now\_sec
              \item 読み取る主なインスタンス属性: self.current\_meas\_filter, self.get\_clock
              \item 更新する主なインスタンス属性: self.current\_meas\_time, self.solar\_i\_meas
              \item 明示的に送出する例外: 明示的raiseはない。
              \item 制御構造の規模: 条件分岐 0、ループ 0、try 1
            \end{itemize}
            \paragraph{この定義を読むためのPython構文}
            \begin{itemize}[leftmargin=1.6em]
            \item 第1引数selfは、このメソッドを呼んだインスタンス自身を受け取る慣習名である。
\item try文は例外が起き得る処理と、異常時または終了時の経路を分ける。
            \end{itemize}
            \paragraph{上から順の処理}
            \begin{enumerate}[leftmargin=1.8em]
            \item 例外処理を伴う try ブロックを実行する。
\item   now\_sec に self.get\_clock().now().nanoseconds / 1000000000.0 の結果を代入する。
\item   self.solar\_i\_meas に float(self.current\_meas\_filter.update(msg.data, now=now\_sec)) の結果を代入する。
\item   self.current\_meas\_time に now\_sec の結果を代入する。
\item   Exceptionを捕捉した場合:
\item   Pass 文を実行する。
            \end{enumerate}
            \paragraph{代表コード断片}
            \begin{lstlisting}[language=Python]
    def _on_i_solar(self, msg: Float32):
        try:
            now_sec = self.get_clock().now().nanoseconds / 1e9
            self.solar_i_meas = float(self.current_meas_filter.update(msg.data, now=now_sec))
            self.current_meas_time = now_sec
        except Exception:
            pass
\end{lstlisting}

\subsection{L242 関数 MPCNode.\_on\_v\_solar}
            \begin{itemize}[leftmargin=1.6em]
              \item 行範囲: L242--L248
              \item 所属: \path|MPCNode|
              \item 定義: \path|_on_v_solar(self, msg: Float32)|
              \item docstring: 明示されていない。
              \item このブロックが直接呼ぶ主な関数・メソッド: float, get\_clock, now, update
              \item 戻り値の要点: 戻り値が無いか、return 文が明示されていない。
              \item 主なローカル代入名: now\_sec
              \item 読み取る主なインスタンス属性: self.get\_clock, self.voltage\_meas\_filter
              \item 更新する主なインスタンス属性: self.solar\_v\_meas, self.voltage\_meas\_time
              \item 明示的に送出する例外: 明示的raiseはない。
              \item 制御構造の規模: 条件分岐 0、ループ 0、try 1
            \end{itemize}
            \paragraph{この定義を読むためのPython構文}
            \begin{itemize}[leftmargin=1.6em]
            \item 第1引数selfは、このメソッドを呼んだインスタンス自身を受け取る慣習名である。
\item try文は例外が起き得る処理と、異常時または終了時の経路を分ける。
            \end{itemize}
            \paragraph{上から順の処理}
            \begin{enumerate}[leftmargin=1.8em]
            \item 例外処理を伴う try ブロックを実行する。
\item   now\_sec に self.get\_clock().now().nanoseconds / 1000000000.0 の結果を代入する。
\item   self.solar\_v\_meas に float(self.voltage\_meas\_filter.update(msg.data, now=now\_sec)) の結果を代入する。
\item   self.voltage\_meas\_time に now\_sec の結果を代入する。
\item   Exceptionを捕捉した場合:
\item   Pass 文を実行する。
            \end{enumerate}
            \paragraph{代表コード断片}
            \begin{lstlisting}[language=Python]
    def _on_v_solar(self, msg: Float32):
        try:
            now_sec = self.get_clock().now().nanoseconds / 1e9
            self.solar_v_meas = float(self.voltage_meas_filter.update(msg.data, now=now_sec))
            self.voltage_meas_time = now_sec
        except Exception:
            pass
\end{lstlisting}

\subsection{L250 関数 MPCNode.\_measured\_distance\_km}
            \begin{itemize}[leftmargin=1.6em]
              \item 行範囲: L250--L254
              \item 所属: \path|MPCNode|
              \item 定義: \path|_measured_distance_km(self) -> float|
              \item docstring: 明示されていない。
              \item このブロックが直接呼ぶ主な関数・メソッド: finite\_float, fresh\_enough, get\_clock, now
              \item 戻り値の要点: math.nan / finite\_float(self.s\_meas)
              \item 主なローカル代入名: now\_sec
              \item 読み取る主なインスタンス属性: self.distance\_meas\_timeout\_sec, self.get\_clock, self.s\_meas, self.s\_meas\_time
              \item 更新する主なインスタンス属性: 特になし。
              \item 明示的に送出する例外: 明示的raiseはない。
              \item 制御構造の規模: 条件分岐 1、ループ 0、try 0
            \end{itemize}
            \paragraph{この定義を読むためのPython構文}
            \begin{itemize}[leftmargin=1.6em]
            \item 第1引数selfは、このメソッドを呼んだインスタンス自身を受け取る慣習名である。
\item `->`以後は戻り値の型注釈であり、実行時の値を自動変換する命令ではない。
            \end{itemize}
            \paragraph{上から順の処理}
            \begin{enumerate}[leftmargin=1.8em]
            \item now\_sec に self.get\_clock().now().nanoseconds / 1000000000.0 の結果を代入する。
\item 条件 fresh\_enough(self.s\_meas\_time, now\_sec, self.distance\_meas\_timeout\_sec) を判定し、真なら内部処理を行う。
\item   finite\_float(self.s\_meas) を返す。
\item math.nan を返す。
            \end{enumerate}
            \paragraph{代表コード断片}
            \begin{lstlisting}[language=Python]
    def _measured_distance_km(self) -> float:
        now_sec = self.get_clock().now().nanoseconds / 1e9
        if fresh_enough(self.s_meas_time, now_sec, self.distance_meas_timeout_sec):
            return finite_float(self.s_meas)
        return math.nan
\end{lstlisting}

\subsection{L256 関数 MPCNode.\_sync\_measured\_state}
            \begin{itemize}[leftmargin=1.6em]
              \item 行範囲: L256--L268
              \item 所属: \path|MPCNode|
              \item 定義: \path|_sync_measured_state(self) -> None|
              \item docstring: 明示されていない。
              \item このブロックが直接呼ぶ主な関数・メソッド: \_measured\_distance\_km, bool, clip, float, fresh\_enough, get\_clock, get\_parameter, isfinite, now
              \item 戻り値の要点: 戻り値が無いか、return 文が明示されていない。
              \item 主なローカル代入名: distance\_km, now\_sec
              \item 読み取る主なインスタンス属性: self.\_measured\_distance\_km, self.battery\_meas\_timeout\_sec, self.get\_clock, self.get\_parameter, self.model, self.soc\_meas\_time, self.solar\_soc\_meas, self.solar\_tb\_meas, self.speed\_meas\_timeout\_sec, self.tb\_meas\_time, self.v\_meas\_time
              \item 更新する主なインスタンス属性: self.Tb, self.s\_km, self.v\_now, self.z
              \item 明示的に送出する例外: 明示的raiseはない。
              \item 制御構造の規模: 条件分岐 6、ループ 0、try 0
            \end{itemize}
            \paragraph{この定義を読むためのPython構文}
            \begin{itemize}[leftmargin=1.6em]
            \item 第1引数selfは、このメソッドを呼んだインスタンス自身を受け取る慣習名である。
\item `->`以後は戻り値の型注釈であり、実行時の値を自動変換する命令ではない。
            \end{itemize}
            \paragraph{上から順の処理}
            \begin{enumerate}[leftmargin=1.8em]
            \item now\_sec に self.get\_clock().now().nanoseconds / 1000000000.0 の結果を代入する。
\item 条件 not fresh\_enough(self.v\_meas\_time, now\_sec, self.speed\_meas\_timeout\_sec) を判定し、真なら内部処理を行う。
\item   self.v\_now に math.nan の結果を代入する。
\item distance\_km に self.\_measured\_distance\_km() の結果を代入する。
\item 条件 bool(self.get\_parameter('use\_measured\_s').value) and math.isfinite(distance\_km) を判定し、真なら内部処理を行う。
\item   self.s\_km に float(distance\_km) の結果を代入する。
\item 条件 fresh\_enough(self.soc\_meas\_time, now\_sec, self.battery\_meas\_timeout\_sec) を判定し、真なら内部処理を行う。
\item   条件 math.isfinite(self.solar\_soc\_meas) を判定し、真なら内部処理を行う。
\item     self.z に float(np.clip(self.solar\_soc\_meas, self.model.p.soc\_min, self.model.p.soc\_max)) の結果を代入する。
\item 条件 fresh\_enough(self.tb\_meas\_time, now\_sec, self.battery\_meas\_timeout\_sec) を判定し、真なら内部処理を行う。
\item   条件 math.isfinite(self.solar\_tb\_meas) を判定し、真なら内部処理を行う。
\item     self.Tb に float(np.clip(self.solar\_tb\_meas, self.model.p.T\_min, self.model.p.T\_max)) の結果を代入する。
            \end{enumerate}
            \paragraph{代表コード断片}
            \begin{lstlisting}[language=Python]
    def _sync_measured_state(self) -> None:
        now_sec = self.get_clock().now().nanoseconds / 1e9
        if not fresh_enough(self.v_meas_time, now_sec, self.speed_meas_timeout_sec):
            self.v_now = math.nan
        distance_km = self._measured_distance_km()
        if bool(self.get_parameter('use_measured_s').value) and math.isfinite(distance_km):
            self.s_km = float(distance_km)
        if fresh_enough(self.soc_meas_time, now_sec, self.battery_meas_timeout_sec):
            if math.isfinite(self.solar_soc_meas):
                self.z = float(np.clip(self.solar_soc_meas, self.model.p.soc_min, self.model.p.soc_max))
        if fresh_enough(self.tb_meas_time, now_sec, self.battery_meas_timeout_sec):
            if math.isfinite(self.solar_tb_meas):
                self.Tb = float(np.clip(self.solar_tb_meas, self.model.p.T_min, self.model.p.T_max))
\end{lstlisting}

\subsection{L270 関数 MPCNode.\_on\_calib\_solar\_gain}
            \begin{itemize}[leftmargin=1.6em]
              \item 行範囲: L270--L274
              \item 所属: \path|MPCNode|
              \item 定義: \path|_on_calib_solar_gain(self, msg: Float32)|
              \item docstring: 明示されていない。
              \item このブロックが直接呼ぶ主な関数・メソッド: clip, float
              \item 戻り値の要点: 戻り値が無いか、return 文が明示されていない。
              \item 主なローカル代入名: 特になし。
              \item 読み取る主なインスタンス属性: 特になし。
              \item 更新する主なインスタンス属性: self.solar\_gain
              \item 明示的に送出する例外: 明示的raiseはない。
              \item 制御構造の規模: 条件分岐 0、ループ 0、try 1
            \end{itemize}
            \paragraph{この定義を読むためのPython構文}
            \begin{itemize}[leftmargin=1.6em]
            \item 第1引数selfは、このメソッドを呼んだインスタンス自身を受け取る慣習名である。
\item try文は例外が起き得る処理と、異常時または終了時の経路を分ける。
            \end{itemize}
            \paragraph{上から順の処理}
            \begin{enumerate}[leftmargin=1.8em]
            \item 例外処理を伴う try ブロックを実行する。
\item   self.solar\_gain に float(np.clip(float(msg.data), 0.2, 2.5)) の結果を代入する。
\item   Exceptionを捕捉した場合:
\item    を返す。
            \end{enumerate}
            \paragraph{代表コード断片}
            \begin{lstlisting}[language=Python]
    def _on_calib_solar_gain(self, msg: Float32):
        try:
            self.solar_gain = float(np.clip(float(msg.data), 0.2, 2.5))
        except Exception:
            return
\end{lstlisting}

\subsection{L276 関数 MPCNode.\_on\_calib\_drive\_power\_gain}
            \begin{itemize}[leftmargin=1.6em]
              \item 行範囲: L276--L282
              \item 所属: \path|MPCNode|
              \item 定義: \path|_on_calib_drive_power_gain(self, msg: Float32)|
              \item docstring: 明示されていない。
              \item このブロックが直接呼ぶ主な関数・メソッド: clip, float, max
              \item 戻り値の要点: 戻り値が無いか、return 文が明示されていない。
              \item 主なローカル代入名: gain
              \item 読み取る主なインスタンス属性: self.base\_drive\_eff\_scale, self.base\_regen\_eff\_scale, self.model
              \item 更新する主なインスタンス属性: 特になし。
              \item 明示的に送出する例外: 明示的raiseはない。
              \item 制御構造の規模: 条件分岐 0、ループ 0、try 1
            \end{itemize}
            \paragraph{この定義を読むためのPython構文}
            \begin{itemize}[leftmargin=1.6em]
            \item 第1引数selfは、このメソッドを呼んだインスタンス自身を受け取る慣習名である。
\item try文は例外が起き得る処理と、異常時または終了時の経路を分ける。
            \end{itemize}
            \paragraph{上から順の処理}
            \begin{enumerate}[leftmargin=1.8em]
            \item 例外処理を伴う try ブロックを実行する。
\item   gain に float(np.clip(float(msg.data), 0.4, 2.5)) の結果を代入する。
\item   self.model.p.drive\_eff\_scale に float(self.base\_drive\_eff\_scale) / max(gain, 0.001) の結果を代入する。
\item   self.model.p.regen\_eff\_scale に float(self.base\_regen\_eff\_scale) / max(gain, 0.001) の結果を代入する。
\item   Exceptionを捕捉した場合:
\item    を返す。
            \end{enumerate}
            \paragraph{代表コード断片}
            \begin{lstlisting}[language=Python]
    def _on_calib_drive_power_gain(self, msg: Float32):
        try:
            gain = float(np.clip(float(msg.data), 0.4, 2.5))
            self.model.p.drive_eff_scale = float(self.base_drive_eff_scale) / max(gain, 1.0e-3)
            self.model.p.regen_eff_scale = float(self.base_regen_eff_scale) / max(gain, 1.0e-3)
        except Exception:
            return
\end{lstlisting}

\subsection{L284 関数 MPCNode.\_on\_calib\_aux\_power}
            \begin{itemize}[leftmargin=1.6em]
              \item 行範囲: L284--L288
              \item 所属: \path|MPCNode|
              \item 定義: \path|_on_calib_aux_power(self, msg: Float32)|
              \item docstring: 明示されていない。
              \item このブロックが直接呼ぶ主な関数・メソッド: float, max
              \item 戻り値の要点: 戻り値が無いか、return 文が明示されていない。
              \item 主なローカル代入名: 特になし。
              \item 読み取る主なインスタンス属性: self.model
              \item 更新する主なインスタンス属性: 特になし。
              \item 明示的に送出する例外: 明示的raiseはない。
              \item 制御構造の規模: 条件分岐 0、ループ 0、try 1
            \end{itemize}
            \paragraph{この定義を読むためのPython構文}
            \begin{itemize}[leftmargin=1.6em]
            \item 第1引数selfは、このメソッドを呼んだインスタンス自身を受け取る慣習名である。
\item try文は例外が起き得る処理と、異常時または終了時の経路を分ける。
            \end{itemize}
            \paragraph{上から順の処理}
            \begin{enumerate}[leftmargin=1.8em]
            \item 例外処理を伴う try ブロックを実行する。
\item   self.model.aux\_power\_override\_w に max(0.0, float(msg.data)) の結果を代入する。
\item   Exceptionを捕捉した場合:
\item    を返す。
            \end{enumerate}
            \paragraph{代表コード断片}
            \begin{lstlisting}[language=Python]
    def _on_calib_aux_power(self, msg: Float32):
        try:
            self.model.aux_power_override_w = max(0.0, float(msg.data))
        except Exception:
            return
\end{lstlisting}

\subsection{L291 関数 MPCNode.\_init\_solar}
            \begin{itemize}[leftmargin=1.6em]
              \item 行範囲: L291--L778
              \item 所属: \path|MPCNode|
              \item 定義: \path|_init_solar(self)|
              \item docstring: 明示されていない。
              \item このブロックが直接呼ぶ主な関数・メソッド: BatteryMHE, Params, RobustScalarFilter, SolarCarModel, \_apply\_params\_yaml, \_load\_forecast\_file, \_load\_stops, abs, astimezone, bool, clip, create\_publisher
              \item 戻り値の要点: 戻り値が無いか、return 文が明示されていない。
              \item 主なローカル代入名: battery\_tau, drive\_map, drive\_map\_eco, drive\_map\_power, drive\_schedule\_yaml, expected\_dt, fcsv, initial\_policy\_csv, maps\_dir, mppt\_map, ocv\_map, panel\_map, params\_yaml, regen\_map, regen\_map\_eco, regen\_map\_power, rint\_map, route\_profile\_csv, speed\_profile\_csv, start\_time\_utc
              \item 読み取る主なインスタンス属性: self.\_apply\_params\_yaml, self.\_load\_forecast\_file, self.\_load\_stops, self.\_on\_calib\_aux\_power, self.\_on\_calib\_drive\_power\_gain, self.\_on\_calib\_solar\_gain, self.\_on\_i\_solar, self.\_on\_s\_km\_solar, self.\_on\_soc\_solar, self.\_on\_speed\_solar, self.\_on\_tb\_solar, self.\_on\_v\_solar, self.\_publish\_lower\_command\_cycle, self.\_step\_lower, self.\_step\_solar, self.\_warned\_lower\_rate, self.base\_drive\_eff\_scale, self.command\_callback\_group, self.create\_publisher, self.create\_subscription, self.create\_timer, self.declare\_parameter, self.drive\_schedule, self.dt
              \item 更新する主なインスタンス属性: self.Np, self.Tb, self.\_forecast\_warned\_out\_of\_range, self.\_lower\_solve\_logged, self.\_lower\_upper\_busy\_logged, self.\_upper\_solve\_in\_progress, self.\_upper\_solve\_logged, self.\_warned\_lower\_rate, self.active\_control\_stop\_index, self.base\_drive\_eff\_scale, self.base\_regen\_eff\_scale, self.battery\_meas\_timeout\_sec, self.completed\_control\_stops, self.control\_stop\_end\_monotonic, self.control\_stop\_hold, self.control\_stop\_position\_initialized, self.control\_stop\_tilt\_fraction, self.current\_forecast\_time\_utc, self.current\_meas\_filter, self.current\_meas\_time, self.distance\_meas\_filter, self.distance\_meas\_max\_backtrack\_km, self.distance\_meas\_timeout\_sec, self.drive\_schedule
              \item 明示的に送出する例外: 明示的raiseはない。
              \item 制御構造の規模: 条件分岐 15、ループ 0、try 5
            \end{itemize}
            \paragraph{この定義を読むためのPython構文}
            \begin{itemize}[leftmargin=1.6em]
            \item 第1引数selfは、このメソッドを呼んだインスタンス自身を受け取る慣習名である。
\item try文は例外が起き得る処理と、異常時または終了時の経路を分ける。
            \end{itemize}
            \paragraph{上から順の処理}
            \begin{enumerate}[leftmargin=1.8em]
            \item self.declare\_parameter(...) を実行する。
\item self.declare\_parameter(...) を実行する。
\item self.declare\_parameter(...) を実行する。
\item self.declare\_parameter(...) を実行する。
\item self.declare\_parameter(...) を実行する。
\item self.declare\_parameter(...) を実行する。
\item self.declare\_parameter(...) を実行する。
\item self.declare\_parameter(...) を実行する。
\item self.declare\_parameter(...) を実行する。
\item self.declare\_parameter(...) を実行する。
\item self.declare\_parameter(...) を実行する。
\item self.declare\_parameter(...) を実行する。
\item self.declare\_parameter(...) を実行する。
\item self.declare\_parameter(...) を実行する。
\item self.declare\_parameter(...) を実行する。
\item self.declare\_parameter(...) を実行する。
\item self.declare\_parameter(...) を実行する。
\item self.declare\_parameter(...) を実行する。
\item self.declare\_parameter(...) を実行する。
\item self.declare\_parameter(...) を実行する。
\item self.declare\_parameter(...) を実行する。
\item self.declare\_parameter(...) を実行する。
\item self.declare\_parameter(...) を実行する。
\item self.declare\_parameter(...) を実行する。
\item self.declare\_parameter(...) を実行する。
\item self.declare\_parameter(...) を実行する。
\item self.declare\_parameter(...) を実行する。
\item self.declare\_parameter(...) を実行する。
\item self.declare\_parameter(...) を実行する。
\item self.declare\_parameter(...) を実行する。
\item self.declare\_parameter(...) を実行する。
\item self.declare\_parameter(...) を実行する。
\item self.declare\_parameter(...) を実行する。
\item self.declare\_parameter(...) を実行する。
\item self.declare\_parameter(...) を実行する。
\item self.declare\_parameter(...) を実行する。
\item self.declare\_parameter(...) を実行する。
\item self.declare\_parameter(...) を実行する。
\item self.declare\_parameter(...) を実行する。
\item self.declare\_parameter(...) を実行する。
\item self.declare\_parameter(...) を実行する。
\item self.declare\_parameter(...) を実行する。
\item self.declare\_parameter(...) を実行する。
\item self.declare\_parameter(...) を実行する。
\item self.declare\_parameter(...) を実行する。
\item self.declare\_parameter(...) を実行する。
\item self.declare\_parameter(...) を実行する。
\item self.declare\_parameter(...) を実行する。
\item self.declare\_parameter(...) を実行する。
\item self.declare\_parameter(...) を実行する。
\item self.declare\_parameter(...) を実行する。
\item self.declare\_parameter(...) を実行する。
\item self.declare\_parameter(...) を実行する。
\item self.declare\_parameter(...) を実行する。
\item self.declare\_parameter(...) を実行する。
\item self.declare\_parameter(...) を実行する。
\item self.declare\_parameter(...) を実行する。
\item self.declare\_parameter(...) を実行する。
\item self.declare\_parameter(...) を実行する。
\item self.declare\_parameter(...) を実行する。
\item self.declare\_parameter(...) を実行する。
\item self.declare\_parameter(...) を実行する。
\item self.declare\_parameter(...) を実行する。
\item self.declare\_parameter(...) を実行する。
\item self.declare\_parameter(...) を実行する。
\item self.declare\_parameter(...) を実行する。
\item self.declare\_parameter(...) を実行する。
\item self.declare\_parameter(...) を実行する。
\item self.declare\_parameter(...) を実行する。
\item self.declare\_parameter(...) を実行する。
\item self.declare\_parameter(...) を実行する。
\item self.declare\_parameter(...) を実行する。
\item self.declare\_parameter(...) を実行する。
\item self.declare\_parameter(...) を実行する。
\item self.declare\_parameter(...) を実行する。
\item self.declare\_parameter(...) を実行する。
\item self.declare\_parameter(...) を実行する。
\item self.declare\_parameter(...) を実行する。
\item self.declare\_parameter(...) を実行する。
\item self.declare\_parameter(...) を実行する。
            \end{enumerate}
            \paragraph{代表コード断片}
            \begin{lstlisting}[language=Python]
    def _init_solar(self):
        self.declare_parameter('forecast_csv', 'inputs/forecast_10min.csv')
        self.declare_parameter('maps_dir', 'maps')
        self.declare_parameter('drive_eff_map', '')
        self.declare_parameter('regen_eff_map', '')
        self.declare_parameter('rint_map', '')
        self.declare_parameter('drive_map_eco', '')
        self.declare_parameter('drive_map_power', '')
        self.declare_parameter('regen_map_eco', '')
        self.declare_parameter('regen_map_power', '')
        self.declare_parameter('panel_eff_map', '')
        self.declare_parameter('mppt_eff_map', '')
        self.declare_parameter('ocv_soc_map', '')
        self.declare_parameter('dt', 600.0)                 # 10 min [s]
        self.declare_parameter('horizon_steps', 9)
        self.declare_parameter('v_max_kmh', 110.0)
        self.declare_parameter('terminal_soc_min', 0.10)
        self.declare_parameter('stop_yaml', 'inputs/stop_points.yaml')
        self.declare_parameter('forecast_time_mode', 'auto')  # auto|absolute|relative|loop
        self.declare_parameter('forecast_time_tz', 'UTC')
        self.declare_parameter('forecast_start_time_utc', '')
        self.declare_parameter('forecast_time_offset_sec', 0.0)
        self.declare_parameter('forecast_reload_sec', 60.0)
        self.declare_parameter('replan_on_forecast_reload', False)
        self.declare_parameter('params_yaml', '')
        self.declare_parameter('profile_runtime_mode', '')
        self.declare_parameter('route_profile_csv', '')
        self.declare_parameter('speed_profile_csv', '')
        self.declare_parameter('drive_schedule_yaml', '')
        self.declare_parameter('initial_upper_policy_csv', '')
        self.declare_parameter('use_measured_s', True)
        self.declare_parameter('use_measured_speed', True)
        self.declare_parameter('soc0', 0.95)
        self.declare_parameter('Tb0', 30.0)
        self.declare_parameter('s0_km', 0.0)
...
\end{lstlisting}

\subsection{L780 関数 MPCNode.\_current\_bin\_index}
            \begin{itemize}[leftmargin=1.6em]
              \item 行範囲: L780--L828
              \item 所属: \path|MPCNode|
              \item 定義: \path|_current_bin_index(self) -> int|
              \item docstring: 明示されていない。
              \item このブロックが直接呼ぶ主な関数・メソッド: clip, datetime64, get\_logger, getattr, int, len, lower, max, now, searchsorted, str, timedelta
              \item 戻り値の要点: int(np.clip(self.k, 0, max(0, bin\_count - 1))) / int(idx \% bin\_count) if mode == 'loop' else int(np.clip(idx, 0, bin\_count - 1)) / int(np.clip(idx, 0, len(self.forecast\_times) - 1)) / 0
              \item 主なローカル代入名: bin\_count, elapsed, has\_time, idx, lookup\_time, loop\_sec, mode, now, t\_first, t\_last, t\_max, t\_min, t\_series
              \item 読み取る主なインスタンス属性: self.df, self.dt, self.forecast\_relative\_wall\_start, self.forecast\_start\_time, self.forecast\_time\_mode, self.forecast\_time\_offset, self.forecast\_times, self.get\_logger, self.k
              \item 更新する主なインスタンス属性: self.\_forecast\_warned\_out\_of\_range, self.current\_forecast\_time\_utc
              \item 明示的に送出する例外: 明示的raiseはない。
              \item 制御構造の規模: 条件分岐 8、ループ 0、try 0
            \end{itemize}
            \paragraph{この定義を読むためのPython構文}
            \begin{itemize}[leftmargin=1.6em]
            \item 第1引数selfは、このメソッドを呼んだインスタンス自身を受け取る慣習名である。
\item `->`以後は戻り値の型注釈であり、実行時の値を自動変換する命令ではない。
            \end{itemize}
            \paragraph{上から順の処理}
            \begin{enumerate}[leftmargin=1.8em]
            \item mode に str(self.forecast\_time\_mode).lower() の結果を代入する。
\item has\_time に len(getattr(self, 'forecast\_times', [])) > 0 の結果を代入する。
\item 条件 mode == 'auto' を判定し、真なら内部処理を行う。
\item   mode に 'absolute' if has\_time else 'relative' の結果を代入する。
\item now に datetime.now(timezone.utc) + timedelta(seconds=self.forecast\_time\_offset) の結果を代入する。
\item lookup\_time に now の結果を代入する。
\item 条件 mode == 'absolute' and has\_time を判定し、真なら内部処理を行う。
\item   t\_series に self.forecast\_times.values の結果を代入する。
\item   t\_min に self.forecast\_times[0] の結果を代入する。
\item   t\_max に self.forecast\_times[-1] の結果を代入する。
\item   条件 now < t\_min or now > t\_max を判定し、真なら内部処理を行う。
\item     条件 not getattr(self, '\_forecast\_warned\_out\_of\_range', False) を判定し、真なら内部処理を行う。
\item       self.get\_logger().warn(...) を実行する。
\item       self.\_forecast\_warned\_out\_of\_range に True の結果を代入する。
\item     mode に 'relative' の結果を代入する。
\item     上の条件が偽の場合:
\item     idx に int(np.searchsorted(t\_series, np.datetime64(now)) - 1) の結果を代入する。
\item     self.current\_forecast\_time\_utc に now の結果を代入する。
\item     int(np.clip(idx, 0, len(self.forecast\_times) - 1)) を返す。
\item 条件 mode in ('relative', 'loop') or not has\_time を判定し、真なら内部処理を行う。
\item   elapsed に (now - self.forecast\_relative\_wall\_start).total\_seconds() の結果を代入する。
\item   elapsed に max(0.0, elapsed) の結果を代入する。
\item   bin\_count に len(self.forecast\_times) if has\_time else len(self.df) の結果を代入する。
\item   条件 bin\_count == 0 を判定し、真なら内部処理を行う。
\item     0 を返す。
\item   条件 has\_time を判定し、真なら内部処理を行う。
\item     lookup\_time に self.forecast\_start\_time + timedelta(seconds=elapsed) の結果を代入する。
\item     条件 mode == 'loop' and len(self.forecast\_times) > 1 を判定し、真なら内部処理を行う。
\item       t\_first に self.forecast\_times[0].to\_pydatetime() の結果を代入する。
\item       t\_last に self.forecast\_times[-1].to\_pydatetime() の結果を代入する。
\item       loop\_sec に max(1.0, (t\_last - t\_first).total\_seconds()) の結果を代入する。
\item       lookup\_time に t\_first + timedelta(seconds=(lookup\_time - t\_first).total\_seconds() \% loop\_sec) の結果を代入する。
\item     self.current\_forecast\_time\_utc に lookup\_time の結果を代入する。
\item     idx に int(np.searchsorted(self.forecast\_times.values, np.datetime64(lookup\_time)) - 1) の結果を代入する。
\item     int(np.clip(idx, 0, bin\_count - 1)) を返す。
\item   idx に int(elapsed / max(self.dt, 0.001)) の結果を代入する。
\item   self.current\_forecast\_time\_utc に self.forecast\_start\_time + timedelta(seconds=elapsed) の結果を代入する。
\item   int(idx \% bin\_count) if mode == 'loop' else int(np.clip(idx, 0, bin\_count - 1)) を返す。
\item bin\_count に len(self.forecast\_times) if has\_time else len(self.df) の結果を代入する。
\item int(np.clip(self.k, 0, max(0, bin\_count - 1))) を返す。
            \end{enumerate}
            \paragraph{代表コード断片}
            \begin{lstlisting}[language=Python]
    def _current_bin_index(self) -> int:
        mode = str(self.forecast_time_mode).lower()
        has_time = len(getattr(self, 'forecast_times', [])) > 0
        if mode == 'auto':
            mode = 'absolute' if has_time else 'relative'

        now = datetime.now(timezone.utc) + timedelta(seconds=self.forecast_time_offset)
        lookup_time = now

        if mode == 'absolute' and has_time:
            t_series = self.forecast_times.values
            t_min = self.forecast_times[0]
            t_max = self.forecast_times[-1]
            if (now < t_min) or (now > t_max):
                if not getattr(self, '_forecast_warned_out_of_range', False):
                    self.get_logger().warn('Forecast time out of range; switching to relative indexing.')
                    self._forecast_warned_out_of_range = True
                mode = 'relative'
            else:
                idx = int(np.searchsorted(t_series, np.datetime64(now)) - 1)
                self.current_forecast_time_utc = now
                return int(np.clip(idx, 0, len(self.forecast_times) - 1))

        if mode in ('relative', 'loop') or not has_time:
            elapsed = (now - self.forecast_relative_wall_start).total_seconds()
            elapsed = max(0.0, elapsed)
            bin_count = len(self.forecast_times) if has_time else len(self.df)
            if bin_count == 0:
                return 0
            if has_time:
                lookup_time = self.forecast_start_time + timedelta(seconds=elapsed)
                if mode == 'loop' and len(self.forecast_times) > 1:
                    t_first = self.forecast_times[0].to_pydatetime()
                    t_last = self.forecast_times[-1].to_pydatetime()
                    loop_sec = max(1.0, (t_last - t_first).total_seconds())
...
\end{lstlisting}

\subsection{L830 関数 MPCNode.\_horizon\_data}
            \begin{itemize}[leftmargin=1.6em]
              \item 行範囲: L830--L870
              \item 所属: \path|MPCNode|
              \item 定義: \path|_horizon_data(self, k0: int)|
              \item docstring: 明示されていない。
              \item このブロックが直接呼ぶ主な関数・メソッド: \_forecast\_at\_time, append, bool, clip, datetime64, float, get, get\_parameter, getattr, int, isfinite, len
              \item 戻り値の要点: data / data
              \item 主なローカル代入名: N, Np, data, env, has\_time, i, is\_drive, j, limits, loop\_sec, mode, row, row\_match, s0\_km, s\_query\_km, speed\_guess\_kmh, t\_first, t\_last, t\_utc
              \item 読み取る主なインスタンス属性: self.Np, self.\_forecast\_at\_time, self.current\_forecast\_time\_utc, self.df, self.drive\_schedule, self.dt, self.forecast\_start\_time, self.forecast\_time\_mode, self.forecast\_times, self.get\_parameter, self.s\_km, self.s\_meas, self.v\_upper\_cmd
              \item 更新する主なインスタンス属性: 特になし。
              \item 明示的に送出する例外: 明示的raiseはない。
              \item 制御構造の規模: 条件分岐 7、ループ 1、try 0
            \end{itemize}
            \paragraph{この定義を読むためのPython構文}
            \begin{itemize}[leftmargin=1.6em]
            \item 第1引数selfは、このメソッドを呼んだインスタンス自身を受け取る慣習名である。
            \end{itemize}
            \paragraph{上から順の処理}
            \begin{enumerate}[leftmargin=1.8em]
            \item data に [] の結果を代入する。
\item has\_time に len(getattr(self, 'forecast\_times', [])) > 0 の結果を代入する。
\item N に len(self.forecast\_times) if has\_time else len(self.df) の結果を代入する。
\item mode に str(self.forecast\_time\_mode).lower() の結果を代入する。
\item 条件 mode == 'auto' を判定し、真なら内部処理を行う。
\item   mode に 'absolute' if has\_time else 'relative' の結果を代入する。
\item 条件 N <= 0 を判定し、真なら内部処理を行う。
\item   data を返す。
\item Np に max(1, self.Np) の結果を代入する。
\item s0\_km に float(self.s\_km) の結果を代入する。
\item 条件 bool(self.get\_parameter('use\_measured\_s').value) and np.isfinite(self.s\_meas) を判定し、真なら内部処理を行う。
\item   s0\_km に float(self.s\_meas) の結果を代入する。
\item speed\_guess\_kmh に max(0.0, float(self.v\_upper\_cmd)) の結果を代入する。
\item range(Np) を順に走査し、各要素を i に入れて処理する。
\item   条件 has\_time を判定し、真なら内部処理を行う。
\item     t\_utc に self.current\_forecast\_time\_utc + timedelta(seconds=self.dt * i) の結果を代入する。
\item     条件 mode == 'loop' and len(self.forecast\_times) > 1 を判定し、真なら内部処理を行う。
\item       t\_first に self.forecast\_times[0].to\_pydatetime() の結果を代入する。
\item       t\_last に self.forecast\_times[-1].to\_pydatetime() の結果を代入する。
\item       loop\_sec に max(1.0, (t\_last - t\_first).total\_seconds()) の結果を代入する。
\item       t\_utc に t\_first + timedelta(seconds=(t\_utc - t\_first).total\_seconds() \% loop\_sec) の結果を代入する。
\item     j に int(np.searchsorted(self.forecast\_times.values, np.datetime64(t\_utc)) - 1) の結果を代入する。
\item     j に int(np.clip(j, 0, N - 1)) の結果を代入する。
\item     row\_match に self.df.loc[self.df['time'] == self.forecast\_times[j]] の結果を代入する。
\item     row に row\_match.iloc[0] if not row\_match.empty else self.df.iloc[0] の結果を代入する。
\item     上の条件が偽の場合:
\item     j に (k0 + i) \% N if mode == 'loop' else min(k0 + i, N - 1) の結果を代入する。
\item     row に self.df.iloc[j] の結果を代入する。
\item     t\_utc に self.forecast\_start\_time + timedelta(seconds=self.dt * (k0 + i)) の結果を代入する。
\item   s\_query\_km に s0\_km + speed\_guess\_kmh * self.dt * i / 3600.0 の結果を代入する。
\item   is\_drive に True の結果を代入する。
\item   条件 self.drive\_schedule is not None を判定し、真なら内部処理を行う。
\item     limits に self.drive\_schedule.speed\_limits(t\_utc) の結果を代入する。
\item     条件 limits is not None and limits[1] <= 0.0 を判定し、真なら内部処理を行う。
\item       is\_drive に False の結果を代入する。
\item   env に self.\_forecast\_at\_time(t\_utc, s\_query\_km, drive=is\_drive) の結果を代入する。
\item   env['slope\_pct'] に float(row.get('slope\_pct', 0.0)) の結果を代入する。
\item   env['t\_utc'] に t\_utc の結果を代入する。
\item   data.append(...) を実行する。
\item data を返す。
            \end{enumerate}
            \paragraph{代表コード断片}
            \begin{lstlisting}[language=Python]
    def _horizon_data(self, k0: int):
        data = []
        has_time = len(getattr(self, 'forecast_times', [])) > 0
        N = len(self.forecast_times) if has_time else len(self.df)
        mode = str(self.forecast_time_mode).lower()
        if mode == 'auto':
            mode = 'absolute' if has_time else 'relative'
        if N <= 0:
            return data
        Np = max(1, self.Np)
        s0_km = float(self.s_km)
        if bool(self.get_parameter('use_measured_s').value) and np.isfinite(self.s_meas):
            s0_km = float(self.s_meas)
        speed_guess_kmh = max(0.0, float(self.v_upper_cmd))
        for i in range(Np):
            if has_time:
                t_utc = self.current_forecast_time_utc + timedelta(seconds=self.dt * i)
                if mode == 'loop' and len(self.forecast_times) > 1:
                    t_first = self.forecast_times[0].to_pydatetime()
                    t_last = self.forecast_times[-1].to_pydatetime()
                    loop_sec = max(1.0, (t_last - t_first).total_seconds())
                    t_utc = t_first + timedelta(seconds=(t_utc - t_first).total_seconds() % loop_sec)
                j = int(np.searchsorted(self.forecast_times.values, np.datetime64(t_utc)) - 1)
                j = int(np.clip(j, 0, N - 1))
                row_match = self.df.loc[self.df['time'] == self.forecast_times[j]]
                row = row_match.iloc[0] if not row_match.empty else self.df.iloc[0]
            else:
                j = (k0 + i) % N if mode == 'loop' else min(k0 + i, N - 1)
                row = self.df.iloc[j]
                t_utc = self.forecast_start_time + timedelta(seconds=self.dt * (k0 + i))
            s_query_km = s0_km + speed_guess_kmh * self.dt * i / 3600.0
            is_drive = True
            if self.drive_schedule is not None:
                limits = self.drive_schedule.speed_limits(t_utc)
                if limits is not None and limits[1] <= 0.0:
...
\end{lstlisting}

\subsection{L872 関数 MPCNode.\_forecast\_at\_time}
            \begin{itemize}[leftmargin=1.6em]
              \item 行範囲: L872--L946
              \item 所属: \path|MPCNode|
              \item 定義: \path!_forecast_at_time(self, t_utc: datetime, s_km: float | None = None, drive: bool = True, control_stop: bool = False) -> dict!
              \item docstring: 明示されていない。
              \item このブロックが直接呼ぶ主な関数・メソッド: \_route\_value, all, clip, datetime64, dict, float, get, getattr, int, interp\_forecast\_grid, isna, len
              \item 戻り値の要点: dict(G\_poa=(poa\_drive + tilt\_fraction * (poa\_stop\_ideal - poa\_drive)) * self.solar\_gain * gain, Tcell\_C=tcell\_drive + tilt\_fraction * (tcell\_stop\_ideal - tcell\_drive), Tamb\_C=float(row.get('Tamb\_C', 30.0)), headwind\_ms=float(row.get('headwind\_ms', 0.0)) if 'headwind\_ms' in row else 0.0, panel\_stop\_mode='drive' if drive else 'control\_stop' if control\_stop else 'camp\_or\_strategy', elevation\_m=self.\_route\_value(float(s\_km or 0.0), 'elev\_m', 0.0)) / dict(G\_poa=0.0, Tcell\_C=40.0, Tamb\_C=30.0, headwind\_ms=0.0, elevation\_m=self.\_route\_value(float(s\_km or 0.0), 'elev\_m', 0.0)) / dict(G\_poa=(poa\_drive + tilt\_fraction * (poa\_stop\_ideal - poa\_drive)) * self.solar\_gain * gain, Tcell\_C=tcell\_drive + tilt\_fraction * (tcell\_stop\_ideal - tcell\_drive), Tamb\_C=tamb, headwind\_ms=headwind, panel\_stop\_mode=stop\_mode, elevation\_m=self.\_route\_value(float(s\_km or 0.0), 'elev\_m', 0.0))
              \item 主なローカル代入名: elapsed, gain, ghi, has\_time, headwind, idx, poa\_drive, poa\_stop\_ideal, row, stop\_mode, t\_series, tamb, tcell\_drive, tcell\_stop\_ideal, tilt\_fraction
              \item 読み取る主なインスタンス属性: self.\_route\_value, self.control\_stop\_tilt\_fraction, self.df, self.dt, self.forecast\_grid, self.forecast\_start\_time, self.poa\_gain\_drive, self.poa\_gain\_stop, self.solar\_gain, self.stop\_tilt\_fraction
              \item 更新する主なインスタンス属性: 特になし。
              \item 明示的に送出する例外: 明示的raiseはない。
              \item 制御構造の規模: 条件分岐 4、ループ 0、try 0
            \end{itemize}
            \paragraph{この定義を読むためのPython構文}
            \begin{itemize}[leftmargin=1.6em]
            \item 第1引数selfは、このメソッドを呼んだインスタンス自身を受け取る慣習名である。
\item `->`以後は戻り値の型注釈であり、実行時の値を自動変換する命令ではない。
            \end{itemize}
            \paragraph{上から順の処理}
            \begin{enumerate}[leftmargin=1.8em]
            \item 条件 len(self.df) == 0 を判定し、真なら内部処理を行う。
\item   dict(G\_poa=0.0, Tcell\_C=40.0, Tamb\_C=30.0, headwind\_ms=0.0, elevation\_m=self.\_route\_value(float(s\_km or 0.0), 'elev\_m', 0.0)) を返す。
\item 条件 getattr(self, 'forecast\_grid', None) を判定し、真なら内部処理を行う。
\item   ghi に interp\_forecast\_grid(self.forecast\_grid, 'GHI', t\_utc, s\_km, 0.0) の結果を代入する。
\item   poa\_drive に interp\_forecast\_grid(self.forecast\_grid, 'POA\_drive', t\_utc, s\_km, ghi) の結果を代入する。
\item   poa\_stop\_ideal に interp\_forecast\_grid(self.forecast\_grid, 'POA\_stop\_ideal', t\_utc, s\_km, poa\_drive) の結果を代入する。
\item   tamb に interp\_forecast\_grid(self.forecast\_grid, 'Tamb\_C', t\_utc, s\_km, 30.0) の結果を代入する。
\item   tcell\_drive に interp\_forecast\_grid(self.forecast\_grid, 'Tcell\_drive\_C', t\_utc, s\_km, interp\_forecast\_grid(self.forecast\_grid, 'Tcell\_C', t\_utc, s\_km, tamb)) の結果を代入する。
\item   tcell\_stop\_ideal に interp\_forecast\_grid(self.forecast\_grid, 'Tcell\_stop\_ideal\_C', t\_utc, s\_km, tcell\_drive) の結果を代入する。
\item   headwind に interp\_forecast\_grid(self.forecast\_grid, 'headwind\_ms', t\_utc, s\_km, 0.0) の結果を代入する。
\item   条件 drive を判定し、真なら内部処理を行う。
\item     tilt\_fraction に 0.0 の結果を代入する。
\item     gain に self.poa\_gain\_drive の結果を代入する。
\item     stop\_mode に 'drive' の結果を代入する。
\item     上の条件が偽の場合:
\item     tilt\_fraction に self.control\_stop\_tilt\_fraction if control\_stop else self.stop\_tilt\_fraction の結果を代入する。
\item     gain に self.poa\_gain\_stop の結果を代入する。
\item     stop\_mode に 'control\_stop' if control\_stop else 'camp\_or\_strategy' の結果を代入する。
\item   dict(G\_poa=(poa\_drive + tilt\_fraction * (poa\_stop\_ideal - poa\_drive)) * self.solar\_gain * gain, Tcell\_C=tcell\_drive + tilt\_fraction * (tcell\_stop\_ideal - tcell\_drive), Tamb\_C=tamb, headwind\_ms=headwind, panel\_stop\_mode=stop\_mode, elevation\_m=self.\_route\_value(float(s\_km or 0.0), 'elev\_m', 0.0)) を返す。
\item has\_time に 'time' in self.df.columns and (not self.df['time'].isna().all()) の結果を代入する。
\item 条件 has\_time を判定し、真なら内部処理を行う。
\item   t\_series に self.df['time'].values の結果を代入する。
\item   idx に int(np.searchsorted(t\_series, np.datetime64(t\_utc)) - 1) の結果を代入する。
\item   idx に int(np.clip(idx, 0, len(self.df) - 1)) の結果を代入する。
\item   上の条件が偽の場合:
\item   elapsed に (t\_utc - self.forecast\_start\_time).total\_seconds() の結果を代入する。
\item   idx に int(np.clip(elapsed / max(self.dt, 0.001), 0, len(self.df) - 1)) の結果を代入する。
\item row に self.df.iloc[idx] の結果を代入する。
\item gain に self.poa\_gain\_drive if drive else self.poa\_gain\_stop の結果を代入する。
\item poa\_drive に float(row.get('POA\_drive', row.get('GHI', 0.0))) の結果を代入する。
\item poa\_stop\_ideal に float(row.get('POA\_stop\_ideal', poa\_drive)) の結果を代入する。
\item tilt\_fraction に 0.0 if drive else self.control\_stop\_tilt\_fraction if control\_stop else self.stop\_tilt\_fraction の結果を代入する。
\item tcell\_drive に float(row.get('Tcell\_drive\_C', row.get('Tcell\_C', 40.0))) の結果を代入する。
\item tcell\_stop\_ideal に float(row.get('Tcell\_stop\_ideal\_C', tcell\_drive)) の結果を代入する。
\item dict(G\_poa=(poa\_drive + tilt\_fraction * (poa\_stop\_ideal - poa\_drive)) * self.solar\_gain * gain, Tcell\_C=tcell\_drive + tilt\_fraction * (tcell\_stop\_ideal - tcell\_drive), Tamb\_C=float(row.get('Tamb\_C', 30.0)), headwind\_ms=float(row.get('headwind\_ms', 0.0)) if 'headwind\_ms' in row else 0.0, panel\_stop\_mode='drive' if drive else 'control\_stop' if control\_stop else 'camp\_or\_strategy', elevation\_m=self.\_route\_value(float(s\_km or 0.0), 'elev\_m', 0.0)) を返す。
            \end{enumerate}
            \paragraph{代表コード断片}
            \begin{lstlisting}[language=Python]
    def _forecast_at_time(
        self,
        t_utc: datetime,
        s_km: float | None = None,
        drive: bool = True,
        control_stop: bool = False,
    ) -> dict:
        if len(self.df) == 0:
            return dict(
                G_poa=0.0, Tcell_C=40.0, Tamb_C=30.0, headwind_ms=0.0,
                elevation_m=self._route_value(float(s_km or 0.0), 'elev_m', 0.0),
            )
        if getattr(self, 'forecast_grid', None):
            ghi = interp_forecast_grid(self.forecast_grid, 'GHI', t_utc, s_km, 0.0)
            poa_drive = interp_forecast_grid(self.forecast_grid, 'POA_drive', t_utc, s_km, ghi)
            poa_stop_ideal = interp_forecast_grid(
                self.forecast_grid, 'POA_stop_ideal', t_utc, s_km, poa_drive
            )
            tamb = interp_forecast_grid(self.forecast_grid, 'Tamb_C', t_utc, s_km, 30.0)
            tcell_drive = interp_forecast_grid(
                self.forecast_grid,
                'Tcell_drive_C',
                t_utc,
                s_km,
                interp_forecast_grid(self.forecast_grid, 'Tcell_C', t_utc, s_km, tamb),
            )
            tcell_stop_ideal = interp_forecast_grid(
                self.forecast_grid, 'Tcell_stop_ideal_C', t_utc, s_km, tcell_drive
            )
            headwind = interp_forecast_grid(self.forecast_grid, 'headwind_ms', t_utc, s_km, 0.0)
            if drive:
                tilt_fraction = 0.0
                gain = self.poa_gain_drive
                stop_mode = 'drive'
            else:
...
\end{lstlisting}

\subsection{L948 関数 MPCNode.\_sample\_plan\_segments}
            \begin{itemize}[leftmargin=1.6em]
              \item 行範囲: L948--L957
              \item 所属: \path|MPCNode|
              \item 定義: \path|_sample_plan_segments(self, dt_sample: float)|
              \item docstring: 明示されていない。
              \item このブロックが直接呼ぶ主な関数・メソッド: ceil, extend, float, int, max
              \item 戻り値の要点: samples / [] / [float(seg['v\_kmh']) for seg in self.v\_plan\_segments]
              \item 主なローカル代入名: n, samples, seg
              \item 読み取る主なインスタンス属性: self.v\_plan\_segments
              \item 更新する主なインスタンス属性: 特になし。
              \item 明示的に送出する例外: 明示的raiseはない。
              \item 制御構造の規模: 条件分岐 2、ループ 1、try 0
            \end{itemize}
            \paragraph{この定義を読むためのPython構文}
            \begin{itemize}[leftmargin=1.6em]
            \item 第1引数selfは、このメソッドを呼んだインスタンス自身を受け取る慣習名である。
\item 内包表記は、反復・条件・値生成を一つの式で記述してcollectionを作る。
            \end{itemize}
            \paragraph{上から順の処理}
            \begin{enumerate}[leftmargin=1.8em]
            \item 条件 not self.v\_plan\_segments を判定し、真なら内部処理を行う。
\item   [] を返す。
\item 条件 dt\_sample <= 0.0 を判定し、真なら内部処理を行う。
\item   [float(seg['v\_kmh']) for seg in self.v\_plan\_segments] を返す。
\item samples に [] の結果を代入する。
\item self.v\_plan\_segments を順に走査し、各要素を seg に入れて処理する。
\item   n に max(1, int(math.ceil(seg['dt\_sec'] / dt\_sample))) の結果を代入する。
\item   samples.extend(...) を実行する。
\item samples を返す。
            \end{enumerate}
            \paragraph{代表コード断片}
            \begin{lstlisting}[language=Python]
    def _sample_plan_segments(self, dt_sample: float):
        if not self.v_plan_segments:
            return []
        if dt_sample <= 0.0:
            return [float(seg['v_kmh']) for seg in self.v_plan_segments]
        samples = []
        for seg in self.v_plan_segments:
            n = max(1, int(math.ceil(seg['dt_sec'] / dt_sample)))
            samples.extend([float(seg['v_kmh'])] * n)
        return samples
\end{lstlisting}

\subsection{L959 関数 MPCNode.\_route\_value}
            \begin{itemize}[leftmargin=1.6em]
              \item 行範囲: L959--L968
              \item 所属: \path|MPCNode|
              \item 定義: \path|_route_value(self, s_km: float, field: str, default: float) -> float|
              \item docstring: 明示されていない。
              \item このブロックが直接呼ぶ主な関数・メソッド: float, interpolate\_profile, isfinite
              \item 戻り値の要点: float(default) / float(val) / float(default) / float(default)
              \item 主なローカル代入名: val
              \item 読み取る主なインスタンス属性: self.route\_profile
              \item 更新する主なインスタンス属性: 特になし。
              \item 明示的に送出する例外: 明示的raiseはない。
              \item 制御構造の規模: 条件分岐 2、ループ 0、try 1
            \end{itemize}
            \paragraph{この定義を読むためのPython構文}
            \begin{itemize}[leftmargin=1.6em]
            \item 第1引数selfは、このメソッドを呼んだインスタンス自身を受け取る慣習名である。
\item `->`以後は戻り値の型注釈であり、実行時の値を自動変換する命令ではない。
\item try文は例外が起き得る処理と、異常時または終了時の経路を分ける。
            \end{itemize}
            \paragraph{上から順の処理}
            \begin{enumerate}[leftmargin=1.8em]
            \item 条件 self.route\_profile is None を判定し、真なら内部処理を行う。
\item   float(default) を返す。
\item 例外処理を伴う try ブロックを実行する。
\item   val に float(interpolate\_profile(self.route\_profile, s\_km, field, default)) の結果を代入する。
\item   条件 not np.isfinite(val) を判定し、真なら内部処理を行う。
\item     float(default) を返す。
\item   float(val) を返す。
\item   Exceptionを捕捉した場合:
\item   float(default) を返す。
            \end{enumerate}
            \paragraph{代表コード断片}
            \begin{lstlisting}[language=Python]
    def _route_value(self, s_km: float, field: str, default: float) -> float:
        if self.route_profile is None:
            return float(default)
        try:
            val = float(interpolate_profile(self.route_profile, s_km, field, default))
            if not np.isfinite(val):
                return float(default)
            return float(val)
        except Exception:
            return float(default)
\end{lstlisting}

\subsection{L970 関数 MPCNode.\_route\_average}
            \begin{itemize}[leftmargin=1.6em]
              \item 行範囲: L970--L977
              \item 所属: \path|MPCNode|
              \item 定義: \path|_route_average(self, s0_km: float, s1_km: float, field: str, default: float) -> float|
              \item docstring: 明示されていない。
              \item このブロックが直接呼ぶ主な関数・メソッド: average\_profile, float, isfinite
              \item 戻り値の要点: float(default) / value if np.isfinite(value) else float(default) / float(default)
              \item 主なローカル代入名: value
              \item 読み取る主なインスタンス属性: self.route\_profile
              \item 更新する主なインスタンス属性: 特になし。
              \item 明示的に送出する例外: 明示的raiseはない。
              \item 制御構造の規模: 条件分岐 1、ループ 0、try 1
            \end{itemize}
            \paragraph{この定義を読むためのPython構文}
            \begin{itemize}[leftmargin=1.6em]
            \item 第1引数selfは、このメソッドを呼んだインスタンス自身を受け取る慣習名である。
\item `->`以後は戻り値の型注釈であり、実行時の値を自動変換する命令ではない。
\item try文は例外が起き得る処理と、異常時または終了時の経路を分ける。
            \end{itemize}
            \paragraph{上から順の処理}
            \begin{enumerate}[leftmargin=1.8em]
            \item 条件 self.route\_profile is None を判定し、真なら内部処理を行う。
\item   float(default) を返す。
\item 例外処理を伴う try ブロックを実行する。
\item   value に float(average\_profile(self.route\_profile, s0\_km, s1\_km, field, default)) の結果を代入する。
\item   value if np.isfinite(value) else float(default) を返す。
\item   Exceptionを捕捉した場合:
\item   float(default) を返す。
            \end{enumerate}
            \paragraph{代表コード断片}
            \begin{lstlisting}[language=Python]
    def _route_average(self, s0_km: float, s1_km: float, field: str, default: float) -> float:
        if self.route_profile is None:
            return float(default)
        try:
            value = float(average_profile(self.route_profile, s0_km, s1_km, field, default))
            return value if np.isfinite(value) else float(default)
        except Exception:
            return float(default)
\end{lstlisting}

\subsection{L979 関数 MPCNode.\_speed\_limit\_at}
            \begin{itemize}[leftmargin=1.6em]
              \item 行範囲: L979--L985
              \item 所属: \path|MPCNode|
              \item 定義: \path|_speed_limit_at(self, s_km: float, default_kmh: float) -> float|
              \item docstring: 明示されていない。
              \item このブロックが直接呼ぶ主な関数・メソッド: float, interpolate\_profile
              \item 戻り値の要点: float(default\_kmh) / float(interpolate\_profile(self.speed\_profile, s\_km, 'v\_max\_kmh', default\_kmh)) / float(default\_kmh)
              \item 主なローカル代入名: 特になし。
              \item 読み取る主なインスタンス属性: self.speed\_profile
              \item 更新する主なインスタンス属性: 特になし。
              \item 明示的に送出する例外: 明示的raiseはない。
              \item 制御構造の規模: 条件分岐 1、ループ 0、try 1
            \end{itemize}
            \paragraph{この定義を読むためのPython構文}
            \begin{itemize}[leftmargin=1.6em]
            \item 第1引数selfは、このメソッドを呼んだインスタンス自身を受け取る慣習名である。
\item `->`以後は戻り値の型注釈であり、実行時の値を自動変換する命令ではない。
\item try文は例外が起き得る処理と、異常時または終了時の経路を分ける。
            \end{itemize}
            \paragraph{上から順の処理}
            \begin{enumerate}[leftmargin=1.8em]
            \item 条件 self.speed\_profile is None を判定し、真なら内部処理を行う。
\item   float(default\_kmh) を返す。
\item 例外処理を伴う try ブロックを実行する。
\item   float(interpolate\_profile(self.speed\_profile, s\_km, 'v\_max\_kmh', default\_kmh)) を返す。
\item   Exceptionを捕捉した場合:
\item   float(default\_kmh) を返す。
            \end{enumerate}
            \paragraph{代表コード断片}
            \begin{lstlisting}[language=Python]
    def _speed_limit_at(self, s_km: float, default_kmh: float) -> float:
        if self.speed_profile is None:
            return float(default_kmh)
        try:
            return float(interpolate_profile(self.speed_profile, s_km, 'v_max_kmh', default_kmh))
        except Exception:
            return float(default_kmh)
\end{lstlisting}

\subsection{L987 関数 MPCNode.\_soc\_guard\_speed}
            \begin{itemize}[leftmargin=1.6em]
              \item 行範囲: L987--L1043
              \item 所属: \path|MPCNode|
              \item 定義: \path|_soc_guard_speed(self, v_kmh: float, s_km: float, d0: dict) -> float|
              \item docstring: 明示されていない。
              \item このブロックが直接呼ぶ主な関数・メソッド: \_route\_value, electrical\_balance, float, get, get\_parameter, lower, max, range, soc\_step, str, z\_next\_for
              \item 戻り値の要点: float(lo) / self.model.soc\_step(self.z, P\_pack, self.model.p.dt, current\_a=float(out['I']), Tbat\_C=self.Tb) / float(lo) / v\_kmh
              \item 主なローカル代入名: P\_pack, \_, headwind\_ms, hi, lo, mid, mode, out, slope\_pct, soc\_guard, target
              \item 読み取る主なインスタンス属性: self.Tb, self.\_route\_value, self.get\_parameter, self.model, self.route\_profile, self.z
              \item 更新する主なインスタンス属性: 特になし。
              \item 明示的に送出する例外: 明示的raiseはない。
              \item 制御構造の規模: 条件分岐 8、ループ 2、try 0
            \end{itemize}
            \paragraph{この定義を読むためのPython構文}
            \begin{itemize}[leftmargin=1.6em]
            \item 第1引数selfは、このメソッドを呼んだインスタンス自身を受け取る慣習名である。
\item `->`以後は戻り値の型注釈であり、実行時の値を自動変換する命令ではない。
            \end{itemize}
            \paragraph{上から順の処理}
            \begin{enumerate}[leftmargin=1.8em]
            \item mode に str(self.get\_parameter('soc\_guard\_mode').value).lower() の結果を代入する。
\item soc\_guard に float(self.get\_parameter('soc\_guard\_margin').value) の結果を代入する。
\item target に self.model.p.soc\_min + soc\_guard の結果を代入する。
\item slope\_pct に d0['slope\_pct'] の結果を代入する。
\item 条件 self.route\_profile is not None を判定し、真なら内部処理を行う。
\item   slope\_pct に self.\_route\_value(s\_km, 'slope\_pct', slope\_pct) の結果を代入する。
\item headwind\_ms に d0.get('headwind\_ms', 0.0) の結果を代入する。
\item 条件 self.route\_profile is not None を判定し、真なら内部処理を行う。
\item   headwind\_ms に self.\_route\_value(s\_km, 'headwind\_ms', headwind\_ms) の結果を代入する。
\item 関数 z\_next\_for を定義する。
\item 条件 self.z <= target を判定し、真なら内部処理を行う。
\item   条件 mode == 'stop' を判定し、真なら内部処理を行う。
\item     0.0 を返す。
\item   条件 mode != 'pv\_only' を判定し、真なら内部処理を行う。
\item     v\_kmh を返す。
\item   lo に 0.0 の結果を代入する。
\item   hi に max(0.0, float(v\_kmh)) の結果を代入する。
\item   range(20) を順に走査し、各要素を \_ に入れて処理する。
\item     mid に 0.5 * (lo + hi) の結果を代入する。
\item     条件 z\_next\_for(mid) < self.z を判定し、真なら内部処理を行う。
\item       hi に mid の結果を代入する。
\item       上の条件が偽の場合:
\item       lo に mid の結果を代入する。
\item   float(lo) を返す。
\item 条件 z\_next\_for(v\_kmh) >= target を判定し、真なら内部処理を行う。
\item   v\_kmh を返す。
\item lo に 0.0 の結果を代入する。
\item hi に max(0.0, float(v\_kmh)) の結果を代入する。
\item range(25) を順に走査し、各要素を \_ に入れて処理する。
\item   mid に 0.5 * (lo + hi) の結果を代入する。
\item   条件 z\_next\_for(mid) < target を判定し、真なら内部処理を行う。
\item     hi に mid の結果を代入する。
\item     上の条件が偽の場合:
\item     lo に mid の結果を代入する。
\item float(lo) を返す。
            \end{enumerate}
            \paragraph{代表コード断片}
            \begin{lstlisting}[language=Python]
    def _soc_guard_speed(self, v_kmh: float, s_km: float, d0: dict) -> float:
        mode = str(self.get_parameter('soc_guard_mode').value).lower()
        soc_guard = float(self.get_parameter('soc_guard_margin').value)
        target = self.model.p.soc_min + soc_guard

        slope_pct = d0['slope_pct']
        if self.route_profile is not None:
            slope_pct = self._route_value(s_km, 'slope_pct', slope_pct)
        headwind_ms = d0.get('headwind_ms', 0.0)
        if self.route_profile is not None:
            headwind_ms = self._route_value(s_km, 'headwind_ms', headwind_ms)

        def z_next_for(v_kmh_local: float) -> float:
            out = self.model.electrical_balance(
                v_kmh_local / 3.6, slope_pct, self.z, self.Tb,
                d0['G_poa'], d0['Tcell_C'], headwind_ms=headwind_ms,
                ambient_temp_c=d0.get('Tamb_C'),
                elevation_m=self._route_value(s_km, 'elev_m', d0.get('elevation_m', 0.0)),
            )
            P_pack = float(out['P_pack'])
            return self.model.soc_step(
                self.z,
                P_pack,
                self.model.p.dt,
                current_a=float(out['I']),
                Tbat_C=self.Tb,
            )

        # If already below target, apply guard mode
        if self.z <= target:
            if mode == 'stop':
                return 0.0
            if mode != 'pv_only':
                return v_kmh
            # find max speed such that P_pack <= 0
...
\end{lstlisting}

\subsection{L999 関数 MPCNode.\_soc\_guard\_speed.z\_next\_for}
            \begin{itemize}[leftmargin=1.6em]
              \item 行範囲: L999--L1013
              \item 所属: \path|MPCNode._soc_guard_speed|
              \item 定義: \path|z_next_for(v_kmh_local: float) -> float|
              \item docstring: 明示されていない。
              \item このブロックが直接呼ぶ主な関数・メソッド: \_route\_value, electrical\_balance, float, get, soc\_step
              \item 戻り値の要点: self.model.soc\_step(self.z, P\_pack, self.model.p.dt, current\_a=float(out['I']), Tbat\_C=self.Tb)
              \item 主なローカル代入名: P\_pack, out
              \item 読み取る主なインスタンス属性: self.Tb, self.\_route\_value, self.model, self.z
              \item 更新する主なインスタンス属性: 特になし。
              \item 明示的に送出する例外: 明示的raiseはない。
              \item 制御構造の規模: 条件分岐 0、ループ 0、try 0
            \end{itemize}
            \paragraph{この定義を読むためのPython構文}
            \begin{itemize}[leftmargin=1.6em]
            \item `->`以後は戻り値の型注釈であり、実行時の値を自動変換する命令ではない。
            \end{itemize}
            \paragraph{上から順の処理}
            \begin{enumerate}[leftmargin=1.8em]
            \item out に self.model.electrical\_balance(v\_kmh\_local / 3.6, slope\_pct, self.z, self.Tb, d0['G\_poa'], d0['Tcell\_C'], headwind\_ms=headwind\_ms, ambient\_temp\_c=d0.get('Tamb\_C'), elevation\_m=self.\_route\_value(s\_km, 'elev\_m', d0.get('elevation\_m', 0.0))) の結果を代入する。
\item P\_pack に float(out['P\_pack']) の結果を代入する。
\item self.model.soc\_step(self.z, P\_pack, self.model.p.dt, current\_a=float(out['I']), Tbat\_C=self.Tb) を返す。
            \end{enumerate}
            \paragraph{代表コード断片}
            \begin{lstlisting}[language=Python]
        def z_next_for(v_kmh_local: float) -> float:
            out = self.model.electrical_balance(
                v_kmh_local / 3.6, slope_pct, self.z, self.Tb,
                d0['G_poa'], d0['Tcell_C'], headwind_ms=headwind_ms,
                ambient_temp_c=d0.get('Tamb_C'),
                elevation_m=self._route_value(s_km, 'elev_m', d0.get('elevation_m', 0.0)),
            )
            P_pack = float(out['P_pack'])
            return self.model.soc_step(
                self.z,
                P_pack,
                self.model.p.dt,
                current_a=float(out['I']),
                Tbat_C=self.Tb,
            )
\end{lstlisting}

\subsection{L1045 関数 MPCNode.\_control\_stop\_catalog}
            \begin{itemize}[leftmargin=1.6em]
              \item 行範囲: L1045--L1063
              \item 所属: \path|MPCNode|
              \item 定義: \path|_control_stop_catalog(self)|
              \item docstring: 明示されていない。
              \item このブロックが直接呼ぶ主な関数・メソッド: append, enumerate, float, get, int, lower, max, sorted, str, strip
              \item 戻り値の要点: sorted(catalog, key=lambda item: item['s\_km'])
              \item 主なローカル代入名: catalog, dwell\_sec, explicit, kind, source\_index, stop
              \item 読み取る主なインスタンス属性: self.stops
              \item 更新する主なインスタンス属性: 特になし。
              \item 明示的に送出する例外: 明示的raiseはない。
              \item 制御構造の規模: 条件分岐 3、ループ 1、try 0
            \end{itemize}
            \paragraph{この定義を読むためのPython構文}
            \begin{itemize}[leftmargin=1.6em]
            \item 第1引数selfは、このメソッドを呼んだインスタンス自身を受け取る慣習名である。
\item lambdaは名前を付けずに短い関数オブジェクトを作る構文である。
            \end{itemize}
            \paragraph{上から順の処理}
            \begin{enumerate}[leftmargin=1.8em]
            \item catalog に [] の結果を代入する。
\item enumerate(self.stops) を順に走査し、各要素を (source\_index, stop) に入れて処理する。
\item   kind に str(stop.get('kind', '')).strip().lower() の結果を代入する。
\item   explicit に stop.get('is\_control\_stop', None) の結果を代入する。
\item   条件 explicit is False and kind != 'control\_stop' を判定し、真なら内部処理を行う。
\item     Continue 文を実行する。
\item   条件 explicit is None and kind in \{'trouble\_stop', 'unscheduled\_stop', 'strategy\_stop'\} を判定し、真なら内部処理を行う。
\item     Continue 文を実行する。
\item   dwell\_sec に max(0.0, float(stop.get('dwell\_sec', stop.get('dwell\_s', 0.0)))) の結果を代入する。
\item   条件 dwell\_sec <= 0.0 を判定し、真なら内部処理を行う。
\item     Continue 文を実行する。
\item   catalog.append(...) を実行する。
\item sorted(catalog, key=lambda item: item['s\_km']) を返す。
            \end{enumerate}
            \paragraph{代表コード断片}
            \begin{lstlisting}[language=Python]
    def _control_stop_catalog(self):
        catalog = []
        for source_index, stop in enumerate(self.stops):
            kind = str(stop.get('kind', '')).strip().lower()
            explicit = stop.get('is_control_stop', None)
            if explicit is False and kind != 'control_stop':
                continue
            if explicit is None and kind in {'trouble_stop', 'unscheduled_stop', 'strategy_stop'}:
                continue
            dwell_sec = max(0.0, float(stop.get('dwell_sec', stop.get('dwell_s', 0.0))))
            if dwell_sec <= 0.0:
                continue
            catalog.append({
                'source_index': int(source_index),
                's_km': float(stop.get('s_km', 0.0)),
                'dwell_sec': dwell_sec,
                'label': str(stop.get('label', f'control_stop_{source_index + 1}')),
            })
        return sorted(catalog, key=lambda item: item['s_km'])
\end{lstlisting}

\subsection{L1065 関数 MPCNode.\_update\_live\_control\_stop\_hold}
            \begin{itemize}[leftmargin=1.6em]
              \item 行範囲: L1065--L1113
              \item 所属: \path|MPCNode|
              \item 定義: \path|_update_live_control_stop_hold(self, s_km: float, v_kmh: float) -> bool|
              \item docstring: Enforce approach, standstill detection, and timed hold at declared control stops.
              \item このブロックが直接呼ぶ主な関数・メソッド: \_control\_stop\_catalog, add, float, get\_logger, get\_parameter, getattr, info, int, max, monotonic
              \item 戻り値の要点: False / False / True / True
              \item 主なローカル代入名: braking\_km, catalog, decel\_ms2, distance\_to\_stop\_km, now\_mono, stationary\_kmh, stop, stop\_index, tolerance\_km, v\_ms
              \item 読み取る主なインスタンス属性: self.\_control\_stop\_catalog, self.active\_control\_stop\_index, self.completed\_control\_stops, self.control\_stop\_end\_monotonic, self.get\_logger, self.get\_parameter
              \item 更新する主なインスタンス属性: self.active\_control\_stop\_index, self.control\_stop\_end\_monotonic, self.control\_stop\_position\_initialized
              \item 明示的に送出する例外: 明示的raiseはない。
              \item 制御構造の規模: 条件分岐 8、ループ 2、try 0
            \end{itemize}
            \paragraph{この定義を読むためのPython構文}
            \begin{itemize}[leftmargin=1.6em]
            \item 第1引数selfは、このメソッドを呼んだインスタンス自身を受け取る慣習名である。
\item `->`以後は戻り値の型注釈であり、実行時の値を自動変換する命令ではない。
            \end{itemize}
            \paragraph{上から順の処理}
            \begin{enumerate}[leftmargin=1.8em]
            \item catalog に self.\_control\_stop\_catalog() の結果を代入する。
\item 条件 not catalog を判定し、真なら内部処理を行う。
\item   False を返す。
\item now\_mono に time.monotonic() の結果を代入する。
\item 条件 self.active\_control\_stop\_index is not None を判定し、真なら内部処理を行う。
\item   条件 self.control\_stop\_end\_monotonic is None or now\_mono < self.control\_stop\_end\_monotonic を判定し、真なら内部処理を行う。
\item     True を返す。
\item   self.completed\_control\_stops.add(...) を実行する。
\item   self.get\_logger().info(...) を実行する。
\item   self.active\_control\_stop\_index に None の結果を代入する。
\item   self.control\_stop\_end\_monotonic に None の結果を代入する。
\item tolerance\_km に max(0.01, float(self.get\_parameter('control\_stop\_arrival\_tolerance\_km').value)) の結果を代入する。
\item 条件 not getattr(self, 'control\_stop\_position\_initialized', False) を判定し、真なら内部処理を行う。
\item   catalog を順に走査し、各要素を stop に入れて処理する。
\item     条件 stop['s\_km'] < float(s\_km) - tolerance\_km を判定し、真なら内部処理を行う。
\item       self.completed\_control\_stops.add(...) を実行する。
\item   self.control\_stop\_position\_initialized に True の結果を代入する。
\item stationary\_kmh に max(0.0, float(self.get\_parameter('control\_stop\_stationary\_speed\_kmh').value)) の結果を代入する。
\item decel\_ms2 に max(0.1, float(self.get\_parameter('control\_stop\_brake\_decel\_kmhps').value) / 3.6) の結果を代入する。
\item v\_ms に max(0.0, float(v\_kmh)) / 3.6 の結果を代入する。
\item braking\_km に v\_ms * v\_ms / (2.0 * decel\_ms2) / 1000.0 の結果を代入する。
\item braking\_km を Add で更新する。
\item catalog を順に走査し、各要素を stop に入れて処理する。
\item   stop\_index に int(stop['source\_index']) の結果を代入する。
\item   条件 stop\_index in self.completed\_control\_stops を判定し、真なら内部処理を行う。
\item     Continue 文を実行する。
\item   distance\_to\_stop\_km に float(stop['s\_km']) - float(s\_km) の結果を代入する。
\item   条件 distance\_to\_stop\_km > braking\_km を判定し、真なら内部処理を行う。
\item     False を返す。
\item   条件 distance\_to\_stop\_km <= tolerance\_km and float(v\_kmh) <= stationary\_kmh を判定し、真なら内部処理を行う。
\item     self.active\_control\_stop\_index に stop\_index の結果を代入する。
\item     self.control\_stop\_end\_monotonic に now\_mono + float(stop['dwell\_sec']) の結果を代入する。
\item     self.get\_logger().info(...) を実行する。
\item   True を返す。
\item False を返す。
            \end{enumerate}
            \paragraph{代表コード断片}
            \begin{lstlisting}[language=Python]
    def _update_live_control_stop_hold(self, s_km: float, v_kmh: float) -> bool:
        """Enforce approach, standstill detection, and timed hold at declared control stops."""
        catalog = self._control_stop_catalog()
        if not catalog:
            return False
        now_mono = time.monotonic()
        if self.active_control_stop_index is not None:
            if self.control_stop_end_monotonic is None or now_mono < self.control_stop_end_monotonic:
                return True
            self.completed_control_stops.add(int(self.active_control_stop_index))
            self.get_logger().info('Control-stop dwell completed; speed command released.')
            self.active_control_stop_index = None
            self.control_stop_end_monotonic = None

        tolerance_km = max(
            0.01, float(self.get_parameter('control_stop_arrival_tolerance_km').value)
        )
        if not getattr(self, 'control_stop_position_initialized', False):
            for stop in catalog:
                if stop['s_km'] < float(s_km) - tolerance_km:
                    self.completed_control_stops.add(int(stop['source_index']))
            self.control_stop_position_initialized = True

        stationary_kmh = max(
            0.0, float(self.get_parameter('control_stop_stationary_speed_kmh').value)
        )
        decel_ms2 = max(
            0.1,
            float(self.get_parameter('control_stop_brake_decel_kmhps').value) / 3.6,
        )
        v_ms = max(0.0, float(v_kmh)) / 3.6
        braking_km = v_ms * v_ms / (2.0 * decel_ms2) / 1000.0
        braking_km += max(0.0, float(self.get_parameter('control_stop_brake_margin_km').value))

        for stop in catalog:
...
\end{lstlisting}

\subsection{L1115 関数 MPCNode.\_dwell\_penalty}
            \begin{itemize}[leftmargin=1.6em]
              \item 行範囲: L1115--L1125
              \item 所属: \path|MPCNode|
              \item 定義: \path|_dwell_penalty(self, s_km: float, v_ms: float) -> float|
              \item docstring: 明示されていない。
              \item このブロックが直接呼ぶ主な関数・メソッド: abs, float, get, get\_parameter, max
              \item 戻り値の要点: pen
              \item 主なローカル代入名: dwell\_s, pen, s\_stop, st, vmax\_kmh, vmax\_ms, width\_km
              \item 読み取る主なインスタンス属性: self.get\_parameter, self.stops
              \item 更新する主なインスタンス属性: 特になし。
              \item 明示的に送出する例外: 明示的raiseはない。
              \item 制御構造の規模: 条件分岐 1、ループ 1、try 0
            \end{itemize}
            \paragraph{この定義を読むためのPython構文}
            \begin{itemize}[leftmargin=1.6em]
            \item 第1引数selfは、このメソッドを呼んだインスタンス自身を受け取る慣習名である。
\item `->`以後は戻り値の型注釈であり、実行時の値を自動変換する命令ではない。
            \end{itemize}
            \paragraph{上から順の処理}
            \begin{enumerate}[leftmargin=1.8em]
            \item vmax\_kmh に float(self.get\_parameter('v\_max\_kmh').value) の結果を代入する。
\item vmax\_ms に vmax\_kmh / 3.6 の結果を代入する。
\item pen に 0.0 の結果を代入する。
\item self.stops を順に走査し、各要素を st に入れて処理する。
\item   s\_stop に float(st.get('s\_km', 0.0)) の結果を代入する。
\item   dwell\_s に float(st.get('dwell\_sec', st.get('dwell\_s', 0.0))) の結果を代入する。
\item   width\_km に max(0.05, dwell\_s * vmax\_ms / 1000.0 * 0.5) の結果を代入する。
\item   条件 abs(s\_km - s\_stop) <= width\_km を判定し、真なら内部処理を行う。
\item     pen を Add で更新する。
\item pen を返す。
            \end{enumerate}
            \paragraph{代表コード断片}
            \begin{lstlisting}[language=Python]
    def _dwell_penalty(self, s_km: float, v_ms: float) -> float:
        vmax_kmh = float(self.get_parameter('v_max_kmh').value)
        vmax_ms = vmax_kmh / 3.6
        pen = 0.0
        for st in self.stops:
            s_stop = float(st.get('s_km', 0.0))
            dwell_s = float(st.get('dwell_sec', st.get('dwell_s', 0.0)))
            width_km = max(0.05, (dwell_s * vmax_ms) / 1000.0 * 0.5)
            if abs(s_km - s_stop) <= width_km:
                pen += 1.0e5 * (v_ms ** 2)
        return pen
\end{lstlisting}

\subsection{L1127 関数 MPCNode.\_mpc\_solve\_solar}
            \begin{itemize}[leftmargin=1.6em]
              \item 行範囲: L1127--L1289
              \item 所属: \path|MPCNode|
              \item 定義: \path|_mpc_solve_solar(self, data)|
              \item docstring: 明示されていない。
              \item このブロックが直接呼ぶ主な関数・メソッド: \_dwell\_penalty, \_route\_value, \_speed\_limit\_at, abs, all, array, bool, clip, dict, electrical\_balance, float, get
              \item 戻り値の要点: (float(v\_seq\_kmh[0]), [float(v) for v in v\_seq\_kmh]) / (self.v\_cmd, [self.v\_cmd]) / x * x / J
              \item 主なローカル代入名: I, J, Np, P\_pack, Tb, Tb\_next, V, aux\_power\_w, bounds, d, dv, dv\_max\_kmhps, dv\_max\_msps, headwind\_ms, inertial\_power\_w, k, lb, limits, loss\_int, out
              \item 読み取る主なインスタンス属性: self.Tb, self.\_dwell\_penalty, self.\_route\_value, self.\_speed\_limit\_at, self.drive\_schedule, self.get\_logger, self.get\_parameter, self.model, self.race\_km, self.route\_profile, self.s\_km, self.s\_meas, self.soc\_band, self.soc\_day\_end\_max, self.soc\_finish\_target, self.soc\_finish\_tol, self.soc\_target, self.upper\_vmin\_kmh, self.v\_cmd, self.v\_now, self.w\_soc\_band, self.w\_soc\_day\_max, self.w\_soc\_progress, self.w\_soc\_target
              \item 更新する主なインスタンス属性: self.last\_upper\_solve\_ok
              \item 明示的に送出する例外: 明示的raiseはない。
              \item 制御構造の規模: 条件分岐 20、ループ 1、try 0
            \end{itemize}
            \paragraph{この定義を読むためのPython構文}
            \begin{itemize}[leftmargin=1.6em]
            \item 第1引数selfは、このメソッドを呼んだインスタンス自身を受け取る慣習名である。
\item 内包表記は、反復・条件・値生成を一つの式で記述してcollectionを作る。
            \end{itemize}
            \paragraph{上から順の処理}
            \begin{enumerate}[leftmargin=1.8em]
            \item p に self.model.p の結果を代入する。
\item Np に len(data) の結果を代入する。
\item 条件 Np <= 0 を判定し、真なら内部処理を行う。
\item   (self.v\_cmd, [self.v\_cmd]) を返す。
\item v0\_guess に self.v\_cmd / 3.6 の結果を代入する。
\item 条件 bool(self.get\_parameter('use\_measured\_speed').value) and np.isfinite(self.v\_now) を判定し、真なら内部処理を行う。
\item   v0\_guess に float(self.v\_now) / 3.6 の結果を代入する。
\item x0 に np.array([v0\_guess] * Np, dtype=float) の結果を代入する。
\item lb に np.zeros(Np, dtype=float) の結果を代入する。
\item v\_max\_kmh に float(self.get\_parameter('v\_max\_kmh').value) の結果を代入する。
\item v\_min\_solver に max(0.1, float(self.upper\_vmin\_kmh)) の結果を代入する。
\item ub に np.ones(Np, dtype=float) * (v\_max\_kmh / 3.6) の結果を代入する。
\item term\_soc\_min に float(self.get\_parameter('terminal\_soc\_min').value) の結果を代入する。
\item w\_dv に float(self.get\_parameter('w\_dv').value) の結果を代入する。
\item w\_dv\_limit に float(self.get\_parameter('w\_dv\_limit').value) の結果を代入する。
\item dv\_max\_kmhps に float(self.get\_parameter('dv\_max\_kmhps').value) の結果を代入する。
\item dv\_max\_msps に dv\_max\_kmhps / 3.6 の結果を代入する。
\item w\_T に float(self.get\_parameter('w\_T').value) の結果を代入する。
\item w\_speed\_limit に float(self.get\_parameter('w\_speed\_limit').value) の結果を代入する。
\item w\_drive\_window に float(self.get\_parameter('w\_drive\_window').value) の結果を代入する。
\item w\_current に float(self.get\_parameter('w\_current').value) の結果を代入する。
\item soc\_target に float(self.soc\_target) の結果を代入する。
\item soc\_band に float(self.soc\_band) の結果を代入する。
\item w\_soc\_target に float(self.w\_soc\_target) の結果を代入する。
\item w\_soc\_band に float(self.w\_soc\_band) の結果を代入する。
\item soc\_day\_end\_max に float(self.soc\_day\_end\_max) の結果を代入する。
\item w\_soc\_day\_max に float(self.w\_soc\_day\_max) の結果を代入する。
\item soc\_finish\_target に float(self.soc\_finish\_target) の結果を代入する。
\item soc\_finish\_tol に float(self.soc\_finish\_tol) の結果を代入する。
\item w\_soc\_progress に float(self.w\_soc\_progress) の結果を代入する。
\item w\_soc\_terminal に float(self.w\_soc\_terminal) の結果を代入する。
\item race\_km に float(self.race\_km) の結果を代入する。
\item z\_start に float(self.z) の結果を代入する。
\item 関数 quad\_penalty を定義する。
\item 関数 cost を定義する。
\item bounds に list(zip(lb, ub)) の結果を代入する。
\item res に minimize(cost, x0, method='L-BFGS-B', bounds=bounds, options=dict(maxiter=150)) の結果を代入する。
\item 条件 np.all(np.isfinite(res.x)) を判定し、真なら内部処理を行う。
\item   self.last\_upper\_solve\_ok に bool(res.success) の結果を代入する。
\item   v\_seq に res.x の結果を代入する。
\item   条件 not res.success を判定し、真なら内部処理を行う。
\item     self.get\_logger().warn(...) を実行する。
\item   上の条件が偽の場合:
\item   self.last\_upper\_solve\_ok に False の結果を代入する。
\item   self.get\_logger().warn(...) を実行する。
\item   v\_seq に x0 の結果を代入する。
\item v\_seq\_kmh に np.clip(v\_seq * 3.6, 0.0, v\_max\_kmh) の結果を代入する。
\item (float(v\_seq\_kmh[0]), [float(v) for v in v\_seq\_kmh]) を返す。
            \end{enumerate}
            \paragraph{代表コード断片}
            \begin{lstlisting}[language=Python]
    def _mpc_solve_solar(self, data):
        p = self.model.p
        Np = len(data)
        if Np <= 0:
            return self.v_cmd, [self.v_cmd]

        v0_guess = self.v_cmd / 3.6
        if bool(self.get_parameter('use_measured_speed').value) and np.isfinite(self.v_now):
            v0_guess = float(self.v_now) / 3.6
        x0 = np.array([v0_guess] * Np, dtype=float)  # m/s
        lb = np.zeros(Np, dtype=float)
        v_max_kmh = float(self.get_parameter('v_max_kmh').value)
        v_min_solver = max(0.1, float(self.upper_vmin_kmh))
        ub = np.ones(Np, dtype=float) * (v_max_kmh / 3.6)

        term_soc_min = float(self.get_parameter('terminal_soc_min').value)
        w_dv = float(self.get_parameter('w_dv').value)
        w_dv_limit = float(self.get_parameter('w_dv_limit').value)
        dv_max_kmhps = float(self.get_parameter('dv_max_kmhps').value)
        dv_max_msps = dv_max_kmhps / 3.6
        w_T = float(self.get_parameter('w_T').value)
        w_speed_limit = float(self.get_parameter('w_speed_limit').value)
        w_drive_window = float(self.get_parameter('w_drive_window').value)
        w_current = float(self.get_parameter('w_current').value)
        soc_target = float(self.soc_target)
        soc_band = float(self.soc_band)
        w_soc_target = float(self.w_soc_target)
        w_soc_band = float(self.w_soc_band)
        soc_day_end_max = float(self.soc_day_end_max)
        w_soc_day_max = float(self.w_soc_day_max)
        soc_finish_target = float(self.soc_finish_target)
        soc_finish_tol = float(self.soc_finish_tol)
        w_soc_progress = float(self.w_soc_progress)
        w_soc_terminal = float(self.w_soc_terminal)
        race_km = float(self.race_km)
...
\end{lstlisting}

\subsection{L1164 関数 MPCNode.\_mpc\_solve\_solar.quad\_penalty}
            \begin{itemize}[leftmargin=1.6em]
              \item 行範囲: L1164--L1169
              \item 所属: \path|MPCNode._mpc_solve_solar|
              \item 定義: \path|quad_penalty(x, cap = 1000.0)|
              \item docstring: 明示されていない。
              \item このブロックが直接呼ぶ主な関数・メソッド: 特に明示されていない。
              \item 戻り値の要点: x * x / 0.0
              \item 主なローカル代入名: x
              \item 読み取る主なインスタンス属性: 特になし。
              \item 更新する主なインスタンス属性: 特になし。
              \item 明示的に送出する例外: 明示的raiseはない。
              \item 制御構造の規模: 条件分岐 2、ループ 0、try 0
            \end{itemize}
            \paragraph{この定義を読むためのPython構文}
            \begin{itemize}[leftmargin=1.6em]
            \item 特別な構文上の補足は少ない。
            \end{itemize}
            \paragraph{上から順の処理}
            \begin{enumerate}[leftmargin=1.8em]
            \item 条件 x <= 0.0 を判定し、真なら内部処理を行う。
\item   0.0 を返す。
\item 条件 x > cap を判定し、真なら内部処理を行う。
\item   x に cap の結果を代入する。
\item x * x を返す。
            \end{enumerate}
            \paragraph{代表コード断片}
            \begin{lstlisting}[language=Python]
        def quad_penalty(x, cap=1.0e3):
            if x <= 0.0:
                return 0.0
            if x > cap:
                x = cap
            return x * x
\end{lstlisting}

\subsection{L1171 関数 MPCNode.\_mpc\_solve\_solar.cost}
            \begin{itemize}[leftmargin=1.6em]
              \item 行範囲: L1171--L1275
              \item 所属: \path|MPCNode._mpc_solve_solar|
              \item 定義: \path|cost(v)|
              \item docstring: 明示されていない。
              \item このブロックが直接呼ぶ主な関数・メソッド: \_dwell\_penalty, \_route\_value, \_speed\_limit\_at, abs, bool, electrical\_balance, float, get, get\_parameter, is\_drive\_time, isfinite, max
              \item 戻り値の要点: J
              \item 主なローカル代入名: I, J, P\_pack, Tb, Tb\_next, V, aux\_power\_w, d, dv, headwind\_ms, inertial\_power\_w, k, limits, loss\_int, out, prog, s\_km, slope\_pct, soc\_line, t\_next
              \item 読み取る主なインスタンス属性: self.Tb, self.\_dwell\_penalty, self.\_route\_value, self.\_speed\_limit\_at, self.drive\_schedule, self.get\_parameter, self.model, self.route\_profile, self.s\_km, self.s\_meas, self.z
              \item 更新する主なインスタンス属性: 特になし。
              \item 明示的に送出する例外: 明示的raiseはない。
              \item 制御構造の規模: 条件分岐 14、ループ 1、try 0
            \end{itemize}
            \paragraph{この定義を読むためのPython構文}
            \begin{itemize}[leftmargin=1.6em]
            \item 特別な構文上の補足は少ない。
            \end{itemize}
            \paragraph{上から順の処理}
            \begin{enumerate}[leftmargin=1.8em]
            \item z に float(self.z) の結果を代入する。
\item Tb に float(self.Tb) の結果を代入する。
\item s\_km に float(self.s\_km) の結果を代入する。
\item 条件 bool(self.get\_parameter('use\_measured\_s').value) and np.isfinite(self.s\_meas) を判定し、真なら内部処理を行う。
\item   s\_km に float(self.s\_meas) の結果を代入する。
\item v\_prev に float(v0\_guess) の結果を代入する。
\item J に 0.0 の結果を代入する。
\item range(Np) を順に走査し、各要素を k に入れて処理する。
\item   d に data[k] の結果を代入する。
\item   v\_k に float(v[k]) の結果を代入する。
\item   slope\_pct に d['slope\_pct'] の結果を代入する。
\item   条件 self.route\_profile is not None を判定し、真なら内部処理を行う。
\item     slope\_pct に self.\_route\_value(s\_km, 'slope\_pct', slope\_pct) の結果を代入する。
\item   headwind\_ms に d.get('headwind\_ms', 0.0) の結果を代入する。
\item   条件 self.route\_profile is not None を判定し、真なら内部処理を行う。
\item     headwind\_ms に self.\_route\_value(s\_km, 'headwind\_ms', headwind\_ms) の結果を代入する。
\item   aux\_power\_w に None の結果を代入する。
\item   条件 self.drive\_schedule is not None and 't\_utc' in d and (not self.drive\_schedule.is\_drive\_time(d['t\_utc'])) を判定し、真なら内部処理を行う。
\item     aux\_power\_w に self.model.scheduled\_auxiliary\_power(is\_driving=False, irradiance\_wm2=d.get('G\_poa', 0.0)) の結果を代入する。
\item   inertial\_power\_w に 0.5 * p.m * (v\_k * v\_k - v\_prev * v\_prev) / max(p.dt, 1e-09) の結果を代入する。
\item   out に self.model.electrical\_balance(v\_k, slope\_pct, z, Tb, d['G\_poa'], d['Tcell\_C'], headwind\_ms=headwind\_ms, aux\_power\_w=aux\_power\_w, inertial\_power\_w=inertial\_power\_w, ambient\_temp\_c=d.get('Tamb\_C'), elevation\_m=d.get('elevation\_m', self.\_route\_value(s\_km, 'elev\_m', 0.0))) の結果を代入する。
\item   I に float(out['I']) の結果を代入する。
\item   V に float(out['V']) の結果を代入する。
\item   P\_pack に float(out['P\_pack']) の結果を代入する。
\item   loss\_int に float(out['losses\_int']) の結果を代入する。
\item   z\_next に self.model.soc\_step(z, P\_pack, p.dt, current\_a=I, Tbat\_C=Tb) の結果を代入する。
\item   Tb\_next に Tb + p.dt / 1800.0 * (d['Tamb\_C'] - Tb) + loss\_int * p.dt / 50000.0 の結果を代入する。
\item   s\_km に s\_km + v\_k * (p.dt / 1000.0) の結果を代入する。
\item   J を Add で更新する。
\item   J を Add で更新する。
\item   条件 self.drive\_schedule is not None and soc\_day\_end\_max > 0.0 and ('t\_utc' in d) を判定し、真なら内部処理を行う。
\item     t\_next に d['t\_utc'] + timedelta(seconds=p.dt) の結果を代入する。
\item     条件 self.drive\_schedule.is\_drive\_time(d['t\_utc']) and (not self.drive\_schedule.is\_drive\_time(t\_next)) を判定し、真なら内部処理を行う。
\item       J を Add で更新する。
\item   条件 soc\_finish\_target > 0.0 を判定し、真なら内部処理を行う。
\item     prog に max(0.0, min(1.0, s\_km / max(race\_km, 1.0))) の結果を代入する。
\item     soc\_line に z\_start + (soc\_finish\_target - z\_start) * prog の結果を代入する。
\item     条件 z\_next > soc\_line + soc\_finish\_tol を判定し、真なら内部処理を行う。
\item       J を Add で更新する。
\item     条件 z\_next < soc\_line - soc\_finish\_tol を判定し、真なら内部処理を行う。
\item       J を Add で更新する。
\item   dv に (v\_k - v\_prev) / max(p.dt, 0.001) の結果を代入する。
\item   条件 dv\_max\_msps > 0.0 を判定し、真なら内部処理を行う。
\item     J を Add で更新する。
\item   条件 self.drive\_schedule is not None and 't\_utc' in d を判定し、真なら内部処理を行う。
\item     limits に self.drive\_schedule.speed\_limits(d['t\_utc']) の結果を代入する。
\item     条件 limits is not None を判定し、真なら内部処理を行う。
\item       (vmin\_kmh, vmax\_kmh) に limits の結果を代入する。
\item       vmin\_ms に vmin\_kmh / 3.6 の結果を代入する。
\item       vmax\_ms に vmax\_kmh / 3.6 の結果を代入する。
\item       J を Add で更新する。
\item       J を Add で更新する。
\item   vmax\_local に self.\_speed\_limit\_at(s\_km, self.get\_parameter('v\_max\_kmh').value) の結果を代入する。
\item   条件 vmax\_local < self.get\_parameter('v\_max\_kmh').value を判定し、真なら内部処理を行う。
\item     J を Add で更新する。
\item   J を Add で更新する。
\item   J を Add で更新する。
\item   J を Add で更新する。
\item   J を Add で更新する。
\item   J を Add で更新する。
\item   J を Add で更新する。
\item   J を Add で更新する。
\item   J を Add で更新する。
\item   J を Add で更新する。
\item   (z, Tb) に (z\_next, Tb\_next) の結果を代入する。
\item   v\_prev に v\_k の結果を代入する。
\item J を Add で更新する。
\item 条件 soc\_finish\_target > 0.0 を判定し、真なら内部処理を行う。
\item   J を Add で更新する。
\item J を返す。
            \end{enumerate}
            \paragraph{代表コード断片}
            \begin{lstlisting}[language=Python]
        def cost(v):
            z = float(self.z)
            Tb = float(self.Tb)
            s_km = float(self.s_km)
            if bool(self.get_parameter('use_measured_s').value) and np.isfinite(self.s_meas):
                s_km = float(self.s_meas)
            v_prev = float(v0_guess)
            J = 0.0

            for k in range(Np):
                d = data[k]
                v_k = float(v[k])
                slope_pct = d['slope_pct']
                if self.route_profile is not None:
                    slope_pct = self._route_value(s_km, 'slope_pct', slope_pct)
                headwind_ms = d.get('headwind_ms', 0.0)
                if self.route_profile is not None:
                    headwind_ms = self._route_value(s_km, 'headwind_ms', headwind_ms)
                aux_power_w = None
                if self.drive_schedule is not None and 't_utc' in d and not self.drive_schedule.is_drive_time(d['t_utc']):
                    aux_power_w = self.model.scheduled_auxiliary_power(
                        is_driving=False,
                        irradiance_wm2=d.get('G_poa', 0.0),
                    )
                inertial_power_w = 0.5 * p.m * (v_k * v_k - v_prev * v_prev) / max(p.dt, 1.0e-9)
                out = self.model.electrical_balance(
                    v_k,
                    slope_pct,
                    z,
                    Tb,
                    d['G_poa'],
                    d['Tcell_C'],
                    headwind_ms=headwind_ms,
                    aux_power_w=aux_power_w,
                    inertial_power_w=inertial_power_w,
...
\end{lstlisting}

\subsection{L1291 関数 MPCNode.\_mpc\_solve\_solar\_distance}
            \begin{itemize}[leftmargin=1.6em]
              \item 行範囲: L1291--L1814
              \item 所属: \path|MPCNode|
              \item 定義: \path|_mpc_solve_solar_distance(self, t0_utc: datetime, s0_km: float, x0 = None)|
              \item docstring: 明示されていない。
              \item このブロックが直接呼ぶ主な関数・メソッド: \_control\_stop\_catalog, \_forecast\_at\_time, \_route\_average, \_route\_value, \_speed\_limit\_at, abs, absolute\_control\_distances, append, apply\_control\_stop\_at, array, asarray, bool
              \item 戻り値の要点: (v0\_kmh, segments, u\_seq) / (self.v\_cmd, [\{'v\_kmh': float(self.v\_cmd), 'dt\_sec': float(p.dt)\}], [float(self.v\_cmd)]) / (1.0 - alpha) * u\_vec[idx] + alpha * u\_vec[idx\_next] / np.array(u\_seed, dtype=float)
              \item 主なローカル代入名: I, J, Nc, Np, P\_mech\_wheel, P\_pack, P\_pv, Tb, Tb\_after, Tb\_next, Tb\_seed, V, \_, alpha, base\_edges, best\_score, best\_v, bounds, candidate\_grid, constraint
              \item 読み取る主なインスタンス属性: self.Tb, self.\_control\_stop\_catalog, self.\_forecast\_at\_time, self.\_route\_average, self.\_route\_value, self.\_speed\_limit\_at, self.drive\_schedule, self.get\_logger, self.get\_parameter, self.initial\_upper\_policy, self.last\_upper\_solve\_ok, self.model, self.race\_km, self.soc\_day\_end\_max, self.soc\_day\_end\_target, self.soc\_finish\_target, self.soc\_finish\_tol, self.upper\_adaptive\_growth, self.upper\_adaptive\_max\_ds\_km, self.upper\_adaptive\_min\_ds\_km, self.upper\_cost\_cfg, self.upper\_ctrl\_km, self.upper\_day\_end\_soc\_min, self.upper\_ds\_km
              \item 更新する主なインスタンス属性: self.last\_upper\_solve\_ok, self.upper\_plan\_ctrl\_s
              \item 明示的に送出する例外: 明示的raiseはない。
              \item 制御構造の規模: 条件分岐 25、ループ 7、try 0
            \end{itemize}
            \paragraph{この定義を読むためのPython構文}
            \begin{itemize}[leftmargin=1.6em]
            \item 第1引数selfは、このメソッドを呼んだインスタンス自身を受け取る慣習名である。
\item 内包表記は、反復・条件・値生成を一つの式で記述してcollectionを作る。
            \end{itemize}
            \paragraph{上から順の処理}
            \begin{enumerate}[leftmargin=1.8em]
            \item p に self.model.p の結果を代入する。
\item horizon に build\_upper\_distance\_horizon(mode=self.upper\_horizon\_mode, s0\_km=s0\_km, race\_km=self.race\_km, ds\_km=self.upper\_ds\_km, horizon\_km=self.upper\_horizon\_km, max\_steps=self.upper\_max\_steps, ctrl\_km=self.upper\_ctrl\_km, adaptive\_min\_ds\_km=self.upper\_adaptive\_min\_ds\_km, adaptive\_max\_ds\_km=self.upper\_adaptive\_max\_ds\_km, adaptive\_growth=self.upper\_adaptive\_growth) の結果を代入する。
\item ds\_seq に np.array(horizon.ds\_seq\_km, dtype=float) の結果を代入する。
\item seg\_s に np.array(horizon.seg\_s\_km, dtype=float) の結果を代入する。
\item Np に int(len(ds\_seq)) の結果を代入する。
\item 条件 Np <= 0 を判定し、真なら内部処理を行う。
\item   (self.v\_cmd, [\{'v\_kmh': float(self.v\_cmd), 'dt\_sec': float(p.dt)\}], [float(self.v\_cmd)]) を返す。
\item control\_stops に self.\_control\_stop\_catalog() の結果を代入する。
\item horizon\_end\_km に float(s0\_km + np.sum(ds\_seq)) の結果を代入する。
\item extra\_edges に [float(stop['s\_km']) for stop in control\_stops if float(s0\_km) + 1e-09 < float(stop['s\_km']) < horizon\_end\_km - 1e-09] の結果を代入する。
\item 条件 extra\_edges を判定し、真なら内部処理を行う。
\item   base\_edges に np.concatenate(([float(s0\_km)], float(s0\_km) + np.cumsum(ds\_seq))) の結果を代入する。
\item   route\_edges に np.array(sorted(set(np.round(np.concatenate((base\_edges, extra\_edges)), 9))), dtype=float) の結果を代入する。
\item   ds\_seq に np.diff(route\_edges) の結果を代入する。
\item   seg\_s に route\_edges[:-1] - float(s0\_km) の結果を代入する。
\item   Np に int(len(ds\_seq)) の結果を代入する。
\item v\_max\_kmh に float(self.get\_parameter('v\_max\_kmh').value) の結果を代入する。
\item v\_min\_solver に max(0.1, float(self.upper\_vmin\_kmh)) の結果を代入する。
\item v0 に float(self.v\_cmd) の結果を代入する。
\item 条件 bool(self.get\_parameter('use\_measured\_speed').value) and np.isfinite(self.v\_now) を判定し、真なら内部処理を行う。
\item   v0 に float(self.v\_now) の結果を代入する。
\item ctrl\_s に np.array(horizon.ctrl\_s\_km, dtype=float) の結果を代入する。
\item ctrl\_s\_abs に absolute\_control\_distances(s0\_km, ctrl\_s) の結果を代入する。
\item Nc に int(len(ctrl\_s)) の結果を代入する。
\item previous\_ctrl\_s に getattr(self, 'upper\_plan\_ctrl\_s', None) の結果を代入する。
\item 条件 x0 is not None and previous\_ctrl\_s is not None and (len(x0) == len(previous\_ctrl\_s)) and (len(previous\_ctrl\_s) >= 1) を判定し、真なら内部処理を行う。
\item   x0 に shift\_upper\_policy\_warm\_start(previous\_ctrl\_s, x0, ctrl\_s\_abs, minimum\_speed\_kmh=v\_min\_solver, maximum\_speed\_kmh=v\_max\_kmh) の結果を代入する。
\item   上の条件が偽の場合:
\item   条件 x0 is not None and len(x0) == Nc を判定し、真なら内部処理を行う。
\item     x0 に np.array(x0, dtype=float) の結果を代入する。
\item     上の条件が偽の場合:
\item     条件 self.initial\_upper\_policy is not None を判定し、真なら内部処理を行う。
\item       x0 に interpolate\_upper\_policy(self.initial\_upper\_policy, ctrl\_s\_abs, minimum\_speed\_kmh=v\_min\_solver, maximum\_speed\_kmh=v\_max\_kmh) の結果を代入する。
\item       上の条件が偽の場合:
\item       x0 に np.full(Nc, float(np.clip(v0, v\_min\_solver, v\_max\_kmh)), dtype=float) の結果を代入する。
\item x0 に np.clip(np.asarray(x0, dtype=float), v\_min\_solver, v\_max\_kmh) の結果を代入する。
\item self.upper\_plan\_ctrl\_s に ctrl\_s\_abs.copy() の結果を代入する。
\item bounds に [(v\_min\_solver, v\_max\_kmh)] * Nc の結果を代入する。
\item idx に np.searchsorted(ctrl\_s, seg\_s, side='right') - 1 の結果を代入する。
\item idx に np.clip(idx, 0, Nc - 1) の結果を代入する。
\item idx\_next に np.clip(idx + 1, 0, Nc - 1) の結果を代入する。
\item denom に np.maximum(ctrl\_s[idx\_next] - ctrl\_s[idx], 1e-06) の結果を代入する。
\item alpha に (seg\_s - ctrl\_s[idx]) / denom の結果を代入する。
\item 関数 expand\_ctrl を定義する。
\item 関数 build\_balance\_seed を定義する。
\item w\_dv\_limit に float(self.get\_parameter('w\_dv\_limit').value) の結果を代入する。
\item dv\_max\_kmhps に float(self.get\_parameter('dv\_max\_kmhps').value) の結果を代入する。
\item term\_soc\_min に float(self.get\_parameter('terminal\_soc\_min').value) の結果を代入する。
\item day\_end\_soc\_min に float(self.upper\_day\_end\_soc\_min) の結果を代入する。
\item soc\_day\_end\_max に float(self.soc\_day\_end\_max) の結果を代入する。
\item soc\_finish\_target に float(self.soc\_finish\_target) の結果を代入する。
\item soc\_day\_end\_tol に float(self.get\_parameter('soc\_day\_end\_tol').value) の結果を代入する。
\item 関数 integrate\_stationary\_duration を定義する。
\item 関数 step\_wait を定義する。
\item 関数 apply\_control\_stop\_at を定義する。
\item 関数 cost を定義する。
\item structured\_seeds に [('balance\_seed', build\_balance\_seed())] の結果を代入する。
\item global\_search\_enabled に bool(self.get\_parameter('upper\_global\_search\_enabled').value) の結果を代入する。
\item global\_search\_mode に str(self.get\_parameter('upper\_global\_search\_mode').value or 'auto').strip().lower() の結果を代入する。
\item (u\_seq, solve\_info) に hybrid\_bounded\_minimize(cost, x0, bounds, maxiter=int(self.upper\_max\_iter), structured\_seeds=structured\_seeds, cem\_enabled=global\_search\_enabled, cem\_mode=global\_search\_mode, cem\_generations=int(self.get\_parameter('upper\_cem\_generations').value), cem\_population=int(self.get\_parameter('upper\_cem\_population').value), cem\_elite=int(self.get\_parameter('upper\_cem\_elite').value), local\_refine\_topk=int(self.get\_parameter('upper\_local\_refine\_topk').value), seed\_library\_mode=str(self.get\_parameter('upper\_seed\_library\_mode').value), rng\_seed=int(max(0.0, round(s0\_km * 10.0))), shgo\_samples=int(self.get\_parameter('upper\_shgo\_samples').value), shgo\_iters=int(self.get\_parameter('upper\_shgo\_iters').value), cert\_grid\_levels=int(self.get\_parameter('upper\_cert\_grid\_levels').value), cert\_max\_evaluations=int(self.get\_parameter('upper\_cert\_max\_evaluations').value), cert\_workers=int(self.get\_parameter('upper\_cert\_workers').value)) の結果を代入する。
\item self.last\_upper\_solve\_ok に bool(solve\_info.get('success', False)) の結果を代入する。
\item 条件 not self.last\_upper\_solve\_ok を判定し、真なら内部処理を行う。
\item   self.get\_logger().warn(...) を実行する。
\item v\_seq に expand\_ctrl(u\_seq) の結果を代入する。
\item segments に [] の結果を代入する。
\item t\_utc に t0\_utc の結果を代入する。
\item s\_km に float(s0\_km) の結果を代入する。
\item z に float(self.z) の結果を代入する。
\item Tb に float(self.Tb) の結果を代入する。
\item v\_prev\_seg に float(v0) の結果を代入する。
\item enumerate(v\_seq) を順に走査し、各要素を (idx\_seg, v\_k) に入れて処理する。
\item   (t\_utc, z, Tb, schedule\_wait, \_) に step\_wait(t\_utc, z, Tb, s\_km) の結果を代入する。
\item   条件 schedule\_wait > 1e-09 を判定し、真なら内部処理を行う。
\item     v\_prev\_seg に 0.0 の結果を代入する。
\item   ds\_step\_km に float(ds\_seq[idx\_seg]) の結果を代入する。
\item   v\_k に float(np.clip(v\_k, 0.0, v\_max\_kmh)) の結果を代入する。
\item   vmax\_local に self.\_speed\_limit\_at(s\_km, v\_max\_kmh) の結果を代入する。
\item   条件 vmax\_local >= v\_min\_solver を判定し、真なら内部処理を行う。
\item     v\_k に max(v\_min\_solver, min(v\_k, vmax\_local)) の結果を代入する。
\item     上の条件が偽の場合:
            \end{enumerate}
            \paragraph{代表コード断片}
            \begin{lstlisting}[language=Python]
    def _mpc_solve_solar_distance(self, t0_utc: datetime, s0_km: float, x0=None):
        p = self.model.p
        horizon = build_upper_distance_horizon(
            mode=self.upper_horizon_mode,
            s0_km=s0_km,
            race_km=self.race_km,
            ds_km=self.upper_ds_km,
            horizon_km=self.upper_horizon_km,
            max_steps=self.upper_max_steps,
            ctrl_km=self.upper_ctrl_km,
            adaptive_min_ds_km=self.upper_adaptive_min_ds_km,
            adaptive_max_ds_km=self.upper_adaptive_max_ds_km,
            adaptive_growth=self.upper_adaptive_growth,
        )
        ds_seq = np.array(horizon.ds_seq_km, dtype=float)
        seg_s = np.array(horizon.seg_s_km, dtype=float)
        Np = int(len(ds_seq))
        if Np <= 0:
            return self.v_cmd, [{'v_kmh': float(self.v_cmd), 'dt_sec': float(p.dt)}], [float(self.v_cmd)]

        # Split the integration mesh at every declared control stop so arrival
        # time, dwell charging, and the next departure are evaluated exactly at
        # the stop coordinate rather than at the end of an arbitrary route bin.
        control_stops = self._control_stop_catalog()
        horizon_end_km = float(s0_km + np.sum(ds_seq))
        extra_edges = [
            float(stop['s_km'])
            for stop in control_stops
            if float(s0_km) + 1.0e-9 < float(stop['s_km']) < horizon_end_km - 1.0e-9
        ]
        if extra_edges:
            base_edges = np.concatenate(([float(s0_km)], float(s0_km) + np.cumsum(ds_seq)))
            route_edges = np.array(sorted(set(np.round(np.concatenate((base_edges, extra_edges)), 9))), dtype=float)
            ds_seq = np.diff(route_edges)
            seg_s = route_edges[:-1] - float(s0_km)
...
\end{lstlisting}

\subsection{L1379 関数 MPCNode.\_mpc\_solve\_solar\_distance.expand\_ctrl}
            \begin{itemize}[leftmargin=1.6em]
              \item 行範囲: L1379--L1380
              \item 所属: \path|MPCNode._mpc_solve_solar_distance|
              \item 定義: \path|expand_ctrl(u_vec)|
              \item docstring: 明示されていない。
              \item このブロックが直接呼ぶ主な関数・メソッド: 特に明示されていない。
              \item 戻り値の要点: (1.0 - alpha) * u\_vec[idx] + alpha * u\_vec[idx\_next]
              \item 主なローカル代入名: 特になし。
              \item 読み取る主なインスタンス属性: 特になし。
              \item 更新する主なインスタンス属性: 特になし。
              \item 明示的に送出する例外: 明示的raiseはない。
              \item 制御構造の規模: 条件分岐 0、ループ 0、try 0
            \end{itemize}
            \paragraph{この定義を読むためのPython構文}
            \begin{itemize}[leftmargin=1.6em]
            \item 特別な構文上の補足は少ない。
            \end{itemize}
            \paragraph{上から順の処理}
            \begin{enumerate}[leftmargin=1.8em]
            \item (1.0 - alpha) * u\_vec[idx] + alpha * u\_vec[idx\_next] を返す。
            \end{enumerate}
            \paragraph{代表コード断片}
            \begin{lstlisting}[language=Python]
        def expand_ctrl(u_vec):
            return (1.0 - alpha) * u_vec[idx] + alpha * u_vec[idx_next]
\end{lstlisting}

\subsection{L1382 関数 MPCNode.\_mpc\_solve\_solar\_distance.build\_balance\_seed}
            \begin{itemize}[leftmargin=1.6em]
              \item 行範囲: L1382--L1473
              \item 所属: \path|MPCNode._mpc_solve_solar_distance|
              \item 定義: \path|build_balance_seed()|
              \item docstring: 明示されていない。
              \item このブロックが直接呼ぶ主な関数・メソッド: \_forecast\_at\_time, \_route\_average, \_route\_value, \_speed\_limit\_at, abs, apply\_control\_stop\_at, array, clip, electrical\_balance, float, get, int
              \item 戻り値の要点: np.array(u\_seed, dtype=float)
              \item 主なローカル代入名: Tb\_seed, \_, best\_score, best\_v, candidate\_grid, control\_wait, ds\_ctrl, ds\_leg\_km, dt\_seed, edge\_s\_km, env, env\_leg, headwind\_leg, headwind\_ms, idx\_ctrl, limits, out, p\_pack, s\_seed, score
              \item 読み取る主なインスタンス属性: self.Tb, self.\_forecast\_at\_time, self.\_route\_average, self.\_route\_value, self.\_speed\_limit\_at, self.drive\_schedule, self.model, self.upper\_ds\_km, self.z
              \item 更新する主なインスタンス属性: 特になし。
              \item 明示的に送出する例外: 明示的raiseはない。
              \item 制御構造の規模: 条件分岐 6、ループ 3、try 0
            \end{itemize}
            \paragraph{この定義を読むためのPython構文}
            \begin{itemize}[leftmargin=1.6em]
            \item 内包表記は、反復・条件・値生成を一つの式で記述してcollectionを作る。
            \end{itemize}
            \paragraph{上から順の処理}
            \begin{enumerate}[leftmargin=1.8em]
            \item u\_seed に np.array(x0, dtype=float) の結果を代入する。
\item z\_seed に float(self.z) の結果を代入する。
\item Tb\_seed に float(self.Tb) の結果を代入する。
\item s\_seed に float(s0\_km) の結果を代入する。
\item t\_seed に t0\_utc の結果を代入する。
\item v\_prev\_seed に float(v0) の結果を代入する。
\item range(Nc) を順に走査し、各要素を idx\_ctrl に入れて処理する。
\item   (t\_seed, z\_seed, Tb\_seed, \_, \_) に step\_wait(t\_seed, z\_seed, Tb\_seed, s\_seed) の結果を代入する。
\item   vmax\_local に self.\_speed\_limit\_at(s\_seed, v\_max\_kmh) の結果を代入する。
\item   vmin\_local に v\_min\_solver の結果を代入する。
\item   条件 self.drive\_schedule is not None を判定し、真なら内部処理を行う。
\item     limits に self.drive\_schedule.speed\_limits(t\_seed) の結果を代入する。
\item     条件 limits is not None を判定し、真なら内部処理を行う。
\item       vmin\_local に max(vmin\_local, float(limits[0])) の結果を代入する。
\item       vmax\_local に min(vmax\_local, float(limits[1])) の結果を代入する。
\item   条件 vmax\_local < vmin\_local を判定し、真なら内部処理を行う。
\item     u\_seed[idx\_ctrl] に max(0.0, vmax\_local) の結果を代入する。
\item     Continue 文を実行する。
\item   candidate\_grid に np.linspace(vmin\_local, vmax\_local, num=max(5, min(13, int((vmax\_local - vmin\_local) / 4.0) + 2))) の結果を代入する。
\item   best\_v に float(np.clip(v\_prev\_seed, vmin\_local, vmax\_local)) の結果を代入する。
\item   best\_score に float('inf') の結果を代入する。
\item   env に self.\_forecast\_at\_time(t\_seed, s\_seed, drive=True) の結果を代入する。
\item   slope\_pct に self.\_route\_value(s\_seed, 'slope\_pct', 0.0) の結果を代入する。
\item   headwind\_ms に self.\_route\_value(s\_seed, 'headwind\_ms', env.get('headwind\_ms', 0.0)) の結果を代入する。
\item   candidate\_grid を順に走査し、各要素を v\_test に入れて処理する。
\item     out に self.model.electrical\_balance(v\_test / 3.6, slope\_pct, z\_seed, Tb\_seed, env['G\_poa'], env['Tcell\_C'], headwind\_ms=headwind\_ms, ambient\_temp\_c=env.get('Tamb\_C'), elevation\_m=env.get('elevation\_m', self.\_route\_value(s\_seed, 'elev\_m', 0.0))) の結果を代入する。
\item     p\_pack に float(out['P\_pack']) の結果を代入する。
\item     score に abs(p\_pack) + 0.35 * (float(v\_test) - float(v\_prev\_seed)) ** 2 の結果を代入する。
\item     条件 score < best\_score を判定し、真なら内部処理を行う。
\item       best\_score に score の結果を代入する。
\item       best\_v に float(v\_test) の結果を代入する。
\item   u\_seed[idx\_ctrl] に best\_v の結果を代入する。
\item   条件 idx\_ctrl + 1 < len(ctrl\_s) を判定し、真なら内部処理を行う。
\item     ds\_ctrl に max(0.0, float(ctrl\_s[idx\_ctrl + 1] - ctrl\_s[idx\_ctrl])) の結果を代入する。
\item     上の条件が偽の場合:
\item     ds\_ctrl に max(float(ds\_seq[-1]), float(self.upper\_ds\_km)) の結果を代入する。
\item   target\_s\_km に s\_seed + ds\_ctrl の結果を代入する。
\item   stop\_edges に [float(stop['s\_km']) for stop in control\_stops if s\_seed + 1e-09 < float(stop['s\_km']) <= target\_s\_km + 1e-09] の結果を代入する。
\item   seed\_edges に sorted(set([*stop\_edges, float(target\_s\_km)])) の結果を代入する。
\item   stopped\_in\_chunk に False の結果を代入する。
\item   seed\_edges を順に走査し、各要素を edge\_s\_km に入れて処理する。
\item     ds\_leg\_km に max(0.0, float(edge\_s\_km) - s\_seed) の結果を代入する。
\item     条件 ds\_leg\_km > 1e-09 を判定し、真なら内部処理を行う。
\item       dt\_seed に ds\_leg\_km / max(best\_v, 0.001) * 3600.0 の結果を代入する。
\item       env\_leg に self.\_forecast\_at\_time(t\_seed, s\_seed, drive=True) の結果を代入する。
\item       slope\_leg に self.\_route\_average(s\_seed, edge\_s\_km, 'slope\_pct', 0.0) の結果を代入する。
\item       headwind\_leg に self.\_route\_value(s\_seed, 'headwind\_ms', env\_leg.get('headwind\_ms', 0.0)) の結果を代入する。
\item       out に self.model.electrical\_balance(best\_v / 3.6, slope\_leg, z\_seed, Tb\_seed, env\_leg['G\_poa'], env\_leg['Tcell\_C'], headwind\_ms=headwind\_leg, ambient\_temp\_c=env\_leg.get('Tamb\_C'), elevation\_m=env\_leg.get('elevation\_m', self.\_route\_value(s\_seed, 'elev\_m', 0.0))) の結果を代入する。
\item       z\_seed に float(np.clip(self.model.soc\_step(z\_seed, float(out['P\_pack']), dt\_seed, current\_a=float(out['I']), Tbat\_C=Tb\_seed), p.soc\_min, p.soc\_max)) の結果を代入する。
\item       Tb\_seed に float(np.clip(Tb\_seed + dt\_seed / 1800.0 * (env\_leg['Tamb\_C'] - Tb\_seed) + float(out['losses\_int']) * dt\_seed / 50000.0, p.T\_min, p.T\_max)) の結果を代入する。
\item       t\_seed を Add で更新する。
\item       s\_seed に float(edge\_s\_km) の結果を代入する。
\item     (t\_seed, z\_seed, Tb\_seed, control\_wait, \_) に apply\_control\_stop\_at(t\_seed, z\_seed, Tb\_seed, s\_seed) の結果を代入する。
\item     stopped\_in\_chunk に stopped\_in\_chunk or control\_wait > 1e-09 の結果を代入する。
\item   v\_prev\_seed に 0.0 if stopped\_in\_chunk else best\_v の結果を代入する。
\item np.array(u\_seed, dtype=float) を返す。
            \end{enumerate}
            \paragraph{代表コード断片}
            \begin{lstlisting}[language=Python]
        def build_balance_seed():
            u_seed = np.array(x0, dtype=float)
            z_seed = float(self.z)
            Tb_seed = float(self.Tb)
            s_seed = float(s0_km)
            t_seed = t0_utc
            v_prev_seed = float(v0)
            for idx_ctrl in range(Nc):
                t_seed, z_seed, Tb_seed, _, _ = step_wait(t_seed, z_seed, Tb_seed, s_seed)
                vmax_local = self._speed_limit_at(s_seed, v_max_kmh)
                vmin_local = v_min_solver
                if self.drive_schedule is not None:
                    limits = self.drive_schedule.speed_limits(t_seed)
                    if limits is not None:
                        vmin_local = max(vmin_local, float(limits[0]))
                        vmax_local = min(vmax_local, float(limits[1]))
                if vmax_local < vmin_local:
                    u_seed[idx_ctrl] = max(0.0, vmax_local)
                    continue
                candidate_grid = np.linspace(vmin_local, vmax_local, num=max(5, min(13, int((vmax_local - vmin_local) / 4.0) + 2)))
                best_v = float(np.clip(v_prev_seed, vmin_local, vmax_local))
                best_score = float("inf")
                env = self._forecast_at_time(t_seed, s_seed, drive=True)
                slope_pct = self._route_value(s_seed, 'slope_pct', 0.0)
                headwind_ms = self._route_value(s_seed, 'headwind_ms', env.get('headwind_ms', 0.0))
                for v_test in candidate_grid:
                    out = self.model.electrical_balance(
                        v_test / 3.6, slope_pct, z_seed, Tb_seed, env['G_poa'], env['Tcell_C'],
                        headwind_ms=headwind_ms, ambient_temp_c=env.get('Tamb_C'),
                        elevation_m=env.get('elevation_m', self._route_value(s_seed, 'elev_m', 0.0)),
                    )
                    p_pack = float(out['P_pack'])
                    score = abs(p_pack) + 0.35 * ((float(v_test) - float(v_prev_seed)) ** 2)
                    if score < best_score:
                        best_score = score
...
\end{lstlisting}

\subsection{L1483 関数 MPCNode.\_mpc\_solve\_solar\_distance.integrate\_stationary\_duration}
            \begin{itemize}[leftmargin=1.6em]
              \item 行範囲: L1483--L1536
              \item 所属: \path|MPCNode._mpc_solve_solar_distance|
              \item 定義: \path|integrate_stationary_duration(t_utc, z, Tb, s_km, dt_wait, *, control_stop)|
              \item docstring: 明示されていない。
              \item このブロックが直接呼ぶ主な関数・メソッド: \_forecast\_at\_time, \_route\_value, clip, electrical\_balance, float, get, max, min, quad\_penalty, scheduled\_auxiliary\_power, soc\_step, timedelta
              \item 戻り値の要点: (t\_cursor, float(z), float(Tb), dt\_wait, float(wait\_cost)) / (t\_utc, float(z), float(Tb), 0.0, 0.0)
              \item 主なローカル代入名: I, P\_pack, Tb, V, constraint, dt\_local, dt\_sample, dt\_wait, env, headwind\_ms, loss\_int, out, remaining, slope\_pct, t\_cursor, wait\_cost, z
              \item 読み取る主なインスタンス属性: self.\_forecast\_at\_time, self.\_route\_value, self.model, self.upper\_cost\_cfg
              \item 更新する主なインスタンス属性: 特になし。
              \item 明示的に送出する例外: 明示的raiseはない。
              \item 制御構造の規模: 条件分岐 1、ループ 1、try 0
            \end{itemize}
            \paragraph{この定義を読むためのPython構文}
            \begin{itemize}[leftmargin=1.6em]
            \item 特別な構文上の補足は少ない。
            \end{itemize}
            \paragraph{上から順の処理}
            \begin{enumerate}[leftmargin=1.8em]
            \item dt\_wait に max(0.0, float(dt\_wait)) の結果を代入する。
\item 条件 dt\_wait <= 0.0 を判定し、真なら内部処理を行う。
\item   (t\_utc, float(z), float(Tb), 0.0, 0.0) を返す。
\item slope\_pct に self.\_route\_value(s\_km, 'slope\_pct', 0.0) の結果を代入する。
\item t\_cursor に t\_utc の結果を代入する。
\item remaining に dt\_wait の結果を代入する。
\item dt\_sample に min(600.0, max(60.0, float(p.dt))) の結果を代入する。
\item wait\_cost に float(self.upper\_cost\_cfg.w\_wait) * dt\_wait の結果を代入する。
\item 条件 remaining > 1e-09 が成り立つ間くり返す。
\item   dt\_local に min(dt\_sample, remaining) の結果を代入する。
\item   env に self.\_forecast\_at\_time(t\_cursor, s\_km, drive=False, control\_stop=control\_stop) の結果を代入する。
\item   headwind\_ms に self.\_route\_value(s\_km, 'headwind\_ms', env.get('headwind\_ms', 0.0)) の結果を代入する。
\item   out に self.model.electrical\_balance(0.0, slope\_pct, z, Tb, env['G\_poa'], env['Tcell\_C'], headwind\_ms=headwind\_ms, aux\_power\_w=self.model.scheduled\_auxiliary\_power(is\_driving=False, irradiance\_wm2=env['G\_poa']), ambient\_temp\_c=env.get('Tamb\_C'), elevation\_m=env.get('elevation\_m', self.\_route\_value(s\_km, 'elev\_m', 0.0))) の結果を代入する。
\item   P\_pack に float(out['P\_pack']) の結果を代入する。
\item   loss\_int に float(out['losses\_int']) の結果を代入する。
\item   I に float(out['I']) の結果を代入する。
\item   V に float(out['V']) の結果を代入する。
\item   z に float(np.clip(self.model.soc\_step(z, P\_pack, dt\_local, current\_a=I, Tbat\_C=Tb), p.soc\_min, p.soc\_max)) の結果を代入する。
\item   Tb に float(np.clip(Tb + dt\_local / 1800.0 * (env['Tamb\_C'] - Tb) + loss\_int * dt\_local / 50000.0, p.T\_min, p.T\_max)) の結果を代入する。
\item   constraint に float(self.upper\_cost\_cfg.constraint\_penalty) の結果を代入する。
\item   wait\_cost を Add で更新する。
\item   t\_cursor を Add で更新する。
\item   remaining を Sub で更新する。
\item (t\_cursor, float(z), float(Tb), dt\_wait, float(wait\_cost)) を返す。
            \end{enumerate}
            \paragraph{代表コード断片}
            \begin{lstlisting}[language=Python]
        def integrate_stationary_duration(t_utc, z, Tb, s_km, dt_wait, *, control_stop):
            dt_wait = max(0.0, float(dt_wait))
            if dt_wait <= 0.0:
                return t_utc, float(z), float(Tb), 0.0, 0.0
            slope_pct = self._route_value(s_km, 'slope_pct', 0.0)
            t_cursor = t_utc
            remaining = dt_wait
            dt_sample = min(600.0, max(60.0, float(p.dt)))
            wait_cost = float(self.upper_cost_cfg.w_wait) * dt_wait
            while remaining > 1.0e-9:
                dt_local = min(dt_sample, remaining)
                env = self._forecast_at_time(
                    t_cursor, s_km, drive=False, control_stop=control_stop
                )
                headwind_ms = self._route_value(s_km, 'headwind_ms', env.get('headwind_ms', 0.0))
                out = self.model.electrical_balance(
                    0.0,
                    slope_pct,
                    z,
                    Tb,
                    env['G_poa'],
                    env['Tcell_C'],
                    headwind_ms=headwind_ms,
                    aux_power_w=self.model.scheduled_auxiliary_power(
                        is_driving=False,
                        irradiance_wm2=env['G_poa'],
                    ),
                    ambient_temp_c=env.get('Tamb_C'),
                    elevation_m=env.get('elevation_m', self._route_value(s_km, 'elev_m', 0.0)),
                )
                P_pack = float(out['P_pack'])
                loss_int = float(out['losses_int'])
                I = float(out['I'])
                V = float(out['V'])
                z = float(np.clip(self.model.soc_step(
...
\end{lstlisting}

\subsection{L1538 関数 MPCNode.\_mpc\_solve\_solar\_distance.step\_wait}
            \begin{itemize}[leftmargin=1.6em]
              \item 行範囲: L1538--L1545
              \item 所属: \path|MPCNode._mpc_solve_solar_distance|
              \item 定義: \path|step_wait(t_utc, z, Tb, s_km)|
              \item docstring: 明示されていない。
              \item このブロックが直接呼ぶ主な関数・メソッド: float, integrate\_stationary\_duration, is\_drive\_time, max, next\_drive\_start, total\_seconds
              \item 戻り値の要点: integrate\_stationary\_duration(t\_utc, z, Tb, s\_km, dt\_wait, control\_stop=False) / (t\_utc, float(z), float(Tb), 0.0, 0.0)
              \item 主なローカル代入名: dt\_wait, t\_start
              \item 読み取る主なインスタンス属性: self.drive\_schedule
              \item 更新する主なインスタンス属性: 特になし。
              \item 明示的に送出する例外: 明示的raiseはない。
              \item 制御構造の規模: 条件分岐 1、ループ 0、try 0
            \end{itemize}
            \paragraph{この定義を読むためのPython構文}
            \begin{itemize}[leftmargin=1.6em]
            \item 特別な構文上の補足は少ない。
            \end{itemize}
            \paragraph{上から順の処理}
            \begin{enumerate}[leftmargin=1.8em]
            \item 条件 self.drive\_schedule is None or self.drive\_schedule.is\_drive\_time(t\_utc) を判定し、真なら内部処理を行う。
\item   (t\_utc, float(z), float(Tb), 0.0, 0.0) を返す。
\item t\_start に self.drive\_schedule.next\_drive\_start(t\_utc) の結果を代入する。
\item dt\_wait に max(0.0, (t\_start - t\_utc).total\_seconds()) の結果を代入する。
\item integrate\_stationary\_duration(t\_utc, z, Tb, s\_km, dt\_wait, control\_stop=False) を返す。
            \end{enumerate}
            \paragraph{代表コード断片}
            \begin{lstlisting}[language=Python]
        def step_wait(t_utc, z, Tb, s_km):
            if self.drive_schedule is None or self.drive_schedule.is_drive_time(t_utc):
                return t_utc, float(z), float(Tb), 0.0, 0.0
            t_start = self.drive_schedule.next_drive_start(t_utc)
            dt_wait = max(0.0, (t_start - t_utc).total_seconds())
            return integrate_stationary_duration(
                t_utc, z, Tb, s_km, dt_wait, control_stop=False
            )
\end{lstlisting}

\subsection{L1547 関数 MPCNode.\_mpc\_solve\_solar\_distance.apply\_control\_stop\_at}
            \begin{itemize}[leftmargin=1.6em]
              \item 行範囲: L1547--L1563
              \item 所属: \path|MPCNode._mpc_solve_solar_distance|
              \item 定義: \path|apply_control_stop_at(t_utc, z, Tb, s_km)|
              \item docstring: 明示されていない。
              \item このブロックが直接呼ぶ主な関数・メソッド: abs, float, integrate\_stationary\_duration
              \item 戻り値の要点: (t\_utc, float(z), float(Tb), total\_wait, total\_cost)
              \item 主なローカル代入名: Tb, dt\_stop, stop, stop\_cost, t\_utc, total\_cost, total\_wait, z
              \item 読み取る主なインスタンス属性: 特になし。
              \item 更新する主なインスタンス属性: 特になし。
              \item 明示的に送出する例外: 明示的raiseはない。
              \item 制御構造の規模: 条件分岐 1、ループ 1、try 0
            \end{itemize}
            \paragraph{この定義を読むためのPython構文}
            \begin{itemize}[leftmargin=1.6em]
            \item 特別な構文上の補足は少ない。
            \end{itemize}
            \paragraph{上から順の処理}
            \begin{enumerate}[leftmargin=1.8em]
            \item total\_wait に 0.0 の結果を代入する。
\item total\_cost に 0.0 の結果を代入する。
\item control\_stops を順に走査し、各要素を stop に入れて処理する。
\item   条件 abs(float(stop['s\_km']) - float(s\_km)) > 1e-07 を判定し、真なら内部処理を行う。
\item     Continue 文を実行する。
\item   (t\_utc, z, Tb, dt\_stop, stop\_cost) に integrate\_stationary\_duration(t\_utc, z, Tb, float(stop['s\_km']), float(stop['dwell\_sec']), control\_stop=True) の結果を代入する。
\item   total\_wait を Add で更新する。
\item   total\_cost を Add で更新する。
\item (t\_utc, float(z), float(Tb), total\_wait, total\_cost) を返す。
            \end{enumerate}
            \paragraph{代表コード断片}
            \begin{lstlisting}[language=Python]
        def apply_control_stop_at(t_utc, z, Tb, s_km):
            total_wait = 0.0
            total_cost = 0.0
            for stop in control_stops:
                if abs(float(stop['s_km']) - float(s_km)) > 1.0e-7:
                    continue
                t_utc, z, Tb, dt_stop, stop_cost = integrate_stationary_duration(
                    t_utc,
                    z,
                    Tb,
                    float(stop['s_km']),
                    float(stop['dwell_sec']),
                    control_stop=True,
                )
                total_wait += dt_stop
                total_cost += stop_cost
            return t_utc, float(z), float(Tb), total_wait, total_cost
\end{lstlisting}

\subsection{L1565 関数 MPCNode.\_mpc\_solve\_solar\_distance.cost}
            \begin{itemize}[leftmargin=1.6em]
              \item 行範囲: L1565--L1714
              \item 所属: \path|MPCNode._mpc_solve_solar_distance|
              \item 定義: \path|cost(u_vec)|
              \item docstring: 明示されていない。
              \item このブロックが直接呼ぶ主な関数・メソッド: \_forecast\_at\_time, \_route\_average, \_route\_value, \_speed\_limit\_at, apply\_control\_stop\_at, bool, current\_drive\_window, electrical\_balance, expand\_ctrl, float, get, is\_drive\_time
              \item 戻り値の要点: J
              \item 主なローカル代入名: I, J, P\_mech\_wheel, P\_pack, P\_pv, Tb, Tb\_after, Tb\_next, V, day\_end\_crossing, ds\_step\_km, dt\_travel, dt\_wait, elapsed\_plan\_sec, env, forces, headwind\_ms, inertial\_power\_w, k, kinetic\_delta\_wh
              \item 読み取る主なインスタンス属性: self.Tb, self.\_forecast\_at\_time, self.\_route\_average, self.\_route\_value, self.\_speed\_limit\_at, self.drive\_schedule, self.model, self.soc\_day\_end\_target, self.soc\_finish\_tol, self.upper\_cost\_cfg, self.z
              \item 更新する主なインスタンス属性: 特になし。
              \item 明示的に送出する例外: 明示的raiseはない。
              \item 制御構造の規模: 条件分岐 7、ループ 1、try 0
            \end{itemize}
            \paragraph{この定義を読むためのPython構文}
            \begin{itemize}[leftmargin=1.6em]
            \item 特別な構文上の補足は少ない。
            \end{itemize}
            \paragraph{上から順の処理}
            \begin{enumerate}[leftmargin=1.8em]
            \item z に float(self.z) の結果を代入する。
\item Tb に float(self.Tb) の結果を代入する。
\item s\_km に float(s0\_km) の結果を代入する。
\item t\_utc に t0\_utc の結果を代入する。
\item v\_prev に v0 の結果を代入する。
\item p\_pack\_prev に None の結果を代入する。
\item elapsed\_plan\_sec に 0.0 の結果を代入する。
\item J に 0.0 の結果を代入する。
\item v\_seq に expand\_ctrl(u\_vec) の結果を代入する。
\item range(Np) を順に走査し、各要素を k に入れて処理する。
\item   (t\_utc, z, Tb, dt\_wait, wait\_cost) に step\_wait(t\_utc, z, Tb, s\_km) の結果を代入する。
\item   J を Add で更新する。
\item   条件 dt\_wait > 1e-09 を判定し、真なら内部処理を行う。
\item     v\_prev に 0.0 の結果を代入する。
\item   v\_k に float(v\_seq[k]) の結果を代入する。
\item   ds\_step\_km に float(ds\_seq[k]) の結果を代入する。
\item   vmax\_local に self.\_speed\_limit\_at(s\_km, v\_max\_kmh) の結果を代入する。
\item   条件 vmax\_local >= v\_min\_solver を判定し、真なら内部処理を行う。
\item     v\_k に max(v\_min\_solver, min(v\_k, vmax\_local)) の結果を代入する。
\item     上の条件が偽の場合:
\item     v\_k に max(0.0, min(v\_k, vmax\_local)) の結果を代入する。
\item   limits に None の結果を代入する。
\item   条件 self.drive\_schedule is not None を判定し、真なら内部処理を行う。
\item     limits に self.drive\_schedule.speed\_limits(t\_utc) の結果を代入する。
\item   dt\_travel に ds\_step\_km / max(v\_k, 0.001) * 3600.0 の結果を代入する。
\item   env に self.\_forecast\_at\_time(t\_utc, s\_km, drive=True) の結果を代入する。
\item   slope\_pct に self.\_route\_average(s\_km, s\_km + ds\_step\_km, 'slope\_pct', 0.0) の結果を代入する。
\item   headwind\_ms に self.\_route\_value(s\_km, 'headwind\_ms', env.get('headwind\_ms', 0.0)) の結果を代入する。
\item   kinetic\_delta\_wh に 0.5 * p.m * ((v\_k / 3.6) ** 2 - (v\_prev / 3.6) ** 2) / 3600.0 の結果を代入する。
\item   inertial\_power\_w に kinetic\_delta\_wh * 3600.0 / max(dt\_travel, 1e-09) の結果を代入する。
\item   out に self.model.electrical\_balance(v\_k / 3.6, slope\_pct, z, Tb, env['G\_poa'], env['Tcell\_C'], headwind\_ms=headwind\_ms, inertial\_power\_w=inertial\_power\_w, ambient\_temp\_c=env.get('Tamb\_C'), elevation\_m=env.get('elevation\_m', self.\_route\_value(s\_km, 'elev\_m', 0.0))) の結果を代入する。
\item   I に float(out['I']) の結果を代入する。
\item   V に float(out['V']) の結果を代入する。
\item   P\_pv に float(out.get('P\_pv', 0.0)) の結果を代入する。
\item   P\_pack に float(out['P\_pack']) の結果を代入する。
\item   loss\_int に float(out['losses\_int']) の結果を代入する。
\item   loss\_line に float(out.get('losses\_line', 0.0)) の結果を代入する。
\item   P\_mech\_wheel に float(out.get('P\_mech\_wheel', 0.0)) の結果を代入する。
\item   kinetic\_step\_wh に max(0.0, kinetic\_delta\_wh) の結果を代入する。
\item   z\_next に self.model.soc\_step(z, P\_pack, dt\_travel, current\_a=I, Tbat\_C=Tb) の結果を代入する。
\item   Tb\_next に Tb + dt\_travel / 1800.0 * (env['Tamb\_C'] - Tb) + loss\_int * dt\_travel / 50000.0 の結果を代入する。
\item   forces に self.model.resistive\_forces(v\_k / 3.6, slope\_pct, headwind\_ms=headwind\_ms, ambient\_temp\_c=env.get('Tamb\_C'), elevation\_m=env.get('elevation\_m', self.\_route\_value(s\_km, 'elev\_m', 0.0))) の結果を代入する。
\item   soc\_line に None の結果を代入する。
\item   条件 self.drive\_schedule is not None and self.soc\_day\_end\_target > 0.0 を判定し、真なら内部処理を行う。
\item     win に self.drive\_schedule.current\_drive\_window(t\_utc) の結果を代入する。
\item     条件 win is not None を判定し、真なら内部処理を行う。
\item       (t\_start, t\_end) に win の結果を代入する。
\item       条件 t\_end > t\_start を判定し、真なら内部処理を行う。
\item   t\_next に t\_utc + timedelta(seconds=dt\_travel) の結果を代入する。
\item   day\_end\_crossing に bool(self.drive\_schedule is not None and self.drive\_schedule.is\_drive\_time(t\_utc) and (not self.drive\_schedule.is\_drive\_time(t\_next))) の結果を代入する。
\item   elapsed\_plan\_sec を Add で更新する。
\item   J を Add で更新する。
\item   s\_next\_km に s\_km + ds\_step\_km の結果を代入する。
\item   (t\_after, z\_after, Tb\_after, stop\_wait, stop\_cost) に apply\_control\_stop\_at(t\_next, z\_next, Tb\_next, s\_next\_km) の結果を代入する。
\item   J を Add で更新する。
\item   elapsed\_plan\_sec を Add で更新する。
\item   t\_utc に t\_after の結果を代入する。
\item   s\_km に s\_next\_km の結果を代入する。
\item   (z, Tb) に (z\_after, Tb\_after) の結果を代入する。
\item   条件 stop\_wait > 1e-09 を判定し、真なら内部処理を行う。
\item     v\_prev に 0.0 の結果を代入する。
\item     p\_pack\_prev に None の結果を代入する。
\item     上の条件が偽の場合:
\item     v\_prev に v\_k の結果を代入する。
\item     p\_pack\_prev に P\_pack の結果を代入する。
\item J を Add で更新する。
\item J を返す。
            \end{enumerate}
            \paragraph{代表コード断片}
            \begin{lstlisting}[language=Python]
        def cost(u_vec):
            z = float(self.z)
            Tb = float(self.Tb)
            s_km = float(s0_km)
            t_utc = t0_utc
            v_prev = v0
            p_pack_prev = None
            elapsed_plan_sec = 0.0
            J = 0.0
            v_seq = expand_ctrl(u_vec)
            for k in range(Np):
                t_utc, z, Tb, dt_wait, wait_cost = step_wait(t_utc, z, Tb, s_km)
                J += wait_cost
                if dt_wait > 1.0e-9:
                    v_prev = 0.0
                v_k = float(v_seq[k])
                ds_step_km = float(ds_seq[k])
                vmax_local = self._speed_limit_at(s_km, v_max_kmh)
                if vmax_local >= v_min_solver:
                    v_k = max(v_min_solver, min(v_k, vmax_local))
                else:
                    v_k = max(0.0, min(v_k, vmax_local))
                limits = None
                if self.drive_schedule is not None:
                    limits = self.drive_schedule.speed_limits(t_utc)

                dt_travel = ds_step_km / max(v_k, 1.0e-3) * 3600.0
                env = self._forecast_at_time(t_utc, s_km, drive=True)
                slope_pct = self._route_average(
                    s_km, s_km + ds_step_km, 'slope_pct', 0.0
                )
                headwind_ms = self._route_value(s_km, 'headwind_ms', env.get('headwind_ms', 0.0))
                kinetic_delta_wh = 0.5 * p.m * (
                    (v_k / 3.6) ** 2 - (v_prev / 3.6) ** 2
                ) / 3600.0
...
\end{lstlisting}

\subsection{L1816 関数 MPCNode.\_publish\_upper\_plan}
            \begin{itemize}[leftmargin=1.6em]
              \item 行範囲: L1816--L1821
              \item 所属: \path|MPCNode|
              \item 定義: \path|_publish_upper_plan(self)|
              \item docstring: 明示されていない。
              \item このブロックが直接呼ぶ主な関数・メソッド: Float32MultiArray, float, publish
              \item 戻り値の要点: 戻り値が無いか、return 文が明示されていない。
              \item 主なローカル代入名: msg, v
              \item 読み取る主なインスタンス属性: self.plan\_dt\_sec, self.pub\_plan, self.v\_plan\_kmh
              \item 更新する主なインスタンス属性: 特になし。
              \item 明示的に送出する例外: 明示的raiseはない。
              \item 制御構造の規模: 条件分岐 1、ループ 0、try 0
            \end{itemize}
            \paragraph{この定義を読むためのPython構文}
            \begin{itemize}[leftmargin=1.6em]
            \item 第1引数selfは、このメソッドを呼んだインスタンス自身を受け取る慣習名である。
\item 内包表記は、反復・条件・値生成を一つの式で記述してcollectionを作る。
            \end{itemize}
            \paragraph{上から順の処理}
            \begin{enumerate}[leftmargin=1.8em]
            \item 条件 not self.v\_plan\_kmh を判定し、真なら内部処理を行う。
\item    を返す。
\item msg に Float32MultiArray() の結果を代入する。
\item msg.data に [float(self.plan\_dt\_sec)] + [float(v) for v in self.v\_plan\_kmh] の結果を代入する。
\item self.pub\_plan.publish(...) を実行する。
            \end{enumerate}
            \paragraph{代表コード断片}
            \begin{lstlisting}[language=Python]
    def _publish_upper_plan(self):
        if not self.v_plan_kmh:
            return
        msg = Float32MultiArray()
        msg.data = [float(self.plan_dt_sec)] + [float(v) for v in self.v_plan_kmh]
        self.pub_plan.publish(msg)
\end{lstlisting}

\subsection{L1823 関数 MPCNode.\_publish\_lower\_plan}
            \begin{itemize}[leftmargin=1.6em]
              \item 行範囲: L1823--L1828
              \item 所属: \path|MPCNode|
              \item 定義: \path|_publish_lower_plan(self, v_seq_ms)|
              \item docstring: 明示されていない。
              \item このブロックが直接呼ぶ主な関数・メソッド: Float32MultiArray, float, publish
              \item 戻り値の要点: 戻り値が無いか、return 文が明示されていない。
              \item 主なローカル代入名: msg, v
              \item 読み取る主なインスタンス属性: self.lower\_dt, self.pub\_lower\_plan
              \item 更新する主なインスタンス属性: 特になし。
              \item 明示的に送出する例外: 明示的raiseはない。
              \item 制御構造の規模: 条件分岐 1、ループ 0、try 0
            \end{itemize}
            \paragraph{この定義を読むためのPython構文}
            \begin{itemize}[leftmargin=1.6em]
            \item 第1引数selfは、このメソッドを呼んだインスタンス自身を受け取る慣習名である。
\item 内包表記は、反復・条件・値生成を一つの式で記述してcollectionを作る。
            \end{itemize}
            \paragraph{上から順の処理}
            \begin{enumerate}[leftmargin=1.8em]
            \item 条件 not v\_seq\_ms を判定し、真なら内部処理を行う。
\item    を返す。
\item msg に Float32MultiArray() の結果を代入する。
\item msg.data に [float(self.lower\_dt)] + [float(v * 3.6) for v in v\_seq\_ms] の結果を代入する。
\item self.pub\_lower\_plan.publish(...) を実行する。
            \end{enumerate}
            \paragraph{代表コード断片}
            \begin{lstlisting}[language=Python]
    def _publish_lower_plan(self, v_seq_ms):
        if not v_seq_ms:
            return
        msg = Float32MultiArray()
        msg.data = [float(self.lower_dt)] + [float(v * 3.6) for v in v_seq_ms]
        self.pub_lower_plan.publish(msg)
\end{lstlisting}

\subsection{L1830 関数 MPCNode.\_publish\_plan\_path}
            \begin{itemize}[leftmargin=1.6em]
              \item 行範囲: L1830--L1856
              \item 所属: \path|MPCNode|
              \item 定義: \path|_publish_plan_path(self, data)|
              \item docstring: 明示されていない。
              \item このブロックが直接呼ぶ主な関数・メソッド: Path, PoseStamped, append, bool, float, get\_clock, get\_parameter, isfinite, len, now, publish, to\_msg
              \item 戻り値の要点: 戻り値が無いか、return 文が明示されていない。
              \item 主なローカル代入名: path, pose, s\_tmp, seg, v\_kmh, v\_list
              \item 読み取る主なインスタンス属性: self.get\_clock, self.get\_parameter, self.model, self.pub\_path, self.s\_km, self.s\_meas, self.upper\_plan\_mode, self.v\_plan\_kmh, self.v\_plan\_segments, self.v\_upper\_cmd
              \item 更新する主なインスタンス属性: 特になし。
              \item 明示的に送出する例外: 明示的raiseはない。
              \item 制御構造の規模: 条件分岐 3、ループ 2、try 0
            \end{itemize}
            \paragraph{この定義を読むためのPython構文}
            \begin{itemize}[leftmargin=1.6em]
            \item 第1引数selfは、このメソッドを呼んだインスタンス自身を受け取る慣習名である。
            \end{itemize}
            \paragraph{上から順の処理}
            \begin{enumerate}[leftmargin=1.8em]
            \item 条件 len(data) == 0 を判定し、真なら内部処理を行う。
\item    を返す。
\item path に Path() の結果を代入する。
\item path.header.stamp に self.get\_clock().now().to\_msg() の結果を代入する。
\item path.header.frame\_id に 'map' の結果を代入する。
\item s\_tmp に float(self.s\_km) の結果を代入する。
\item 条件 bool(self.get\_parameter('use\_measured\_s').value) and np.isfinite(self.s\_meas) を判定し、真なら内部処理を行う。
\item   s\_tmp に float(self.s\_meas) の結果を代入する。
\item 条件 self.upper\_plan\_mode == 'distance' and self.v\_plan\_segments を判定し、真なら内部処理を行う。
\item   self.v\_plan\_segments を順に走査し、各要素を seg に入れて処理する。
\item     s\_tmp を Add で更新する。
\item     pose に PoseStamped() の結果を代入する。
\item     pose.header に path.header の結果を代入する。
\item     pose.pose.position.x に s\_tmp の結果を代入する。
\item     pose.pose.position.y に 0.0 の結果を代入する。
\item     path.poses.append(...) を実行する。
\item   上の条件が偽の場合:
\item   v\_list に self.v\_plan\_kmh if self.v\_plan\_kmh else [float(self.v\_upper\_cmd)] * len(data) の結果を代入する。
\item   v\_list[:len(data)] を順に走査し、各要素を v\_kmh に入れて処理する。
\item     s\_tmp を Add で更新する。
\item     pose に PoseStamped() の結果を代入する。
\item     pose.header に path.header の結果を代入する。
\item     pose.pose.position.x に s\_tmp の結果を代入する。
\item     pose.pose.position.y に 0.0 の結果を代入する。
\item     path.poses.append(...) を実行する。
\item self.pub\_path.publish(...) を実行する。
            \end{enumerate}
            \paragraph{代表コード断片}
            \begin{lstlisting}[language=Python]
    def _publish_plan_path(self, data):
        if len(data) == 0:
            return
        path = Path()
        path.header.stamp = self.get_clock().now().to_msg()
        path.header.frame_id = 'map'
        s_tmp = float(self.s_km)
        if bool(self.get_parameter('use_measured_s').value) and np.isfinite(self.s_meas):
            s_tmp = float(self.s_meas)
        if self.upper_plan_mode == 'distance' and self.v_plan_segments:
            for seg in self.v_plan_segments:
                s_tmp += float(seg['v_kmh']) * (seg['dt_sec'] / 3600.0)
                pose = PoseStamped()
                pose.header = path.header
                pose.pose.position.x = s_tmp
                pose.pose.position.y = 0.0
                path.poses.append(pose)
        else:
            v_list = self.v_plan_kmh if self.v_plan_kmh else [float(self.v_upper_cmd)] * len(data)
            for v_kmh in v_list[:len(data)]:
                s_tmp += float(v_kmh) * (self.model.p.dt / 3600.0)
                pose = PoseStamped()
                pose.header = path.header
                pose.pose.position.x = s_tmp
                pose.pose.position.y = 0.0
                path.poses.append(pose)
        self.pub_path.publish(path)
\end{lstlisting}

\subsection{L1858 関数 MPCNode.\_publish\_metrics}
            \begin{itemize}[leftmargin=1.6em]
              \item 行範囲: L1858--L1895
              \item 所属: \path|MPCNode|
              \item 定義: \path|_publish_metrics(self, d0: dict, v_exec_kmh: float, s_for_profile: float)|
              \item docstring: 明示されていない。
              \item このブロックが直接呼ぶ主な関数・メソッド: Float32MultiArray, \_route\_value, abs, electrical\_balance, float, get, publish
              \item 戻り値の要点: 戻り値が無いか、return 文が明示されていない。
              \item 主なローカル代入名: I, I\_motor, P\_mech\_wheel, P\_motor\_elec, P\_pack, P\_pv, V, headwind\_ms, msg, out, slope\_pct
              \item 読み取る主なインスタンス属性: self.Tb, self.\_route\_value, self.model, self.pub\_metrics, self.route\_profile, self.z
              \item 更新する主なインスタンス属性: 特になし。
              \item 明示的に送出する例外: 明示的raiseはない。
              \item 制御構造の規模: 条件分岐 2、ループ 0、try 0
            \end{itemize}
            \paragraph{この定義を読むためのPython構文}
            \begin{itemize}[leftmargin=1.6em]
            \item 第1引数selfは、このメソッドを呼んだインスタンス自身を受け取る慣習名である。
            \end{itemize}
            \paragraph{上から順の処理}
            \begin{enumerate}[leftmargin=1.8em]
            \item slope\_pct に d0.get('slope\_pct', 0.0) の結果を代入する。
\item 条件 self.route\_profile is not None を判定し、真なら内部処理を行う。
\item   slope\_pct に self.\_route\_value(s\_for\_profile, 'slope\_pct', slope\_pct) の結果を代入する。
\item headwind\_ms に d0.get('headwind\_ms', 0.0) の結果を代入する。
\item 条件 self.route\_profile is not None を判定し、真なら内部処理を行う。
\item   headwind\_ms に self.\_route\_value(s\_for\_profile, 'headwind\_ms', headwind\_ms) の結果を代入する。
\item out に self.model.electrical\_balance(v\_exec\_kmh / 3.6, slope\_pct, self.z, self.Tb, d0.get('G\_poa', 0.0), d0.get('Tcell\_C', 40.0), headwind\_ms=headwind\_ms, ambient\_temp\_c=d0.get('Tamb\_C'), elevation\_m=d0.get('elevation\_m', self.\_route\_value(s\_for\_profile, 'elev\_m', 0.0))) の結果を代入する。
\item V に float(out['V']) の結果を代入する。
\item I に float(out['I']) の結果を代入する。
\item P\_pv に float(out['P\_pv']) の結果を代入する。
\item P\_pack に float(out['P\_pack']) の結果を代入する。
\item P\_mech\_wheel に float(out.get('P\_mech\_wheel', out.get('P\_mech', 0.0))) の結果を代入する。
\item P\_motor\_elec に float(out.get('P\_dc\_to\_drv', 0.0)) - float(out.get('P\_reg\_to\_dc', 0.0)) の結果を代入する。
\item I\_motor に P\_motor\_elec / V if abs(V) > 0.001 else 0.0 の結果を代入する。
\item msg に Float32MultiArray() の結果を代入する。
\item msg.data に [float(V), float(I), float(self.z), float(P\_motor\_elec), float(I\_motor), float(P\_pv), float(v\_exec\_kmh), float(P\_mech\_wheel), float(P\_pack)] の結果を代入する。
\item self.pub\_metrics.publish(...) を実行する。
            \end{enumerate}
            \paragraph{代表コード断片}
            \begin{lstlisting}[language=Python]
    def _publish_metrics(self, d0: dict, v_exec_kmh: float, s_for_profile: float):
        slope_pct = d0.get('slope_pct', 0.0)
        if self.route_profile is not None:
            slope_pct = self._route_value(s_for_profile, 'slope_pct', slope_pct)
        headwind_ms = d0.get('headwind_ms', 0.0)
        if self.route_profile is not None:
            headwind_ms = self._route_value(s_for_profile, 'headwind_ms', headwind_ms)
        out = self.model.electrical_balance(
            v_exec_kmh / 3.6,
            slope_pct,
            self.z,
            self.Tb,
            d0.get('G_poa', 0.0),
            d0.get('Tcell_C', 40.0),
            headwind_ms=headwind_ms,
            ambient_temp_c=d0.get('Tamb_C'),
            elevation_m=d0.get('elevation_m', self._route_value(s_for_profile, 'elev_m', 0.0)),
        )
        V = float(out['V'])
        I = float(out['I'])
        P_pv = float(out['P_pv'])
        P_pack = float(out['P_pack'])
        P_mech_wheel = float(out.get('P_mech_wheel', out.get('P_mech', 0.0)))
        P_motor_elec = float(out.get('P_dc_to_drv', 0.0)) - float(out.get('P_reg_to_dc', 0.0))
        I_motor = P_motor_elec / V if abs(V) > 1e-3 else 0.0
        msg = Float32MultiArray()
        msg.data = [
            float(V),
            float(I),
            float(self.z),
            float(P_motor_elec),
            float(I_motor),
            float(P_pv),
            float(v_exec_kmh),
            float(P_mech_wheel),
...
\end{lstlisting}

\subsection{L1897 関数 MPCNode.\_publish\_summary}
            \begin{itemize}[leftmargin=1.6em]
              \item 行範囲: L1897--L1927
              \item 所属: \path|MPCNode|
              \item 定義: \path|_publish_summary(self, _v_exec_kmh: float = math.nan) -> None|
              \item docstring: Publish current state even while a long upper solve is in progress.
              \item このブロックが直接呼ぶ主な関数・メソッド: Float32MultiArray, all, bool, float, isna, len, lower, max, monotonic, notna, now, publish
              \item 戻り値の要点: 戻り値が無いか、return 文が明示されていない。
              \item 主なローカル代入名: elapsed, has\_time, mode, now, sec\_to\_next, st, stop\_remaining\_sec, t\_next
              \item 読み取る主なインスタンス属性: self.Tb, self.completed\_control\_stops, self.control\_stop\_end\_monotonic, self.control\_stop\_hold, self.df, self.forecast\_start\_time, self.forecast\_time\_mode, self.forecast\_time\_offset, self.k, self.model, self.pub\_status, self.s\_km, self.z
              \item 更新する主なインスタンス属性: 特になし。
              \item 明示的に送出する例外: 明示的raiseはない。
              \item 制御構造の規模: 条件分岐 4、ループ 0、try 0
            \end{itemize}
            \paragraph{この定義を読むためのPython構文}
            \begin{itemize}[leftmargin=1.6em]
            \item 第1引数selfは、このメソッドを呼んだインスタンス自身を受け取る慣習名である。
\item `->`以後は戻り値の型注釈であり、実行時の値を自動変換する命令ではない。
            \end{itemize}
            \paragraph{上から順の処理}
            \begin{enumerate}[leftmargin=1.8em]
            \item mode に str(self.forecast\_time\_mode).lower() の結果を代入する。
\item has\_time に 'time' in self.df.columns and (not self.df['time'].isna().all()) の結果を代入する。
\item 条件 mode == 'auto' を判定し、真なら内部処理を行う。
\item   mode に 'absolute' if has\_time else 'relative' の結果を代入する。
\item 条件 has\_time and mode == 'absolute' を判定し、真なら内部処理を行う。
\item   条件 self.k + 1 < len(self.df) and pd.notna(self.df['time'].iloc[self.k + 1]) を判定し、真なら内部処理を行う。
\item     t\_next に self.df['time'].iloc[self.k + 1].to\_pydatetime() の結果を代入する。
\item     sec\_to\_next に max(0.0, (t\_next - datetime.now(timezone.utc)).total\_seconds()) の結果を代入する。
\item     上の条件が偽の場合:
\item     sec\_to\_next に 0.0 の結果を代入する。
\item   上の条件が偽の場合:
\item   now に datetime.now(timezone.utc) + timedelta(seconds=self.forecast\_time\_offset) の結果を代入する。
\item   elapsed に max(0.0, (now - self.forecast\_start\_time).total\_seconds()) の結果を代入する。
\item   sec\_to\_next に max(0.0, self.model.p.dt - elapsed \% self.model.p.dt) の結果を代入する。
\item stop\_remaining\_sec に 0.0 の結果を代入する。
\item 条件 self.control\_stop\_end\_monotonic is not None を判定し、真なら内部処理を行う。
\item   stop\_remaining\_sec に max(0.0, self.control\_stop\_end\_monotonic - time.monotonic()) の結果を代入する。
\item st に Float32MultiArray() の結果を代入する。
\item st.data に [float(self.z), float(self.Tb), float(self.s\_km), float(self.k), float(sec\_to\_next), float(bool(self.control\_stop\_hold)), float(stop\_remaining\_sec), float(len(self.completed\_control\_stops))] の結果を代入する。
\item self.pub\_status.publish(...) を実行する。
            \end{enumerate}
            \paragraph{代表コード断片}
            \begin{lstlisting}[language=Python]
    def _publish_summary(self, _v_exec_kmh: float = math.nan) -> None:
        """Publish current state even while a long upper solve is in progress."""
        mode = str(self.forecast_time_mode).lower()
        has_time = ('time' in self.df.columns) and (not self.df['time'].isna().all())
        if mode == 'auto':
            mode = 'absolute' if has_time else 'relative'
        if has_time and mode == 'absolute':
            if self.k + 1 < len(self.df) and pd.notna(self.df['time'].iloc[self.k + 1]):
                t_next = self.df['time'].iloc[self.k + 1].to_pydatetime()
                sec_to_next = max(0.0, (t_next - datetime.now(timezone.utc)).total_seconds())
            else:
                sec_to_next = 0.0
        else:
            now = datetime.now(timezone.utc) + timedelta(seconds=self.forecast_time_offset)
            elapsed = max(0.0, (now - self.forecast_start_time).total_seconds())
            sec_to_next = max(0.0, self.model.p.dt - (elapsed % self.model.p.dt))
        stop_remaining_sec = 0.0
        if self.control_stop_end_monotonic is not None:
            stop_remaining_sec = max(0.0, self.control_stop_end_monotonic - time.monotonic())
        st = Float32MultiArray()
        st.data = [
            float(self.z),
            float(self.Tb),
            float(self.s_km),
            float(self.k),
            float(sec_to_next),
            float(bool(self.control_stop_hold)),
            float(stop_remaining_sec),
            float(len(self.completed_control_stops)),
        ]
        self.pub_status.publish(st)
\end{lstlisting}

\subsection{L1929 関数 MPCNode.\_interp\_upper\_speed}
            \begin{itemize}[leftmargin=1.6em]
              \item 行範囲: L1929--L1950
              \item 所属: \path|MPCNode|
              \item 定義: \path|_interp_upper_speed(self, t_sec: float) -> float|
              \item docstring: 明示されていない。
              \item このブロックが直接呼ぶ主な関数・メソッド: float, floor, int, len
              \item 戻り値の要点: float((1.0 - alpha) * self.v\_plan\_kmh[i] + alpha * self.v\_plan\_kmh[i + 1]) / float(self.v\_plan\_segments[-1]['v\_kmh']) / float(self.v\_upper\_cmd) / float(self.v\_plan\_kmh[0])
              \item 主なローカル代入名: acc, acc\_next, alpha, dt, i, idx, seg
              \item 読み取る主なインスタンス属性: self.plan\_dt\_sec, self.upper\_plan\_mode, self.v\_plan\_kmh, self.v\_plan\_segments, self.v\_upper\_cmd
              \item 更新する主なインスタンス属性: 特になし。
              \item 明示的に送出する例外: 明示的raiseはない。
              \item 制御構造の規模: 条件分岐 6、ループ 1、try 0
            \end{itemize}
            \paragraph{この定義を読むためのPython構文}
            \begin{itemize}[leftmargin=1.6em]
            \item 第1引数selfは、このメソッドを呼んだインスタンス自身を受け取る慣習名である。
\item `->`以後は戻り値の型注釈であり、実行時の値を自動変換する命令ではない。
            \end{itemize}
            \paragraph{上から順の処理}
            \begin{enumerate}[leftmargin=1.8em]
            \item 条件 self.upper\_plan\_mode == 'distance' and self.v\_plan\_segments を判定し、真なら内部処理を行う。
\item   acc に 0.0 の結果を代入する。
\item   self.v\_plan\_segments を順に走査し、各要素を seg に入れて処理する。
\item     acc\_next に acc + seg['dt\_sec'] の結果を代入する。
\item     条件 t\_sec <= acc\_next を判定し、真なら内部処理を行う。
\item       float(seg['v\_kmh']) を返す。
\item     acc に acc\_next の結果を代入する。
\item   float(self.v\_plan\_segments[-1]['v\_kmh']) を返す。
\item 条件 not self.v\_plan\_kmh を判定し、真なら内部処理を行う。
\item   float(self.v\_upper\_cmd) を返す。
\item dt に float(self.plan\_dt\_sec) の結果を代入する。
\item 条件 dt <= 0.0 を判定し、真なら内部処理を行う。
\item   float(self.v\_plan\_kmh[0]) を返す。
\item idx に t\_sec / dt の結果を代入する。
\item i に int(math.floor(idx)) の結果を代入する。
\item 条件 i <= 0 を判定し、真なら内部処理を行う。
\item   float(self.v\_plan\_kmh[0]) を返す。
\item 条件 i >= len(self.v\_plan\_kmh) - 1 を判定し、真なら内部処理を行う。
\item   float(self.v\_plan\_kmh[-1]) を返す。
\item alpha に idx - i の結果を代入する。
\item float((1.0 - alpha) * self.v\_plan\_kmh[i] + alpha * self.v\_plan\_kmh[i + 1]) を返す。
            \end{enumerate}
            \paragraph{代表コード断片}
            \begin{lstlisting}[language=Python]
    def _interp_upper_speed(self, t_sec: float) -> float:
        if self.upper_plan_mode == 'distance' and self.v_plan_segments:
            acc = 0.0
            for seg in self.v_plan_segments:
                acc_next = acc + seg['dt_sec']
                if t_sec <= acc_next:
                    return float(seg['v_kmh'])
                acc = acc_next
            return float(self.v_plan_segments[-1]['v_kmh'])
        if not self.v_plan_kmh:
            return float(self.v_upper_cmd)
        dt = float(self.plan_dt_sec)
        if dt <= 0.0:
            return float(self.v_plan_kmh[0])
        idx = t_sec / dt
        i = int(math.floor(idx))
        if i <= 0:
            return float(self.v_plan_kmh[0])
        if i >= len(self.v_plan_kmh) - 1:
            return float(self.v_plan_kmh[-1])
        alpha = idx - i
        return float((1.0 - alpha) * self.v_plan_kmh[i] + alpha * self.v_plan_kmh[i + 1])
\end{lstlisting}

\subsection{L1952 関数 MPCNode.\_distance\_plan\_speed}
            \begin{itemize}[leftmargin=1.6em]
              \item 行範囲: L1952--L1956
              \item 所属: \path|MPCNode|
              \item 定義: \path|_distance_plan_speed(self, s_km: float) -> float|
              \item docstring: 明示されていない。
              \item このブロックが直接呼ぶ主な関数・メソッド: float, plan\_segment\_index
              \item 戻り値の要点: float(self.v\_plan\_segments[idx]['v\_kmh']) / float(self.v\_upper\_cmd)
              \item 主なローカル代入名: idx
              \item 読み取る主なインスタンス属性: self.v\_plan\_segments, self.v\_upper\_cmd
              \item 更新する主なインスタンス属性: 特になし。
              \item 明示的に送出する例外: 明示的raiseはない。
              \item 制御構造の規模: 条件分岐 1、ループ 0、try 0
            \end{itemize}
            \paragraph{この定義を読むためのPython構文}
            \begin{itemize}[leftmargin=1.6em]
            \item 第1引数selfは、このメソッドを呼んだインスタンス自身を受け取る慣習名である。
\item `->`以後は戻り値の型注釈であり、実行時の値を自動変換する命令ではない。
            \end{itemize}
            \paragraph{上から順の処理}
            \begin{enumerate}[leftmargin=1.8em]
            \item idx に plan\_segment\_index(self.v\_plan\_segments or [], s\_km) の結果を代入する。
\item 条件 idx < 0 を判定し、真なら内部処理を行う。
\item   float(self.v\_upper\_cmd) を返す。
\item float(self.v\_plan\_segments[idx]['v\_kmh']) を返す。
            \end{enumerate}
            \paragraph{代表コード断片}
            \begin{lstlisting}[language=Python]
    def _distance_plan_speed(self, s_km: float) -> float:
        idx = plan_segment_index(self.v_plan_segments or [], s_km)
        if idx < 0:
            return float(self.v_upper_cmd)
        return float(self.v_plan_segments[idx]['v_kmh'])
\end{lstlisting}

\subsection{L1958 関数 MPCNode.\_tau\_max\_for\_mode}
            \begin{itemize}[leftmargin=1.6em]
              \item 行範囲: L1958--L1963
              \item 所属: \path|MPCNode|
              \item 定義: \path|_tau_max_for_mode(self, maps, mode: str) -> float|
              \item docstring: 明示されていない。
              \item このブロックが直接呼ぶ主な関数・メソッド: float, max
              \item 戻り値の要点: float(max(maps[key][1])) / 0.0
              \item 主なローカル代入名: key
              \item 読み取る主なインスタンス属性: 特になし。
              \item 更新する主なインスタンス属性: 特になし。
              \item 明示的に送出する例外: 明示的raiseはない。
              \item 制御構造の規模: 条件分岐 0、ループ 0、try 1
            \end{itemize}
            \paragraph{この定義を読むためのPython構文}
            \begin{itemize}[leftmargin=1.6em]
            \item 第1引数selfは、このメソッドを呼んだインスタンス自身を受け取る慣習名である。
\item `->`以後は戻り値の型注釈であり、実行時の値を自動変換する命令ではない。
\item try文は例外が起き得る処理と、異常時または終了時の経路を分ける。
            \end{itemize}
            \paragraph{上から順の処理}
            \begin{enumerate}[leftmargin=1.8em]
            \item key に mode if mode in maps else 'default' の結果を代入する。
\item 例外処理を伴う try ブロックを実行する。
\item   float(max(maps[key][1])) を返す。
\item   Exceptionを捕捉した場合:
\item   0.0 を返す。
            \end{enumerate}
            \paragraph{代表コード断片}
            \begin{lstlisting}[language=Python]
    def _tau_max_for_mode(self, maps, mode: str) -> float:
        key = mode if mode in maps else 'default'
        try:
            return float(max(maps[key][1]))
        except Exception:
            return 0.0
\end{lstlisting}

\subsection{L1965 関数 MPCNode.\_tau\_limits}
            \begin{itemize}[leftmargin=1.6em]
              \item 行範囲: L1965--L1979
              \item 所属: \path|MPCNode|
              \item 定義: \path|_tau_limits(self)|
              \item docstring: 明示されていない。
              \item このブロックが直接呼ぶ主な関数・メソッド: \_tau\_max\_for\_mode, lower, str
              \item 戻り値の要点: (tau\_drive, tau\_regen)
              \item 主なローカル代入名: mode, tau\_drive, tau\_regen
              \item 読み取る主なインスタンス属性: self.\_tau\_max\_for\_mode, self.model
              \item 更新する主なインスタンス属性: 特になし。
              \item 明示的に送出する例外: 明示的raiseはない。
              \item 制御構造の規模: 条件分岐 2、ループ 0、try 0
            \end{itemize}
            \paragraph{この定義を読むためのPython構文}
            \begin{itemize}[leftmargin=1.6em]
            \item 第1引数selfは、このメソッドを呼んだインスタンス自身を受け取る慣習名である。
            \end{itemize}
            \paragraph{上から順の処理}
            \begin{enumerate}[leftmargin=1.8em]
            \item mode に str(self.model.drive\_mode or 'default').lower() の結果を代入する。
\item 条件 mode in ('eco', 'power') を判定し、真なら内部処理を行う。
\item   tau\_drive に self.\_tau\_max\_for\_mode(self.model.maps\_drive, mode) の結果を代入する。
\item   tau\_regen に self.\_tau\_max\_for\_mode(self.model.maps\_regen, mode) の結果を代入する。
\item   上の条件が偽の場合:
\item   tau\_drive に self.\_tau\_max\_for\_mode(self.model.maps\_drive, 'power') if 'power' in self.model.maps\_drive else self.\_tau\_max\_for\_mode(self.model.maps\_drive, 'eco') if 'eco' in self.model.maps\_drive else self.\_tau\_max\_for\_mode(self.model.maps\_drive, 'default') の結果を代入する。
\item   tau\_regen に self.\_tau\_max\_for\_mode(self.model.maps\_regen, 'power') if 'power' in self.model.maps\_regen else self.\_tau\_max\_for\_mode(self.model.maps\_regen, 'eco') if 'eco' in self.model.maps\_regen else self.\_tau\_max\_for\_mode(self.model.maps\_regen, 'default') の結果を代入する。
\item 条件 tau\_regen <= 0.0 を判定し、真なら内部処理を行う。
\item   tau\_regen に tau\_drive の結果を代入する。
\item (tau\_drive, tau\_regen) を返す。
            \end{enumerate}
            \paragraph{代表コード断片}
            \begin{lstlisting}[language=Python]
    def _tau_limits(self):
        mode = str(self.model.drive_mode or 'default').lower()
        if mode in ('eco', 'power'):
            tau_drive = self._tau_max_for_mode(self.model.maps_drive, mode)
            tau_regen = self._tau_max_for_mode(self.model.maps_regen, mode)
        else:
            tau_drive = self._tau_max_for_mode(self.model.maps_drive, 'power') if 'power' in self.model.maps_drive else \
                self._tau_max_for_mode(self.model.maps_drive, 'eco') if 'eco' in self.model.maps_drive else \
                self._tau_max_for_mode(self.model.maps_drive, 'default')
            tau_regen = self._tau_max_for_mode(self.model.maps_regen, 'power') if 'power' in self.model.maps_regen else \
                self._tau_max_for_mode(self.model.maps_regen, 'eco') if 'eco' in self.model.maps_regen else \
                self._tau_max_for_mode(self.model.maps_regen, 'default')
        if tau_regen <= 0.0:
            tau_regen = tau_drive
        return tau_drive, tau_regen
\end{lstlisting}

\subsection{L1981 関数 MPCNode.\_traction\_force}
            \begin{itemize}[leftmargin=1.6em]
              \item 行範囲: L1981--L1986
              \item 所属: \path|MPCNode|
              \item 定義: \path|_traction_force(self, tau_nm: float) -> float|
              \item docstring: 明示されていない。
              \item このブロックが直接呼ぶ主な関数・メソッド: float, int, max
              \item 戻り値の要点: float(tau\_nm) * float(p.gear\_ratio) * float(p.gear\_eta) * motor\_count / wheel\_r
              \item 主なローカル代入名: motor\_count, p, wheel\_r
              \item 読み取る主なインスタンス属性: self.model
              \item 更新する主なインスタンス属性: 特になし。
              \item 明示的に送出する例外: 明示的raiseはない。
              \item 制御構造の規模: 条件分岐 0、ループ 0、try 0
            \end{itemize}
            \paragraph{この定義を読むためのPython構文}
            \begin{itemize}[leftmargin=1.6em]
            \item 第1引数selfは、このメソッドを呼んだインスタンス自身を受け取る慣習名である。
\item `->`以後は戻り値の型注釈であり、実行時の値を自動変換する命令ではない。
            \end{itemize}
            \paragraph{上から順の処理}
            \begin{enumerate}[leftmargin=1.8em]
            \item p に self.model.p の結果を代入する。
\item wheel\_r に max(0.001, float(p.wheel\_radius)) の結果を代入する。
\item motor\_count に int(p.motor\_count) if int(p.motor\_count) > 0 else int(p.driven\_wheel\_count or 1) の結果を代入する。
\item motor\_count に max(1, motor\_count) の結果を代入する。
\item float(tau\_nm) * float(p.gear\_ratio) * float(p.gear\_eta) * motor\_count / wheel\_r を返す。
            \end{enumerate}
            \paragraph{代表コード断片}
            \begin{lstlisting}[language=Python]
    def _traction_force(self, tau_nm: float) -> float:
        p = self.model.p
        wheel_r = max(1e-3, float(p.wheel_radius))
        motor_count = int(p.motor_count) if int(p.motor_count) > 0 else int(p.driven_wheel_count or 1)
        motor_count = max(1, motor_count)
        return float(tau_nm) * float(p.gear_ratio) * float(p.gear_eta) * motor_count / wheel_r
\end{lstlisting}

\subsection{L1988 関数 MPCNode.\_pack\_from\_tau}
            \begin{itemize}[leftmargin=1.6em]
              \item 行範囲: L1988--L2010
              \item 所属: \path|MPCNode|
              \item 定義: \path|_pack_from_tau(self, v_ms: float, tau_nm: float, z: float, Tb: float, env: dict)|
              \item docstring: 明示されていない。
              \item このブロックが直接呼ぶ主な関数・メソッド: battery\_iv, dict, eff\_drive, eff\_regen, float, int, max, pv\_power\_mppt
              \item 戻り値の要点: dict(P\_pack=P\_pack, I=I, V=V, loss\_int=loss\_int, eff=eff, P\_pv=P\_pv, P\_elec=P\_elec)
              \item 主なローカル代入名: I, P\_elec, P\_mech\_motor, P\_pack, P\_pv, Rint, V, eff, iv, loss\_int, motor\_count, omega\_m, p, wheel\_r
              \item 読み取る主なインスタンス属性: self.model
              \item 更新する主なインスタンス属性: 特になし。
              \item 明示的に送出する例外: 明示的raiseはない。
              \item 制御構造の規模: 条件分岐 1、ループ 0、try 0
            \end{itemize}
            \paragraph{この定義を読むためのPython構文}
            \begin{itemize}[leftmargin=1.6em]
            \item 第1引数selfは、このメソッドを呼んだインスタンス自身を受け取る慣習名である。
            \end{itemize}
            \paragraph{上から順の処理}
            \begin{enumerate}[leftmargin=1.8em]
            \item p に self.model.p の結果を代入する。
\item wheel\_r に max(0.001, float(p.wheel\_radius)) の結果を代入する。
\item motor\_count に int(p.motor\_count) if int(p.motor\_count) > 0 else int(p.driven\_wheel\_count or 1) の結果を代入する。
\item motor\_count に max(1, motor\_count) の結果を代入する。
\item omega\_m に float(v\_ms) / wheel\_r * float(p.gear\_ratio) の結果を代入する。
\item P\_mech\_motor に float(tau\_nm) * omega\_m * motor\_count の結果を代入する。
\item 条件 tau\_nm >= 0.0 を判定し、真なら内部処理を行う。
\item   eff に float(self.model.eff\_drive(v\_ms, tau\_nm)) の結果を代入する。
\item   eff に max(0.001, eff) の結果を代入する。
\item   P\_elec に P\_mech\_motor / (eff * float(p.inverter\_eta)) の結果を代入する。
\item   上の条件が偽の場合:
\item   eff に float(self.model.eff\_regen(v\_ms, tau\_nm)) の結果を代入する。
\item   eff に max(0.001, eff) の結果を代入する。
\item   P\_elec に P\_mech\_motor * eff * float(p.inverter\_eta) の結果を代入する。
\item P\_pv に float(self.model.pv\_power\_mppt(env['G\_poa'], env['Tcell\_C'])) の結果を代入する。
\item P\_pack に P\_elec + float(p.P\_aux) - P\_pv の結果を代入する。
\item iv に self.model.battery\_iv(P\_pack, z, Tb) の結果を代入する。
\item I に float(iv['I']) の結果を代入する。
\item V に float(iv['V']) の結果を代入する。
\item Rint に float(iv['Rint']) の結果を代入する。
\item loss\_int に I * I * Rint の結果を代入する。
\item dict(P\_pack=P\_pack, I=I, V=V, loss\_int=loss\_int, eff=eff, P\_pv=P\_pv, P\_elec=P\_elec) を返す。
            \end{enumerate}
            \paragraph{代表コード断片}
            \begin{lstlisting}[language=Python]
    def _pack_from_tau(self, v_ms: float, tau_nm: float, z: float, Tb: float, env: dict):
        p = self.model.p
        wheel_r = max(1e-3, float(p.wheel_radius))
        motor_count = int(p.motor_count) if int(p.motor_count) > 0 else int(p.driven_wheel_count or 1)
        motor_count = max(1, motor_count)
        omega_m = float(v_ms) / wheel_r * float(p.gear_ratio)
        P_mech_motor = float(tau_nm) * omega_m * motor_count
        if tau_nm >= 0.0:
            eff = float(self.model.eff_drive(v_ms, tau_nm))
            eff = max(1.0e-3, eff)
            P_elec = P_mech_motor / (eff * float(p.inverter_eta))
        else:
            eff = float(self.model.eff_regen(v_ms, tau_nm))
            eff = max(1.0e-3, eff)
            P_elec = P_mech_motor * eff * float(p.inverter_eta)
        P_pv = float(self.model.pv_power_mppt(env['G_poa'], env['Tcell_C']))
        P_pack = P_elec + float(p.P_aux) - P_pv
        iv = self.model.battery_iv(P_pack, z, Tb)
        I = float(iv['I'])
        V = float(iv['V'])
        Rint = float(iv['Rint'])
        loss_int = I * I * Rint
        return dict(P_pack=P_pack, I=I, V=V, loss_int=loss_int, eff=eff, P_pv=P_pv, P_elec=P_elec)
\end{lstlisting}

\subsection{L2012 関数 MPCNode.\_build\_lower\_ref}
            \begin{itemize}[leftmargin=1.6em]
              \item 行範囲: L2012--L2037
              \item 所属: \path|MPCNode|
              \item 定義: \path|_build_lower_ref(self, base_time_utc, s_km: float, d0: dict)|
              \item docstring: 明示されていない。
              \item このブロックが直接呼ぶ主な関数・メソッド: \_interp\_upper\_speed, \_soc\_guard\_speed, \_speed\_limit\_at, append, clip, float, get\_parameter, max, min, monotonic, range, speed\_limits
              \item 戻り値の要点: ref
              \item 主なローカル代入名: guard\_speed, i, limits, offset, ref, s\_tmp, soc\_guard, t\_sec, t\_utc, v\_ref, vmax\_kmh, vmax\_local, vmin\_kmh
              \item 読み取る主なインスタンス属性: self.\_interp\_upper\_speed, self.\_soc\_guard\_speed, self.\_speed\_limit\_at, self.drive\_schedule, self.get\_parameter, self.lower\_N, self.lower\_dt, self.model, self.plan\_start\_monotonic, self.v\_upper\_cmd, self.z
              \item 更新する主なインスタンス属性: 特になし。
              \item 明示的に送出する例外: 明示的raiseはない。
              \item 制御構造の規模: 条件分岐 5、ループ 1、try 0
            \end{itemize}
            \paragraph{この定義を読むためのPython構文}
            \begin{itemize}[leftmargin=1.6em]
            \item 第1引数selfは、このメソッドを呼んだインスタンス自身を受け取る慣習名である。
            \end{itemize}
            \paragraph{上から順の処理}
            \begin{enumerate}[leftmargin=1.8em]
            \item ref に [] の結果を代入する。
\item s\_tmp に float(s\_km) の結果を代入する。
\item offset に 0.0 の結果を代入する。
\item 条件 self.plan\_start\_monotonic is not None を判定し、真なら内部処理を行う。
\item   offset に max(0.0, time.monotonic() - self.plan\_start\_monotonic) の結果を代入する。
\item soc\_guard に float(self.get\_parameter('soc\_guard\_margin').value) の結果を代入する。
\item guard\_speed に None の結果を代入する。
\item 条件 self.z <= self.model.p.soc\_min + soc\_guard を判定し、真なら内部処理を行う。
\item   guard\_speed に self.\_soc\_guard\_speed(self.v\_upper\_cmd, s\_tmp, d0) の結果を代入する。
\item range(max(1, self.lower\_N)) を順に走査し、各要素を i に入れて処理する。
\item   t\_sec に offset + i * self.lower\_dt の結果を代入する。
\item   v\_ref に self.\_interp\_upper\_speed(t\_sec) の結果を代入する。
\item   条件 guard\_speed is not None を判定し、真なら内部処理を行う。
\item     v\_ref に min(v\_ref, guard\_speed) の結果を代入する。
\item   vmax\_local に self.\_speed\_limit\_at(s\_tmp, self.get\_parameter('v\_max\_kmh').value) の結果を代入する。
\item   v\_ref に min(v\_ref, vmax\_local) の結果を代入する。
\item   条件 self.drive\_schedule is not None and base\_time\_utc is not None を判定し、真なら内部処理を行う。
\item     t\_utc に base\_time\_utc + timedelta(seconds=t\_sec) の結果を代入する。
\item     limits に self.drive\_schedule.speed\_limits(t\_utc) の結果を代入する。
\item     条件 limits is not None を判定し、真なら内部処理を行う。
\item       (vmin\_kmh, vmax\_kmh) に limits の結果を代入する。
\item       v\_ref に float(np.clip(v\_ref, vmin\_kmh, vmax\_kmh)) の結果を代入する。
\item   ref.append(...) を実行する。
\item   s\_tmp を Add で更新する。
\item ref を返す。
            \end{enumerate}
            \paragraph{代表コード断片}
            \begin{lstlisting}[language=Python]
    def _build_lower_ref(self, base_time_utc, s_km: float, d0: dict):
        ref = []
        s_tmp = float(s_km)
        offset = 0.0
        if self.plan_start_monotonic is not None:
            offset = max(0.0, time.monotonic() - self.plan_start_monotonic)
        soc_guard = float(self.get_parameter('soc_guard_margin').value)
        guard_speed = None
        if self.z <= (self.model.p.soc_min + soc_guard):
            guard_speed = self._soc_guard_speed(self.v_upper_cmd, s_tmp, d0)
        for i in range(max(1, self.lower_N)):
            t_sec = offset + i * self.lower_dt
            v_ref = self._interp_upper_speed(t_sec)
            if guard_speed is not None:
                v_ref = min(v_ref, guard_speed)
            vmax_local = self._speed_limit_at(s_tmp, self.get_parameter('v_max_kmh').value)
            v_ref = min(v_ref, vmax_local)
            if self.drive_schedule is not None and base_time_utc is not None:
                t_utc = base_time_utc + timedelta(seconds=t_sec)
                limits = self.drive_schedule.speed_limits(t_utc)
                if limits is not None:
                    vmin_kmh, vmax_kmh = limits
                    v_ref = float(np.clip(v_ref, vmin_kmh, vmax_kmh))
            ref.append(float(v_ref) / 3.6)
            s_tmp += float(v_ref) * (self.lower_dt / 3600.0)
        return ref
\end{lstlisting}

\subsection{L2039 関数 MPCNode.\_lower\_rollout}
            \begin{itemize}[leftmargin=1.6em]
              \item 行範囲: L2039--L2070
              \item 所属: \path|MPCNode|
              \item 定義: \path|_lower_rollout(self, v0_ms: float, u_seq: np.ndarray, env: dict, z0: float, Tb0: float, tau_drive: float, tau_regen: float)|
              \item docstring: 明示されていない。
              \item このブロックが直接呼ぶ主な関数・メソッド: \_pack\_from\_tau, \_traction\_force, append, float, get, max, resistive\_forces, soc\_step
              \item 戻り値の要点: v\_seq
              \item 主なローカル代入名: F\_res, F\_trac, Tb, a, forces, p, pack, tau, u, v, v\_seq, z
              \item 読み取る主なインスタンス属性: self.\_pack\_from\_tau, self.\_traction\_force, self.lower\_dt, self.model
              \item 更新する主なインスタンス属性: 特になし。
              \item 明示的に送出する例外: 明示的raiseはない。
              \item 制御構造の規模: 条件分岐 1、ループ 1、try 0
            \end{itemize}
            \paragraph{この定義を読むためのPython構文}
            \begin{itemize}[leftmargin=1.6em]
            \item 第1引数selfは、このメソッドを呼んだインスタンス自身を受け取る慣習名である。
            \end{itemize}
            \paragraph{上から順の処理}
            \begin{enumerate}[leftmargin=1.8em]
            \item p に self.model.p の結果を代入する。
\item v に float(v0\_ms) の結果を代入する。
\item z に float(z0) の結果を代入する。
\item Tb に float(Tb0) の結果を代入する。
\item v\_seq に [] の結果を代入する。
\item u\_seq を順に走査し、各要素を u に入れて処理する。
\item   u に float(u) の結果を代入する。
\item   条件 u >= 0.0 を判定し、真なら内部処理を行う。
\item     tau に u * tau\_drive の結果を代入する。
\item     上の条件が偽の場合:
\item     tau に u * tau\_regen の結果を代入する。
\item   forces に self.model.resistive\_forces(v, env['slope\_pct'], env['headwind\_ms'], ambient\_temp\_c=env.get('Tamb\_C'), elevation\_m=env.get('elevation\_m', 0.0)) の結果を代入する。
\item   F\_res に float(forces['F\_total']) の結果を代入する。
\item   F\_trac に self.\_traction\_force(tau) の結果を代入する。
\item   a に (F\_trac - F\_res) / float(p.m) の結果を代入する。
\item   v に max(0.0, v + a * self.lower\_dt) の結果を代入する。
\item   pack に self.\_pack\_from\_tau(v, tau, z, Tb, env) の結果を代入する。
\item   z に self.model.soc\_step(z, float(pack['P\_pack']), self.lower\_dt, current\_a=float(pack['I']), Tbat\_C=Tb) の結果を代入する。
\item   Tb に Tb + self.lower\_dt / 1800.0 * (env['Tamb\_C'] - Tb) + float(pack['loss\_int']) * self.lower\_dt / 50000.0 の結果を代入する。
\item   v\_seq.append(...) を実行する。
\item v\_seq を返す。
            \end{enumerate}
            \paragraph{代表コード断片}
            \begin{lstlisting}[language=Python]
    def _lower_rollout(self, v0_ms: float, u_seq: np.ndarray, env: dict, z0: float, Tb0: float,
                       tau_drive: float, tau_regen: float):
        p = self.model.p
        v = float(v0_ms)
        z = float(z0)
        Tb = float(Tb0)
        v_seq = []
        for u in u_seq:
            u = float(u)
            if u >= 0.0:
                tau = u * tau_drive
            else:
                tau = u * tau_regen
            forces = self.model.resistive_forces(
                v, env['slope_pct'], env['headwind_ms'],
                ambient_temp_c=env.get('Tamb_C'), elevation_m=env.get('elevation_m', 0.0),
            )
            F_res = float(forces['F_total'])
            F_trac = self._traction_force(tau)
            a = (F_trac - F_res) / float(p.m)
            v = max(0.0, v + a * self.lower_dt)
            pack = self._pack_from_tau(v, tau, z, Tb, env)
            z = self.model.soc_step(
                z,
                float(pack['P_pack']),
                self.lower_dt,
                current_a=float(pack['I']),
                Tbat_C=Tb,
            )
            Tb = Tb + (self.lower_dt / 1800.0) * (env['Tamb_C'] - Tb) + (float(pack['loss_int']) * self.lower_dt) / 50000.0
            v_seq.append(float(v))
        return v_seq
\end{lstlisting}

\subsection{L2072 関数 MPCNode.\_lower\_mpc\_solve}
            \begin{itemize}[leftmargin=1.6em]
              \item 行範囲: L2072--L2271
              \item 所属: \path|MPCNode|
              \item 定義: \path|_lower_mpc_solve(self, v0_ms: float, s0_km: float, z0: float, Tb0: float, env: dict, v_ref_seq)|
              \item docstring: 明示されていない。
              \item このブロックが直接呼ぶ主な関数・メソッド: \_lower\_rollout, \_pack\_from\_tau, \_tau\_limits, \_traction\_force, abs, all, append, array, asarray, bool, clip, concatenate
              \item 戻り値の要点: (u\_seq, v\_pred, mode) / (np.zeros(1), [v0\_ms], 'eco') / J
              \item 主なローカル代入名: F\_res, F\_trac, I, J, N, P\_pack, Tb, V, a, best\_fallback, best\_feasible, best\_feasible\_score, best\_violation\_score, bounds, candidate\_controls, candidate\_u, current, desired\_accel, du, feasible
              \item 読み取る主なインスタンス属性: self.\_lower\_rollout, self.\_lower\_upper\_busy\_logged, self.\_pack\_from\_tau, self.\_tau\_limits, self.\_traction\_force, self.\_upper\_solve\_in\_progress, self.get\_logger, self.get\_parameter, self.lower\_N, self.lower\_dt, self.lower\_last\_u, self.lower\_max\_iter, self.lower\_maxfun, self.lower\_skip\_optimization\_during\_upper\_solve, self.model, self.throttle\_rate\_limit, self.w\_throttle, self.w\_track
              \item 更新する主なインスタンス属性: self.\_lower\_upper\_busy\_logged, self.last\_lower\_solve\_ok
              \item 明示的に送出する例外: 明示的raiseはない。
              \item 制御構造の規模: 条件分岐 9、ループ 3、try 0
            \end{itemize}
            \paragraph{この定義を読むためのPython構文}
            \begin{itemize}[leftmargin=1.6em]
            \item 第1引数selfは、このメソッドを呼んだインスタンス自身を受け取る慣習名である。
            \end{itemize}
            \paragraph{上から順の処理}
            \begin{enumerate}[leftmargin=1.8em]
            \item N に max(1, min(self.lower\_N, len(v\_ref\_seq))) の結果を代入する。
\item 条件 N <= 0 を判定し、真なら内部処理を行う。
\item   (np.zeros(1), [v0\_ms], 'eco') を返す。
\item v\_ref\_seq に v\_ref\_seq[:N] の結果を代入する。
\item (tau\_drive, tau\_regen) に self.\_tau\_limits() の結果を代入する。
\item p に self.model.p の結果を代入する。
\item w\_track に float(self.w\_track) の結果を代入する。
\item w\_throttle に float(self.w\_throttle) の結果を代入する。
\item w\_current に float(self.get\_parameter('w\_current').value) の結果を代入する。
\item w\_T に float(self.get\_parameter('w\_T').value) の結果を代入する。
\item rate\_lim に float(self.throttle\_rate\_limit) / 100.0 if self.throttle\_rate\_limit > 0.0 else 0.0 の結果を代入する。
\item u0 に float(np.clip(self.lower\_last\_u, -1.0, 1.0)) の結果を代入する。
\item bounds に [(-1.0, 1.0)] * N の結果を代入する。
\item seed に [] の結果を代入する。
\item v\_seed に float(v0\_ms) の結果を代入する。
\item z\_seed に float(z0) の結果を代入する。
\item tb\_seed に float(Tb0) の結果を代入する。
\item u\_prev\_seed に u0 の結果を代入する。
\item motor\_count に int(p.motor\_count) if int(p.motor\_count) > 0 else int(p.driven\_wheel\_count or 1) の結果を代入する。
\item motor\_count に max(1, motor\_count) の結果を代入する。
\item torque\_factor に max(0.001, float(p.wheel\_radius)) / max(1e-06, float(p.gear\_ratio) * float(p.gear\_eta) * motor\_count) の結果を代入する。
\item max\_du に rate\_lim * max(self.lower\_dt, 0.001) if rate\_lim > 0.0 else 2.0 の結果を代入する。
\item v\_ref\_seq を順に走査し、各要素を v\_ref に入れて処理する。
\item   forces に self.model.resistive\_forces(v\_seed, env['slope\_pct'], env['headwind\_ms'], ambient\_temp\_c=env.get('Tamb\_C'), elevation\_m=env.get('elevation\_m', 0.0)) の結果を代入する。
\item   desired\_accel に (float(v\_ref) - v\_seed) / max(self.lower\_dt, 0.001) の結果を代入する。
\item   tau\_request に (float(p.m) * desired\_accel + float(forces['F\_total'])) * torque\_factor の結果を代入する。
\item   tau\_limit に tau\_drive if tau\_request >= 0.0 else tau\_regen の結果を代入する。
\item   u\_request に tau\_request / max(abs(float(tau\_limit)), 1e-06) の結果を代入する。
\item   u\_request に float(np.clip(u\_request, u\_prev\_seed - max\_du, u\_prev\_seed + max\_du)) の結果を代入する。
\item   u\_request に float(np.clip(u\_request, -1.0, 1.0)) の結果を代入する。
\item   u\_min に max(-1.0, u\_prev\_seed - max\_du) の結果を代入する。
\item   u\_max に min(1.0, u\_prev\_seed + max\_du) の結果を代入する。
\item   candidate\_controls に np.unique(np.concatenate((np.linspace(u\_min, u\_max, 25, dtype=float), np.array([u\_request, np.clip(0.0, u\_min, u\_max)], dtype=float)))) の結果を代入する。
\item   best\_feasible に None の結果を代入する。
\item   best\_feasible\_score に float('inf') の結果を代入する。
\item   best\_fallback に float(np.clip(u\_request, u\_min, u\_max)) の結果を代入する。
\item   best\_violation\_score に float('inf') の結果を代入する。
\item   candidate\_controls を順に走査し、各要素を candidate\_u に入れて処理する。
\item     candidate\_u に float(np.clip(candidate\_u, u\_min, u\_max)) の結果を代入する。
\item     tau に candidate\_u * (tau\_drive if candidate\_u >= 0.0 else tau\_regen) の結果を代入する。
\item     F\_trac に self.\_traction\_force(tau) の結果を代入する。
\item     v\_trial に max(0.0, v\_seed + (F\_trac - float(forces['F\_total'])) / float(p.m) * self.lower\_dt) の結果を代入する。
\item     pack に self.\_pack\_from\_tau(v\_trial, tau, z\_seed, tb\_seed, env) の結果を代入する。
\item     current に float(pack['I']) の結果を代入する。
\item     voltage に float(pack['V']) の結果を代入する。
\item     tracking\_score に (v\_trial - float(v\_ref)) ** 2 + 0.01 * (candidate\_u - u\_request) ** 2 の結果を代入する。
\item     feasible に float(p.I\_chg\_min) <= current <= float(p.I\_max) and float(p.V\_min) <= voltage <= float(p.V\_max) の結果を代入する。
\item     条件 feasible and tracking\_score < best\_feasible\_score を判定し、真なら内部処理を行う。
\item       best\_feasible に candidate\_u の結果を代入する。
\item       best\_feasible\_score に tracking\_score の結果を代入する。
\item     violation\_score に (max(0.0, current - float(p.I\_max)) / max(abs(float(p.I\_max)), 1.0)) ** 2 + (max(0.0, float(p.I\_chg\_min) - current) / max(abs(float(p.I\_chg\_min)), 1.0)) ** 2 + (max(0.0, float(p.V\_min) - voltage) / max(abs(float(p.V\_min)), 1.0)) ** 2 + (max(0.0, voltage - float(p.V\_max)) / max(abs(float(p.V\_max)), 1.0)) ** 2 + 0.001 * tracking\_score の結果を代入する。
\item     条件 violation\_score < best\_violation\_score を判定し、真なら内部処理を行う。
\item       best\_violation\_score に violation\_score の結果を代入する。
\item       best\_fallback に candidate\_u の結果を代入する。
\item   u\_request に float(best\_feasible if best\_feasible is not None else best\_fallback) の結果を代入する。
\item   seed.append(...) を実行する。
\item   tau に u\_request * (tau\_drive if u\_request >= 0.0 else tau\_regen) の結果を代入する。
\item   F\_trac に self.\_traction\_force(tau) の結果を代入する。
\item   v\_seed に max(0.0, v\_seed + (F\_trac - float(forces['F\_total'])) / float(p.m) * self.lower\_dt) の結果を代入する。
\item   pack に self.\_pack\_from\_tau(v\_seed, tau, z\_seed, tb\_seed, env) の結果を代入する。
\item   z\_seed に self.model.soc\_step(z\_seed, float(pack['P\_pack']), self.lower\_dt, current\_a=float(pack['I']), Tbat\_C=tb\_seed) の結果を代入する。
\item   tb\_seed を Add で更新する。
\item   u\_prev\_seed に u\_request の結果を代入する。
\item x0 に np.asarray(seed, dtype=float) の結果を代入する。
\item 関数 cost を定義する。
\item skip\_refine に bool(self.\_upper\_solve\_in\_progress and self.lower\_skip\_optimization\_during\_upper\_solve) の結果を代入する。
\item 条件 skip\_refine or self.lower\_max\_iter <= 0 を判定し、真なら内部処理を行う。
\item   self.last\_lower\_solve\_ok に True の結果を代入する。
\item   u\_seq に x0 の結果を代入する。
\item   条件 skip\_refine and (not self.\_lower\_upper\_busy\_logged) を判定し、真なら内部処理を行う。
\item     self.get\_logger().info(...) を実行する。
\item     self.\_lower\_upper\_busy\_logged に True の結果を代入する。
\item   上の条件が偽の場合:
\item   res に minimize(cost, x0, method='L-BFGS-B', bounds=bounds, options=dict(maxiter=self.lower\_max\_iter, maxfun=self.lower\_maxfun, ftol=1e-06)) の結果を代入する。
\item   条件 np.all(np.isfinite(res.x)) を判定し、真なら内部処理を行う。
\item     self.last\_lower\_solve\_ok に bool(res.success) の結果を代入する。
\item     u\_seq に res.x の結果を代入する。
\item     条件 not res.success を判定し、真なら内部処理を行う。
\item       self.get\_logger().warn(...) を実行する。
\item     上の条件が偽の場合:
            \end{enumerate}
            \paragraph{代表コード断片}
            \begin{lstlisting}[language=Python]
    def _lower_mpc_solve(self, v0_ms: float, s0_km: float, z0: float, Tb0: float, env: dict, v_ref_seq):
        N = max(1, min(self.lower_N, len(v_ref_seq)))
        if N <= 0:
            return np.zeros(1), [v0_ms], 'eco'
        v_ref_seq = v_ref_seq[:N]
        tau_drive, tau_regen = self._tau_limits()
        p = self.model.p
        w_track = float(self.w_track)
        w_throttle = float(self.w_throttle)
        w_current = float(self.get_parameter('w_current').value)
        w_T = float(self.get_parameter('w_T').value)
        rate_lim = float(self.throttle_rate_limit) / 100.0 if self.throttle_rate_limit > 0.0 else 0.0
        u0 = float(np.clip(self.lower_last_u, -1.0, 1.0))
        bounds = [(-1.0, 1.0)] * N

        # Inverse dynamics gives a physically meaningful, bounded horizon seed in
        # deterministic time. It also remains the safe lower policy while the
        # hourly full-race optimizer is using the other executor thread.
        seed = []
        v_seed = float(v0_ms)
        z_seed = float(z0)
        tb_seed = float(Tb0)
        u_prev_seed = u0
        motor_count = int(p.motor_count) if int(p.motor_count) > 0 else int(p.driven_wheel_count or 1)
        motor_count = max(1, motor_count)
        torque_factor = (
            max(1.0e-3, float(p.wheel_radius))
            / max(1.0e-6, float(p.gear_ratio) * float(p.gear_eta) * motor_count)
        )
        max_du = rate_lim * max(self.lower_dt, 1.0e-3) if rate_lim > 0.0 else 2.0
        for v_ref in v_ref_seq:
            forces = self.model.resistive_forces(
                v_seed, env['slope_pct'], env['headwind_ms'],
                ambient_temp_c=env.get('Tamb_C'), elevation_m=env.get('elevation_m', 0.0),
            )
...
\end{lstlisting}

\subsection{L2185 関数 MPCNode.\_lower\_mpc\_solve.cost}
            \begin{itemize}[leftmargin=1.6em]
              \item 行範囲: L2185--L2237
              \item 所属: \path|MPCNode._lower_mpc_solve|
              \item 定義: \path|cost(u_vec)|
              \item docstring: 明示されていない。
              \item このブロックが直接呼ぶ主な関数・メソッド: \_pack\_from\_tau, \_traction\_force, abs, float, get, max, range, resistive\_forces, soc\_step
              \item 戻り値の要点: J
              \item 主なローカル代入名: F\_res, F\_trac, I, J, P\_pack, Tb, V, a, du, forces, i, loss\_int, pack, tau, u, u\_prev, v, v\_ref, z
              \item 読み取る主なインスタンス属性: self.\_pack\_from\_tau, self.\_traction\_force, self.lower\_dt, self.model
              \item 更新する主なインスタンス属性: 特になし。
              \item 明示的に送出する例外: 明示的raiseはない。
              \item 制御構造の規模: 条件分岐 2、ループ 1、try 0
            \end{itemize}
            \paragraph{この定義を読むためのPython構文}
            \begin{itemize}[leftmargin=1.6em]
            \item 特別な構文上の補足は少ない。
            \end{itemize}
            \paragraph{上から順の処理}
            \begin{enumerate}[leftmargin=1.8em]
            \item v に float(v0\_ms) の結果を代入する。
\item z に float(z0) の結果を代入する。
\item Tb に float(Tb0) の結果を代入する。
\item u\_prev に u0 の結果を代入する。
\item J に 0.0 の結果を代入する。
\item range(N) を順に走査し、各要素を i に入れて処理する。
\item   u に float(u\_vec[i]) の結果を代入する。
\item   条件 u >= 0.0 を判定し、真なら内部処理を行う。
\item     tau に u * tau\_drive の結果を代入する。
\item     上の条件が偽の場合:
\item     tau に u * tau\_regen の結果を代入する。
\item   forces に self.model.resistive\_forces(v, env['slope\_pct'], env['headwind\_ms'], ambient\_temp\_c=env.get('Tamb\_C'), elevation\_m=env.get('elevation\_m', 0.0)) の結果を代入する。
\item   F\_res に float(forces['F\_total']) の結果を代入する。
\item   F\_trac に self.\_traction\_force(tau) の結果を代入する。
\item   a に (F\_trac - F\_res) / float(p.m) の結果を代入する。
\item   v に max(0.0, v + a * self.lower\_dt) の結果を代入する。
\item   pack に self.\_pack\_from\_tau(v, tau, z, Tb, env) の結果を代入する。
\item   P\_pack に float(pack['P\_pack']) の結果を代入する。
\item   I に float(pack['I']) の結果を代入する。
\item   V に float(pack['V']) の結果を代入する。
\item   loss\_int に float(pack['loss\_int']) の結果を代入する。
\item   z に self.model.soc\_step(z, P\_pack, self.lower\_dt, current\_a=I, Tbat\_C=Tb) の結果を代入する。
\item   Tb に Tb + self.lower\_dt / 1800.0 * (env['Tamb\_C'] - Tb) + loss\_int * self.lower\_dt / 50000.0 の結果を代入する。
\item   v\_ref に float(v\_ref\_seq[i]) の結果を代入する。
\item   J を Add で更新する。
\item   J を Add で更新する。
\item   du に (u - u\_prev) / max(self.lower\_dt, 0.001) の結果を代入する。
\item   条件 rate\_lim > 0.0 を判定し、真なら内部処理を行う。
\item     J を Add で更新する。
\item   J を Add で更新する。
\item   J を Add で更新する。
\item   J を Add で更新する。
\item   J を Add で更新する。
\item   J を Add で更新する。
\item   J を Add で更新する。
\item   J を Add で更新する。
\item   J を Add で更新する。
\item   J を Add で更新する。
\item   u\_prev に u の結果を代入する。
\item J を返す。
            \end{enumerate}
            \paragraph{代表コード断片}
            \begin{lstlisting}[language=Python]
        def cost(u_vec):
            v = float(v0_ms)
            z = float(z0)
            Tb = float(Tb0)
            u_prev = u0
            J = 0.0
            for i in range(N):
                u = float(u_vec[i])
                if u >= 0.0:
                    tau = u * tau_drive
                else:
                    tau = u * tau_regen
                forces = self.model.resistive_forces(
                    v, env['slope_pct'], env['headwind_ms'],
                    ambient_temp_c=env.get('Tamb_C'), elevation_m=env.get('elevation_m', 0.0),
                )
                F_res = float(forces['F_total'])
                F_trac = self._traction_force(tau)
                a = (F_trac - F_res) / float(p.m)
                v = max(0.0, v + a * self.lower_dt)
                pack = self._pack_from_tau(v, tau, z, Tb, env)
                P_pack = float(pack['P_pack'])
                I = float(pack['I'])
                V = float(pack['V'])
                loss_int = float(pack['loss_int'])
                z = self.model.soc_step(
                    z,
                    P_pack,
                    self.lower_dt,
                    current_a=I,
                    Tbat_C=Tb,
                )
                Tb = Tb + (self.lower_dt / 1800.0) * (env['Tamb_C'] - Tb) + (loss_int * self.lower_dt) / 50000.0

                v_ref = float(v_ref_seq[i])
...
\end{lstlisting}

\subsection{L2273 関数 MPCNode.\_step\_solar}
            \begin{itemize}[leftmargin=1.6em]
              \item 行範囲: L2273--L2547
              \item 所属: \path|MPCNode|
              \item 定義: \path|_step_solar(self)|
              \item docstring: 明示されていない。
              \item このブロックが直接呼ぶ主な関数・メソッド: Float32, Float32MultiArray, MheInput, MheMeas, \_current\_bin\_index, \_distance\_plan\_speed, \_forecast\_at\_time, \_horizon\_data, \_interp\_upper\_speed, \_maybe\_reload\_forecast, \_measured\_distance\_km, \_mpc\_solve\_solar
              \item 戻り値の要点: 戻り値が無いか、return 文が明示されていない。
              \item 主なローカル代入名: P\_pack, current\_meas\_fresh, d0, data, distance\_meas\_km, env\_msg, env\_stop, hard\_stop, headwind\_ms, inertial\_power\_w, k\_now, limits, loss\_int, meas, meas\_now\_sec, moved\_to\_new\_bin, need\_plan, out, plan\_end\_km, s\_for\_plan
              \item 読み取る主なインスタンス属性: self.Tb, self.\_current\_bin\_index, self.\_distance\_plan\_speed, self.\_forecast\_at\_time, self.\_horizon\_data, self.\_interp\_upper\_speed, self.\_maybe\_reload\_forecast, self.\_measured\_distance\_km, self.\_mpc\_solve\_solar, self.\_mpc\_solve\_solar\_distance, self.\_publish\_metrics, self.\_publish\_plan\_path, self.\_publish\_summary, self.\_publish\_upper\_plan, self.\_route\_value, self.\_sample\_plan\_segments, self.\_soc\_guard\_speed, self.\_sync\_measured\_state, self.\_update\_live\_control\_stop\_hold, self.\_upper\_solve\_logged, self.battery\_meas\_timeout\_sec, self.control\_stop\_hold, self.current\_meas\_time, self.drive\_schedule
              \item 更新する主なインスタンス属性: self.Tb, self.\_upper\_solve\_in\_progress, self.\_upper\_solve\_logged, self.control\_stop\_hold, self.forecast\_reloaded, self.k, self.last\_bin, self.last\_data, self.last\_plan\_time, self.model\_prev\_speed\_ms, self.plan\_dt\_sec, self.plan\_start\_monotonic, self.s\_km, self.upper\_plan\_id, self.upper\_plan\_mode, self.upper\_plan\_s\_km, self.upper\_plan\_seq, self.upper\_plan\_soc, self.upper\_plan\_stop\_hold, self.upper\_plan\_tb\_c, self.upper\_plan\_time, self.v\_cmd, self.v\_plan\_kmh, self.v\_plan\_segments
              \item 明示的に送出する例外: 明示的raiseはない。
              \item 制御構造の規模: 条件分岐 52、ループ 0、try 1
            \end{itemize}
            \paragraph{この定義を読むためのPython構文}
            \begin{itemize}[leftmargin=1.6em]
            \item 第1引数selfは、このメソッドを呼んだインスタンス自身を受け取る慣習名である。
\item try文は例外が起き得る処理と、異常時または終了時の経路を分ける。
            \end{itemize}
            \paragraph{上から順の処理}
            \begin{enumerate}[leftmargin=1.8em]
            \item self.\_maybe\_reload\_forecast(...) を実行する。
\item self.\_sync\_measured\_state(...) を実行する。
\item s\_for\_stop に float(self.s\_km) の結果を代入する。
\item 条件 bool(self.get\_parameter('use\_measured\_s').value) and np.isfinite(self.s\_meas) を判定し、真なら内部処理を行う。
\item   s\_for\_stop に float(self.s\_meas) の結果を代入する。
\item v\_for\_stop に float(self.v\_upper\_cmd) の結果を代入する。
\item 条件 bool(self.get\_parameter('use\_measured\_speed').value) and np.isfinite(self.v\_now) を判定し、真なら内部処理を行う。
\item   v\_for\_stop に float(self.v\_now) の結果を代入する。
\item self.control\_stop\_hold に self.\_update\_live\_control\_stop\_hold(s\_for\_stop, v\_for\_stop) の結果を代入する。
\item k\_now に self.\_current\_bin\_index() の結果を代入する。
\item moved\_to\_new\_bin に self.last\_bin is None or k\_now != self.last\_bin の結果を代入する。
\item self.k に k\_now の結果を代入する。
\item data に self.\_horizon\_data(self.k) の結果を代入する。
\item 条件 self.control\_stop\_hold and data を判定し、真なら内部処理を行う。
\item   env\_stop に self.\_forecast\_at\_time(data[0].get('t\_utc', datetime.now(timezone.utc)), s\_for\_stop, drive=False, control\_stop=True) の結果を代入する。
\item   env\_stop['slope\_pct'] に self.\_route\_value(s\_for\_stop, 'slope\_pct', 0.0) の結果を代入する。
\item   env\_stop['t\_utc'] に data[0].get('t\_utc', datetime.now(timezone.utc)) の結果を代入する。
\item   data[0] に env\_stop の結果を代入する。
\item self.last\_data に data の結果を代入する。
\item 条件 len(data) > 0 を判定し、真なら内部処理を行う。
\item   d0 に data[0] の結果を代入する。
\item   s\_for\_profile に self.s\_km の結果を代入する。
\item   条件 bool(self.get\_parameter('use\_measured\_s').value) and np.isfinite(self.s\_meas) を判定し、真なら内部処理を行う。
\item     s\_for\_profile に float(self.s\_meas) の結果を代入する。
\item   slope\_pct に d0.get('slope\_pct', 0.0) の結果を代入する。
\item   条件 self.route\_profile is not None を判定し、真なら内部処理を行う。
\item     slope\_pct に self.\_route\_value(s\_for\_profile, 'slope\_pct', slope\_pct) の結果を代入する。
\item   headwind\_ms に d0.get('headwind\_ms', 0.0) の結果を代入する。
\item   条件 self.route\_profile is not None を判定し、真なら内部処理を行う。
\item     headwind\_ms に self.\_route\_value(s\_for\_profile, 'headwind\_ms', headwind\_ms) の結果を代入する。
\item   env\_msg に Float32MultiArray() の結果を代入する。
\item   env\_msg.data に [float(d0.get('G\_poa', 0.0)), float(d0.get('Tcell\_C', 40.0)), float(d0.get('Tamb\_C', 30.0)), float(slope\_pct), float(headwind\_ms)] の結果を代入する。
\item   self.pub\_env.publish(...) を実行する。
\item   v\_exec\_kmh に self.v\_upper\_cmd の結果を代入する。
\item   条件 self.hierarchical and self.timer\_lower is not None を判定し、真なら内部処理を行う。
\item     v\_exec\_kmh に self.v\_lower\_cmd の結果を代入する。
\item   条件 bool(self.get\_parameter('use\_measured\_speed').value) and np.isfinite(self.v\_now) を判定し、真なら内部処理を行う。
\item     v\_exec\_kmh に float(self.v\_now) の結果を代入する。
\item   self.\_publish\_metrics(...) を実行する。
\item   self.\_publish\_summary(...) を実行する。
\item need\_plan に self.v\_plan\_kmh is None or (self.forecast\_reloaded and self.replan\_on\_forecast\_reload) の結果を代入する。
\item 条件 self.upper\_mode == 'distance' and self.v\_plan\_segments is None を判定し、真なら内部処理を行う。
\item   need\_plan に True の結果を代入する。
\item 条件 self.upper\_mode == 'distance' and self.v\_plan\_segments を判定し、真なら内部処理を行う。
\item   plan\_end\_km に float(self.v\_plan\_segments[-1].get('s\_end\_km', self.s\_km)) の結果を代入する。
\item   条件 self.s\_km >= plan\_end\_km - 1e-06 を判定し、真なら内部処理を行う。
\item     need\_plan に True の結果を代入する。
\item 条件 self.upper\_replan\_km > 0.0 and self.upper\_plan\_s\_km is not None を判定し、真なら内部処理を行う。
\item   条件 self.s\_km - self.upper\_plan\_s\_km >= self.upper\_replan\_km を判定し、真なら内部処理を行う。
\item     need\_plan に True の結果を代入する。
\item 条件 self.upper\_replan\_sec > 0.0 and self.upper\_plan\_time is not None を判定し、真なら内部処理を行う。
\item   条件 data and data[0].get('t\_utc') and ((data[0]['t\_utc'] - self.upper\_plan\_time).total\_seconds() >= self.upper\_replan\_sec) を判定し、真なら内部処理を行う。
\item     need\_plan に True の結果を代入する。
\item 条件 self.upper\_replan\_soc\_delta > 0.0 and self.upper\_plan\_soc is not None を判定し、真なら内部処理を行う。
\item   条件 abs(float(self.z) - float(self.upper\_plan\_soc)) >= self.upper\_replan\_soc\_delta を判定し、真なら内部処理を行う。
\item     need\_plan に True の結果を代入する。
\item 条件 self.upper\_replan\_tb\_delta\_c > 0.0 and self.upper\_plan\_tb\_c is not None を判定し、真なら内部処理を行う。
\item   条件 abs(float(self.Tb) - float(self.upper\_plan\_tb\_c)) >= self.upper\_replan\_tb\_delta\_c を判定し、真なら内部処理を行う。
\item     need\_plan に True の結果を代入する。
\item 条件 self.upper\_replan\_on\_stop\_transition and self.upper\_plan\_stop\_hold is not None and bool(self.upper\_plan\_stop\_hold) and (not bool(self.control\_stop\_hold)) を判定し、真なら内部処理を行う。
\item   need\_plan に True の結果を代入する。
\item self.upper\_plan\_stop\_hold に bool(self.control\_stop\_hold) の結果を代入する。
\item 条件 self.upper\_replan\_km <= 0.0 and self.upper\_replan\_sec <= 0.0 を判定し、真なら内部処理を行う。
\item   need\_plan に need\_plan or moved\_to\_new\_bin の結果を代入する。
\item 条件 need\_plan and self.upper\_mode == 'distance' and self.v\_plan\_segments and (self.drive\_schedule is not None) and (len(data) > 0) and data[0].get('t\_utc') and (not self.drive\_schedule.is\_drive\_time(data[0]['t\_utc'])) を判定し、真なら内部処理を行う。
\item   plan\_end\_km に float(self.v\_plan\_segments[-1].get('s\_end\_km', self.s\_km)) の結果を代入する。
\item   条件 self.s\_km < plan\_end\_km - 1e-06 を判定し、真なら内部処理を行う。
\item     need\_plan に False の結果を代入する。
\item 条件 need\_plan and len(data) > 0 を判定し、真なら内部処理を行う。
\item   s\_for\_profile に self.s\_km の結果を代入する。
\item   条件 bool(self.get\_parameter('use\_measured\_s').value) and np.isfinite(self.s\_meas) を判定し、真なら内部処理を行う。
\item     s\_for\_profile に float(self.s\_meas) の結果を代入する。
\item   t0 に data[0].get('t\_utc', datetime.now(timezone.utc)) の結果を代入する。
\item   upper\_started に time.monotonic() の結果を代入する。
\item   self.\_upper\_solve\_in\_progress に True の結果を代入する。
\item   例外処理を伴う try ブロックを実行する。
\item     条件 self.upper\_mode == 'distance' を判定し、真なら内部処理を行う。
\item       (self.v\_upper\_cmd, self.v\_plan\_segments, self.upper\_plan\_seq) に self.\_mpc\_solve\_solar\_distance(t0, s\_for\_profile, self.upper\_plan\_seq) の結果を代入する。
\item       self.upper\_plan\_mode に 'distance' の結果を代入する。
\item       self.plan\_dt\_sec に float(self.model.p.dt) の結果を代入する。
            \end{enumerate}
            \paragraph{代表コード断片}
            \begin{lstlisting}[language=Python]
    def _step_solar(self):
        self._maybe_reload_forecast()
        self._sync_measured_state()
        s_for_stop = float(self.s_km)
        if bool(self.get_parameter('use_measured_s').value) and np.isfinite(self.s_meas):
            s_for_stop = float(self.s_meas)
        v_for_stop = float(self.v_upper_cmd)
        if bool(self.get_parameter('use_measured_speed').value) and np.isfinite(self.v_now):
            v_for_stop = float(self.v_now)
        self.control_stop_hold = self._update_live_control_stop_hold(s_for_stop, v_for_stop)
        k_now = self._current_bin_index()
        moved_to_new_bin = (self.last_bin is None) or (k_now != self.last_bin)
        self.k = k_now

        data = self._horizon_data(self.k)
        if self.control_stop_hold and data:
            env_stop = self._forecast_at_time(
                data[0].get('t_utc', datetime.now(timezone.utc)),
                s_for_stop,
                drive=False,
                control_stop=True,
            )
            env_stop['slope_pct'] = self._route_value(s_for_stop, 'slope_pct', 0.0)
            env_stop['t_utc'] = data[0].get('t_utc', datetime.now(timezone.utc))
            data[0] = env_stop
        self.last_data = data
        if len(data) > 0:
            d0 = data[0]
            s_for_profile = self.s_km
            if bool(self.get_parameter('use_measured_s').value) and np.isfinite(self.s_meas):
                s_for_profile = float(self.s_meas)
            slope_pct = d0.get('slope_pct', 0.0)
            if self.route_profile is not None:
                slope_pct = self._route_value(s_for_profile, 'slope_pct', slope_pct)
            headwind_ms = d0.get('headwind_ms', 0.0)
...
\end{lstlisting}

\subsection{L2549 関数 MPCNode.\_step\_lower}
            \begin{itemize}[leftmargin=1.6em]
              \item 行範囲: L2549--L2598
              \item 所属: \path|MPCNode|
              \item 定義: \path|_step_lower(self)|
              \item docstring: 明示されていない。
              \item このブロックが直接呼ぶ主な関数・メソッド: \_build\_lower\_ref, \_lower\_mpc\_solve, \_publish\_lower\_command\_cycle, \_route\_value, bool, dict, float, get, get\_logger, get\_parameter, info, isfinite
              \item 戻り値の要点: 戻り値が無いか、return 文が明示されていない。
              \item 主なローカル代入名: base\_time, d0, env, headwind\_ms, mode, s\_for\_profile, slope\_pct, solve\_started, u\_seq, v0\_ms, v\_cmd\_ms, v\_pred, v\_ref\_seq
              \item 読み取る主なインスタンス属性: self.Tb, self.\_build\_lower\_ref, self.\_lower\_mpc\_solve, self.\_lower\_solve\_logged, self.\_publish\_lower\_command\_cycle, self.\_route\_value, self.get\_logger, self.get\_parameter, self.hierarchical, self.last\_data, self.route\_profile, self.s\_km, self.s\_meas, self.timer\_lower, self.v\_lower\_cmd, self.v\_now, self.v\_upper\_cmd, self.z
              \item 更新する主なインスタンス属性: self.\_lower\_solve\_logged, self.lower\_last\_mode, self.lower\_last\_u, self.lower\_last\_v\_pred, self.v\_cmd, self.v\_lower\_cmd
              \item 明示的に送出する例外: 明示的raiseはない。
              \item 制御構造の規模: 条件分岐 10、ループ 0、try 0
            \end{itemize}
            \paragraph{この定義を読むためのPython構文}
            \begin{itemize}[leftmargin=1.6em]
            \item 第1引数selfは、このメソッドを呼んだインスタンス自身を受け取る慣習名である。
            \end{itemize}
            \paragraph{上から順の処理}
            \begin{enumerate}[leftmargin=1.8em]
            \item 条件 not self.hierarchical or self.timer\_lower is None を判定し、真なら内部処理を行う。
\item    を返す。
\item self.\_publish\_lower\_command\_cycle(...) を実行する。
\item 条件 not self.last\_data を判定し、真なら内部処理を行う。
\item    を返す。
\item d0 に self.last\_data[0] の結果を代入する。
\item s\_for\_profile に self.s\_km の結果を代入する。
\item 条件 bool(self.get\_parameter('use\_measured\_s').value) and np.isfinite(self.s\_meas) を判定し、真なら内部処理を行う。
\item   s\_for\_profile に float(self.s\_meas) の結果を代入する。
\item slope\_pct に d0.get('slope\_pct', 0.0) の結果を代入する。
\item 条件 self.route\_profile is not None を判定し、真なら内部処理を行う。
\item   slope\_pct に self.\_route\_value(s\_for\_profile, 'slope\_pct', slope\_pct) の結果を代入する。
\item headwind\_ms に d0.get('headwind\_ms', 0.0) の結果を代入する。
\item 条件 self.route\_profile is not None を判定し、真なら内部処理を行う。
\item   headwind\_ms に self.\_route\_value(s\_for\_profile, 'headwind\_ms', headwind\_ms) の結果を代入する。
\item env に dict(slope\_pct=float(slope\_pct), headwind\_ms=float(headwind\_ms), G\_poa=float(d0.get('G\_poa', 0.0)), Tcell\_C=float(d0.get('Tcell\_C', 40.0)), Tamb\_C=float(d0.get('Tamb\_C', 30.0))) の結果を代入する。
\item base\_time に d0.get('t\_utc', None) の結果を代入する。
\item 条件 base\_time is None を判定し、真なら内部処理を行う。
\item   base\_time に datetime.now(timezone.utc) の結果を代入する。
\item v\_ref\_seq に self.\_build\_lower\_ref(base\_time, s\_for\_profile, d0) の結果を代入する。
\item 条件 bool(self.get\_parameter('use\_measured\_speed').value) and np.isfinite(self.v\_now) を判定し、真なら内部処理を行う。
\item   v0\_ms に float(self.v\_now) / 3.6 の結果を代入する。
\item   上の条件が偽の場合:
\item   v0\_ms に float(self.v\_lower\_cmd) / 3.6 の結果を代入する。
\item   条件 not np.isfinite(v0\_ms) or v0\_ms <= 0.0 を判定し、真なら内部処理を行う。
\item     v0\_ms に float(self.v\_upper\_cmd) / 3.6 の結果を代入する。
\item solve\_started に time.monotonic() の結果を代入する。
\item (u\_seq, v\_pred, mode) に self.\_lower\_mpc\_solve(v0\_ms, s\_for\_profile, self.z, self.Tb, env, v\_ref\_seq) の結果を代入する。
\item 条件 not self.\_lower\_solve\_logged を判定し、真なら内部処理を行う。
\item   self.get\_logger().info(...) を実行する。
\item   self.\_lower\_solve\_logged に True の結果を代入する。
\item v\_cmd\_ms に v\_pred[0] if v\_pred else v0\_ms の結果を代入する。
\item self.v\_lower\_cmd に float(v\_cmd\_ms * 3.6) の結果を代入する。
\item self.v\_cmd に float(self.v\_lower\_cmd) の結果を代入する。
\item 条件 len(u\_seq) > 0 を判定し、真なら内部処理を行う。
\item   self.lower\_last\_u に float(u\_seq[0]) の結果を代入する。
\item self.lower\_last\_mode に str(mode) の結果を代入する。
\item self.lower\_last\_v\_pred に list(v\_pred) if v\_pred else [self.v\_lower\_cmd / 3.6] の結果を代入する。
\item self.\_publish\_lower\_command\_cycle(...) を実行する。
            \end{enumerate}
            \paragraph{代表コード断片}
            \begin{lstlisting}[language=Python]
    def _step_lower(self):
        if not self.hierarchical or self.timer_lower is None:
            return
        self._publish_lower_command_cycle()
        if not self.last_data:
            return
        d0 = self.last_data[0]
        s_for_profile = self.s_km
        if bool(self.get_parameter('use_measured_s').value) and np.isfinite(self.s_meas):
            s_for_profile = float(self.s_meas)
        slope_pct = d0.get('slope_pct', 0.0)
        if self.route_profile is not None:
            slope_pct = self._route_value(s_for_profile, 'slope_pct', slope_pct)
        headwind_ms = d0.get('headwind_ms', 0.0)
        if self.route_profile is not None:
            headwind_ms = self._route_value(s_for_profile, 'headwind_ms', headwind_ms)
        env = dict(
            slope_pct=float(slope_pct),
            headwind_ms=float(headwind_ms),
            G_poa=float(d0.get('G_poa', 0.0)),
            Tcell_C=float(d0.get('Tcell_C', 40.0)),
            Tamb_C=float(d0.get('Tamb_C', 30.0)),
        )
        base_time = d0.get('t_utc', None)
        if base_time is None:
            base_time = datetime.now(timezone.utc)
        v_ref_seq = self._build_lower_ref(base_time, s_for_profile, d0)

        if bool(self.get_parameter('use_measured_speed').value) and np.isfinite(self.v_now):
            v0_ms = float(self.v_now) / 3.6
        else:
            v0_ms = float(self.v_lower_cmd) / 3.6
            if not np.isfinite(v0_ms) or v0_ms <= 0.0:
                v0_ms = float(self.v_upper_cmd) / 3.6

...
\end{lstlisting}

\subsection{L2600 関数 MPCNode.\_publish\_lower\_command\_cycle}
            \begin{itemize}[leftmargin=1.6em]
              \item 行範囲: L2600--L2615
              \item 所属: \path|MPCNode|
              \item 定義: \path|_publish_lower_command_cycle(self) -> None|
              \item docstring: Non-blocking 1 Hz output path, independent of both optimizers.
              \item このブロックが直接呼ぶ主な関数・メソッド: Float32, String, \_publish\_lower\_plan, \_publish\_summary, \_sync\_measured\_state, clip, float, get\_parameter, min, publish, str
              \item 戻り値の要点: 戻り値が無いか、return 文が明示されていない。
              \item 主なローカル代入名: speed\_kmh, throttle\_pct
              \item 読み取る主なインスタンス属性: self.\_publish\_lower\_plan, self.\_publish\_summary, self.\_sync\_measured\_state, self.control\_stop\_hold, self.get\_parameter, self.hierarchical, self.lower\_last\_mode, self.lower\_last\_u, self.lower\_last\_v\_pred, self.pub\_drive\_mode, self.pub\_speed, self.pub\_throttle, self.timer\_lower, self.v\_lower\_cmd
              \item 更新する主なインスタンス属性: 特になし。
              \item 明示的に送出する例外: 明示的raiseはない。
              \item 制御構造の規模: 条件分岐 3、ループ 0、try 0
            \end{itemize}
            \paragraph{この定義を読むためのPython構文}
            \begin{itemize}[leftmargin=1.6em]
            \item 第1引数selfは、このメソッドを呼んだインスタンス自身を受け取る慣習名である。
\item `->`以後は戻り値の型注釈であり、実行時の値を自動変換する命令ではない。
            \end{itemize}
            \paragraph{上から順の処理}
            \begin{enumerate}[leftmargin=1.8em]
            \item 条件 not self.hierarchical or self.timer\_lower is None を判定し、真なら内部処理を行う。
\item    を返す。
\item self.\_sync\_measured\_state(...) を実行する。
\item speed\_kmh に float(np.clip(self.v\_lower\_cmd, 0.0, self.get\_parameter('v\_max\_kmh').value)) の結果を代入する。
\item 条件 self.control\_stop\_hold を判定し、真なら内部処理を行う。
\item   speed\_kmh に 0.0 の結果を代入する。
\item self.pub\_speed.publish(...) を実行する。
\item throttle\_pct に float(np.clip(self.lower\_last\_u * 100.0, -100.0, 100.0)) の結果を代入する。
\item 条件 self.control\_stop\_hold を判定し、真なら内部処理を行う。
\item   throttle\_pct に min(0.0, throttle\_pct) の結果を代入する。
\item self.pub\_throttle.publish(...) を実行する。
\item self.pub\_drive\_mode.publish(...) を実行する。
\item self.\_publish\_lower\_plan(...) を実行する。
\item self.\_publish\_summary(...) を実行する。
            \end{enumerate}
            \paragraph{代表コード断片}
            \begin{lstlisting}[language=Python]
    def _publish_lower_command_cycle(self) -> None:
        """Non-blocking 1 Hz output path, independent of both optimizers."""
        if not self.hierarchical or self.timer_lower is None:
            return
        self._sync_measured_state()
        speed_kmh = float(np.clip(self.v_lower_cmd, 0.0, self.get_parameter('v_max_kmh').value))
        if self.control_stop_hold:
            speed_kmh = 0.0
        self.pub_speed.publish(Float32(data=speed_kmh))
        throttle_pct = float(np.clip(self.lower_last_u * 100.0, -100.0, 100.0))
        if self.control_stop_hold:
            throttle_pct = min(0.0, throttle_pct)
        self.pub_throttle.publish(Float32(data=throttle_pct))
        self.pub_drive_mode.publish(String(data=str(self.lower_last_mode)))
        self._publish_lower_plan(self.lower_last_v_pred)
        self._publish_summary()
\end{lstlisting}

\subsection{L2618 関数 MPCNode.\_init\_passo}
            \begin{itemize}[leftmargin=1.6em]
              \item 行範囲: L2618--L2707
              \item 所属: \path|MPCNode|
              \item 定義: \path|_init_passo(self)|
              \item docstring: 明示されていない。
              \item このブロックが直接呼ぶ主な関数・メソッド: \_load\_stops, array, bool, create\_publisher, create\_subscription, create\_timer, declare\_parameter, deque, float, get\_logger, get\_parameter, info
              \item 戻り値の要点: 戻り値が無いか、return 文が明示されていない。
              \item 主なローカル代入名: stop\_yaml
              \item 読み取る主なインスタンス属性: self.\_load\_stops, self.\_on\_config, self.\_on\_config\_ready, self.\_on\_fuel, self.\_on\_grade, self.\_on\_idle\_fuel, self.\_on\_obd\_ok, self.\_on\_s\_km, self.\_on\_speed, self.\_on\_system\_state, self.\_on\_throttle, self.\_step\_passo, self.create\_publisher, self.create\_subscription, self.create\_timer, self.declare\_parameter, self.get\_logger, self.get\_parameter
              \item 更新する主なインスタンス属性: self.Np, self.config\_ready, self.dt, self.fuel\_rate\_lph, self.grade, self.id\_ema\_alpha, self.id\_min\_samples, self.id\_r2, self.id\_rmse, self.id\_samples, self.id\_window\_sec, self.idle\_fuel\_lph, self.last\_step\_time, self.max\_acc\_dt\_sec, self.model\_coeffs, self.mpc\_state, self.obd\_ok, self.online\_id\_enabled, self.prev\_speed\_kmh, self.prev\_speed\_time, self.pub\_mpc\_state, self.pub\_path, self.pub\_speed, self.pub\_status
              \item 明示的に送出する例外: 明示的raiseはない。
              \item 制御構造の規模: 条件分岐 0、ループ 0、try 0
            \end{itemize}
            \paragraph{この定義を読むためのPython構文}
            \begin{itemize}[leftmargin=1.6em]
            \item 第1引数selfは、このメソッドを呼んだインスタンス自身を受け取る慣習名である。
            \end{itemize}
            \paragraph{上から順の処理}
            \begin{enumerate}[leftmargin=1.8em]
            \item self.declare\_parameter(...) を実行する。
\item self.declare\_parameter(...) を実行する。
\item self.declare\_parameter(...) を実行する。
\item self.declare\_parameter(...) を実行する。
\item self.declare\_parameter(...) を実行する。
\item self.declare\_parameter(...) を実行する。
\item self.declare\_parameter(...) を実行する。
\item self.declare\_parameter(...) を実行する。
\item self.declare\_parameter(...) を実行する。
\item self.declare\_parameter(...) を実行する。
\item self.declare\_parameter(...) を実行する。
\item self.declare\_parameter(...) を実行する。
\item self.declare\_parameter(...) を実行する。
\item self.declare\_parameter(...) を実行する。
\item self.declare\_parameter(...) を実行する。
\item self.declare\_parameter(...) を実行する。
\item self.declare\_parameter(...) を実行する。
\item self.declare\_parameter(...) を実行する。
\item self.declare\_parameter(...) を実行する。
\item self.declare\_parameter(...) を実行する。
\item self.declare\_parameter(...) を実行する。
\item self.declare\_parameter(...) を実行する。
\item self.declare\_parameter(...) を実行する。
\item self.dt に float(self.get\_parameter('passo\_dt').value) の結果を代入する。
\item self.Np に int(self.get\_parameter('passo\_horizon\_steps').value) の結果を代入する。
\item self.v\_cmd に float(self.get\_parameter('v\_ref\_kmh').value) の結果を代入する。
\item stop\_yaml に self.get\_parameter('stop\_yaml').value の結果を代入する。
\item self.\_load\_stops(...) を実行する。
\item self.v\_now に math.nan の結果を代入する。
\item self.s\_km に 0.0 の結果を代入する。
\item self.fuel\_rate\_lph に math.nan の結果を代入する。
\item self.throttle\_pct に math.nan の結果を代入する。
\item self.obd\_ok に 0.0 の結果を代入する。
\item self.config\_ready に False の結果を代入する。
\item self.mpc\_state に 'IDLE' の結果を代入する。
\item self.system\_state に '' の結果を代入する。
\item self.grade に math.nan の結果を代入する。
\item self.idle\_fuel\_lph に math.nan の結果を代入する。
\item self.model\_coeffs に np.array([float(self.get\_parameter('model\_a0').value), float(self.get\_parameter('model\_a1').value), float(self.get\_parameter('model\_a2').value), float(self.get\_parameter('model\_a3').value), float(self.get\_parameter('model\_a4').value)], dtype=float) の結果を代入する。
\item self.online\_id\_enabled に bool(self.get\_parameter('online\_id\_enabled').value) の結果を代入する。
\item self.id\_window\_sec に float(self.get\_parameter('id\_window\_sec').value) の結果を代入する。
\item self.id\_min\_samples に int(self.get\_parameter('id\_min\_samples').value) の結果を代入する。
\item self.id\_ema\_alpha に float(self.get\_parameter('id\_ema\_alpha').value) の結果を代入する。
\item self.max\_acc\_dt\_sec に float(self.get\_parameter('max\_acc\_dt\_sec').value) の結果を代入する。
\item self.id\_samples に deque() の結果を代入する。
\item self.id\_rmse に math.nan の結果を代入する。
\item self.id\_r2 に math.nan の結果を代入する。
\item self.prev\_speed\_kmh に math.nan の結果を代入する。
\item self.prev\_speed\_time に None の結果を代入する。
\item self.valid\_obd\_sec に 0.0 の結果を代入する。
\item self.valid\_speed\_sec に 0.0 の結果を代入する。
\item self.valid\_fuel\_sec に 0.0 の結果を代入する。
\item self.last\_step\_time に None の結果を代入する。
\item self.create\_subscription(...) を実行する。
\item self.create\_subscription(...) を実行する。
\item self.create\_subscription(...) を実行する。
\item self.create\_subscription(...) を実行する。
\item self.create\_subscription(...) を実行する。
\item self.create\_subscription(...) を実行する。
\item self.create\_subscription(...) を実行する。
\item self.create\_subscription(...) を実行する。
\item self.create\_subscription(...) を実行する。
\item self.create\_subscription(...) を実行する。
\item self.pub\_speed に self.create\_publisher(Float32, '/planner/speed\_cmd', 10) の結果を代入する。
\item self.pub\_path に self.create\_publisher(Path, '/planner/trajectory', 10) の結果を代入する。
\item self.pub\_status に self.create\_publisher(Float32MultiArray, '/planner/status', 10) の結果を代入する。
\item self.pub\_mpc\_state に self.create\_publisher(String, '/system/mpc\_state', 10) の結果を代入する。
\item self.timer に self.create\_timer(1.0, self.\_step\_passo) の結果を代入する。
\item self.get\_logger().info(...) を実行する。
            \end{enumerate}
            \paragraph{代表コード断片}
            \begin{lstlisting}[language=Python]
    def _init_passo(self):
        self.declare_parameter('stop_yaml', 'inputs/stop_points.yaml')
        self.declare_parameter('passo_dt', 1.0)
        self.declare_parameter('passo_horizon_steps', 10)
        self.declare_parameter('v_min_kmh', 0.0)
        self.declare_parameter('v_max_kmh', 110.0)
        self.declare_parameter('v_ref_kmh', 40.0)
        self.declare_parameter('w_fuel', 1.0)
        self.declare_parameter('w_speed', 0.3)
        self.declare_parameter('w_dv', 0.2)
        self.declare_parameter('w_dv_limit', 2.0)
        self.declare_parameter('dv_max_kmhps', 4.0)
        self.declare_parameter('w_stop', 1.0e4)
        self.declare_parameter('model_a0', 0.4)
        self.declare_parameter('model_a1', 0.02)
        self.declare_parameter('model_a2', 0.001)
        self.declare_parameter('model_a3', 0.08)
        self.declare_parameter('model_a4', 0.02)
        self.declare_parameter('online_id_enabled', True)
        self.declare_parameter('id_window_sec', 60.0)
        self.declare_parameter('id_min_samples', 30)
        self.declare_parameter('id_ema_alpha', 0.2)
        self.declare_parameter('max_acc_dt_sec', 2.0)
        self.declare_parameter('run_ready_sec', 3.0)

        self.dt = float(self.get_parameter('passo_dt').value)
        self.Np = int(self.get_parameter('passo_horizon_steps').value)
        self.v_cmd = float(self.get_parameter('v_ref_kmh').value)

        stop_yaml = self.get_parameter('stop_yaml').value
        self._load_stops(stop_yaml)

        # Inputs
        self.v_now = math.nan
        self.s_km = 0.0
...
\end{lstlisting}

\subsection{L2709 関数 MPCNode.\_on\_s\_km}
            \begin{itemize}[leftmargin=1.6em]
              \item 行範囲: L2709--L2710
              \item 所属: \path|MPCNode|
              \item 定義: \path|_on_s_km(self, msg: Float32)|
              \item docstring: 明示されていない。
              \item このブロックが直接呼ぶ主な関数・メソッド: float
              \item 戻り値の要点: 戻り値が無いか、return 文が明示されていない。
              \item 主なローカル代入名: 特になし。
              \item 読み取る主なインスタンス属性: 特になし。
              \item 更新する主なインスタンス属性: self.s\_km
              \item 明示的に送出する例外: 明示的raiseはない。
              \item 制御構造の規模: 条件分岐 0、ループ 0、try 0
            \end{itemize}
            \paragraph{この定義を読むためのPython構文}
            \begin{itemize}[leftmargin=1.6em]
            \item 第1引数selfは、このメソッドを呼んだインスタンス自身を受け取る慣習名である。
            \end{itemize}
            \paragraph{上から順の処理}
            \begin{enumerate}[leftmargin=1.8em]
            \item self.s\_km に float(msg.data) の結果を代入する。
            \end{enumerate}
            \paragraph{代表コード断片}
            \begin{lstlisting}[language=Python]
    def _on_s_km(self, msg: Float32):
        self.s_km = float(msg.data)
\end{lstlisting}

\subsection{L2712 関数 MPCNode.\_on\_speed}
            \begin{itemize}[leftmargin=1.6em]
              \item 行範囲: L2712--L2713
              \item 所属: \path|MPCNode|
              \item 定義: \path|_on_speed(self, msg: Float32)|
              \item docstring: 明示されていない。
              \item このブロックが直接呼ぶ主な関数・メソッド: float
              \item 戻り値の要点: 戻り値が無いか、return 文が明示されていない。
              \item 主なローカル代入名: 特になし。
              \item 読み取る主なインスタンス属性: 特になし。
              \item 更新する主なインスタンス属性: self.v\_now
              \item 明示的に送出する例外: 明示的raiseはない。
              \item 制御構造の規模: 条件分岐 0、ループ 0、try 0
            \end{itemize}
            \paragraph{この定義を読むためのPython構文}
            \begin{itemize}[leftmargin=1.6em]
            \item 第1引数selfは、このメソッドを呼んだインスタンス自身を受け取る慣習名である。
            \end{itemize}
            \paragraph{上から順の処理}
            \begin{enumerate}[leftmargin=1.8em]
            \item self.v\_now に float(msg.data) の結果を代入する。
            \end{enumerate}
            \paragraph{代表コード断片}
            \begin{lstlisting}[language=Python]
    def _on_speed(self, msg: Float32):
        self.v_now = float(msg.data)
\end{lstlisting}

\subsection{L2715 関数 MPCNode.\_on\_fuel}
            \begin{itemize}[leftmargin=1.6em]
              \item 行範囲: L2715--L2716
              \item 所属: \path|MPCNode|
              \item 定義: \path|_on_fuel(self, msg: Float32)|
              \item docstring: 明示されていない。
              \item このブロックが直接呼ぶ主な関数・メソッド: float
              \item 戻り値の要点: 戻り値が無いか、return 文が明示されていない。
              \item 主なローカル代入名: 特になし。
              \item 読み取る主なインスタンス属性: 特になし。
              \item 更新する主なインスタンス属性: self.fuel\_rate\_lph
              \item 明示的に送出する例外: 明示的raiseはない。
              \item 制御構造の規模: 条件分岐 0、ループ 0、try 0
            \end{itemize}
            \paragraph{この定義を読むためのPython構文}
            \begin{itemize}[leftmargin=1.6em]
            \item 第1引数selfは、このメソッドを呼んだインスタンス自身を受け取る慣習名である。
            \end{itemize}
            \paragraph{上から順の処理}
            \begin{enumerate}[leftmargin=1.8em]
            \item self.fuel\_rate\_lph に float(msg.data) の結果を代入する。
            \end{enumerate}
            \paragraph{代表コード断片}
            \begin{lstlisting}[language=Python]
    def _on_fuel(self, msg: Float32):
        self.fuel_rate_lph = float(msg.data)
\end{lstlisting}

\subsection{L2718 関数 MPCNode.\_on\_throttle}
            \begin{itemize}[leftmargin=1.6em]
              \item 行範囲: L2718--L2719
              \item 所属: \path|MPCNode|
              \item 定義: \path|_on_throttle(self, msg: Float32)|
              \item docstring: 明示されていない。
              \item このブロックが直接呼ぶ主な関数・メソッド: float
              \item 戻り値の要点: 戻り値が無いか、return 文が明示されていない。
              \item 主なローカル代入名: 特になし。
              \item 読み取る主なインスタンス属性: 特になし。
              \item 更新する主なインスタンス属性: self.throttle\_pct
              \item 明示的に送出する例外: 明示的raiseはない。
              \item 制御構造の規模: 条件分岐 0、ループ 0、try 0
            \end{itemize}
            \paragraph{この定義を読むためのPython構文}
            \begin{itemize}[leftmargin=1.6em]
            \item 第1引数selfは、このメソッドを呼んだインスタンス自身を受け取る慣習名である。
            \end{itemize}
            \paragraph{上から順の処理}
            \begin{enumerate}[leftmargin=1.8em]
            \item self.throttle\_pct に float(msg.data) の結果を代入する。
            \end{enumerate}
            \paragraph{代表コード断片}
            \begin{lstlisting}[language=Python]
    def _on_throttle(self, msg: Float32):
        self.throttle_pct = float(msg.data)
\end{lstlisting}

\subsection{L2721 関数 MPCNode.\_on\_obd\_ok}
            \begin{itemize}[leftmargin=1.6em]
              \item 行範囲: L2721--L2722
              \item 所属: \path|MPCNode|
              \item 定義: \path|_on_obd_ok(self, msg: Float32)|
              \item docstring: 明示されていない。
              \item このブロックが直接呼ぶ主な関数・メソッド: float
              \item 戻り値の要点: 戻り値が無いか、return 文が明示されていない。
              \item 主なローカル代入名: 特になし。
              \item 読み取る主なインスタンス属性: 特になし。
              \item 更新する主なインスタンス属性: self.obd\_ok
              \item 明示的に送出する例外: 明示的raiseはない。
              \item 制御構造の規模: 条件分岐 0、ループ 0、try 0
            \end{itemize}
            \paragraph{この定義を読むためのPython構文}
            \begin{itemize}[leftmargin=1.6em]
            \item 第1引数selfは、このメソッドを呼んだインスタンス自身を受け取る慣習名である。
            \end{itemize}
            \paragraph{上から順の処理}
            \begin{enumerate}[leftmargin=1.8em]
            \item self.obd\_ok に float(msg.data) の結果を代入する。
            \end{enumerate}
            \paragraph{代表コード断片}
            \begin{lstlisting}[language=Python]
    def _on_obd_ok(self, msg: Float32):
        self.obd_ok = float(msg.data)
\end{lstlisting}

\subsection{L2724 関数 MPCNode.\_on\_grade}
            \begin{itemize}[leftmargin=1.6em]
              \item 行範囲: L2724--L2725
              \item 所属: \path|MPCNode|
              \item 定義: \path|_on_grade(self, msg: Float32)|
              \item docstring: 明示されていない。
              \item このブロックが直接呼ぶ主な関数・メソッド: float
              \item 戻り値の要点: 戻り値が無いか、return 文が明示されていない。
              \item 主なローカル代入名: 特になし。
              \item 読み取る主なインスタンス属性: 特になし。
              \item 更新する主なインスタンス属性: self.grade
              \item 明示的に送出する例外: 明示的raiseはない。
              \item 制御構造の規模: 条件分岐 0、ループ 0、try 0
            \end{itemize}
            \paragraph{この定義を読むためのPython構文}
            \begin{itemize}[leftmargin=1.6em]
            \item 第1引数selfは、このメソッドを呼んだインスタンス自身を受け取る慣習名である。
            \end{itemize}
            \paragraph{上から順の処理}
            \begin{enumerate}[leftmargin=1.8em]
            \item self.grade に float(msg.data) の結果を代入する。
            \end{enumerate}
            \paragraph{代表コード断片}
            \begin{lstlisting}[language=Python]
    def _on_grade(self, msg: Float32):
        self.grade = float(msg.data)
\end{lstlisting}

\subsection{L2727 関数 MPCNode.\_on\_idle\_fuel}
            \begin{itemize}[leftmargin=1.6em]
              \item 行範囲: L2727--L2730
              \item 所属: \path|MPCNode|
              \item 定義: \path|_on_idle_fuel(self, msg: Float32)|
              \item docstring: 明示されていない。
              \item このブロックが直接呼ぶ主な関数・メソッド: float, isfinite
              \item 戻り値の要点: 戻り値が無いか、return 文が明示されていない。
              \item 主なローカル代入名: 特になし。
              \item 読み取る主なインスタンス属性: self.idle\_fuel\_lph, self.model\_coeffs
              \item 更新する主なインスタンス属性: self.idle\_fuel\_lph
              \item 明示的に送出する例外: 明示的raiseはない。
              \item 制御構造の規模: 条件分岐 1、ループ 0、try 0
            \end{itemize}
            \paragraph{この定義を読むためのPython構文}
            \begin{itemize}[leftmargin=1.6em]
            \item 第1引数selfは、このメソッドを呼んだインスタンス自身を受け取る慣習名である。
            \end{itemize}
            \paragraph{上から順の処理}
            \begin{enumerate}[leftmargin=1.8em]
            \item self.idle\_fuel\_lph に float(msg.data) の結果を代入する。
\item 条件 np.isfinite(self.idle\_fuel\_lph) を判定し、真なら内部処理を行う。
\item   self.model\_coeffs[0] に float(self.idle\_fuel\_lph) の結果を代入する。
            \end{enumerate}
            \paragraph{代表コード断片}
            \begin{lstlisting}[language=Python]
    def _on_idle_fuel(self, msg: Float32):
        self.idle_fuel_lph = float(msg.data)
        if np.isfinite(self.idle_fuel_lph):
            self.model_coeffs[0] = float(self.idle_fuel_lph)
\end{lstlisting}

\subsection{L2732 関数 MPCNode.\_on\_config\_ready}
            \begin{itemize}[leftmargin=1.6em]
              \item 行範囲: L2732--L2733
              \item 所属: \path|MPCNode|
              \item 定義: \path|_on_config_ready(self, msg: Bool)|
              \item docstring: 明示されていない。
              \item このブロックが直接呼ぶ主な関数・メソッド: bool
              \item 戻り値の要点: 戻り値が無いか、return 文が明示されていない。
              \item 主なローカル代入名: 特になし。
              \item 読み取る主なインスタンス属性: 特になし。
              \item 更新する主なインスタンス属性: self.config\_ready
              \item 明示的に送出する例外: 明示的raiseはない。
              \item 制御構造の規模: 条件分岐 0、ループ 0、try 0
            \end{itemize}
            \paragraph{この定義を読むためのPython構文}
            \begin{itemize}[leftmargin=1.6em]
            \item 第1引数selfは、このメソッドを呼んだインスタンス自身を受け取る慣習名である。
            \end{itemize}
            \paragraph{上から順の処理}
            \begin{enumerate}[leftmargin=1.8em]
            \item self.config\_ready に bool(msg.data) の結果を代入する。
            \end{enumerate}
            \paragraph{代表コード断片}
            \begin{lstlisting}[language=Python]
    def _on_config_ready(self, msg: Bool):
        self.config_ready = bool(msg.data)
\end{lstlisting}

\subsection{L2735 関数 MPCNode.\_on\_system\_state}
            \begin{itemize}[leftmargin=1.6em]
              \item 行範囲: L2735--L2736
              \item 所属: \path|MPCNode|
              \item 定義: \path|_on_system_state(self, msg: String)|
              \item docstring: 明示されていない。
              \item このブロックが直接呼ぶ主な関数・メソッド: str
              \item 戻り値の要点: 戻り値が無いか、return 文が明示されていない。
              \item 主なローカル代入名: 特になし。
              \item 読み取る主なインスタンス属性: 特になし。
              \item 更新する主なインスタンス属性: self.system\_state
              \item 明示的に送出する例外: 明示的raiseはない。
              \item 制御構造の規模: 条件分岐 0、ループ 0、try 0
            \end{itemize}
            \paragraph{この定義を読むためのPython構文}
            \begin{itemize}[leftmargin=1.6em]
            \item 第1引数selfは、このメソッドを呼んだインスタンス自身を受け取る慣習名である。
            \end{itemize}
            \paragraph{上から順の処理}
            \begin{enumerate}[leftmargin=1.8em]
            \item self.system\_state に str(msg.data) の結果を代入する。
            \end{enumerate}
            \paragraph{代表コード断片}
            \begin{lstlisting}[language=Python]
    def _on_system_state(self, msg: String):
        self.system_state = str(msg.data)
\end{lstlisting}

\subsection{L2738 関数 MPCNode.\_on\_config}
            \begin{itemize}[leftmargin=1.6em]
              \item 行範囲: L2738--L2743
              \item 所属: \path|MPCNode|
              \item 定義: \path|_on_config(self, msg: String)|
              \item docstring: 明示されていない。
              \item このブロックが直接呼ぶ主な関数・メソッド: \_apply\_config, safe\_load
              \item 戻り値の要点: 戻り値が無いか、return 文が明示されていない。
              \item 主なローカル代入名: cfg
              \item 読み取る主なインスタンス属性: self.\_apply\_config
              \item 更新する主なインスタンス属性: 特になし。
              \item 明示的に送出する例外: 明示的raiseはない。
              \item 制御構造の規模: 条件分岐 0、ループ 0、try 1
            \end{itemize}
            \paragraph{この定義を読むためのPython構文}
            \begin{itemize}[leftmargin=1.6em]
            \item 第1引数selfは、このメソッドを呼んだインスタンス自身を受け取る慣習名である。
\item try文は例外が起き得る処理と、異常時または終了時の経路を分ける。
            \end{itemize}
            \paragraph{上から順の処理}
            \begin{enumerate}[leftmargin=1.8em]
            \item 例外処理を伴う try ブロックを実行する。
\item   cfg に yaml.safe\_load(msg.data) or \{\} の結果を代入する。
\item   Exceptionを捕捉した場合:
\item    を返す。
\item self.\_apply\_config(...) を実行する。
            \end{enumerate}
            \paragraph{代表コード断片}
            \begin{lstlisting}[language=Python]
    def _on_config(self, msg: String):
        try:
            cfg = yaml.safe_load(msg.data) or {}
        except Exception:
            return
        self._apply_config(cfg)
\end{lstlisting}

\subsection{L2745 関数 MPCNode.\_apply\_config}
            \begin{itemize}[leftmargin=1.6em]
              \item 行範囲: L2745--L2776
              \item 所属: \path|MPCNode|
              \item 定義: \path|_apply_config(self, cfg: dict)|
              \item docstring: 明示されていない。
              \item このブロックが直接呼ぶ主な関数・メソッド: Parameter, \_load\_stops, append, bool, float, int, set\_parameters, str
              \item 戻り値の要点: 戻り値が無いか、return 文が明示されていない。
              \item 主なローカル代入名: key, params
              \item 読み取る主なインスタンス属性: self.\_load\_stops, self.model\_coeffs, self.set\_parameters
              \item 更新する主なインスタンス属性: self.Np, self.dt, self.online\_id\_enabled
              \item 明示的に送出する例外: 明示的raiseはない。
              \item 制御構造の規模: 条件分岐 14、ループ 1、try 0
            \end{itemize}
            \paragraph{この定義を読むためのPython構文}
            \begin{itemize}[leftmargin=1.6em]
            \item 第1引数selfは、このメソッドを呼んだインスタンス自身を受け取る慣習名である。
            \end{itemize}
            \paragraph{上から順の処理}
            \begin{enumerate}[leftmargin=1.8em]
            \item params に [] の結果を代入する。
\item ['v\_min\_kmh', 'v\_max\_kmh', 'v\_ref\_kmh', 'w\_fuel', 'w\_speed', 'w\_dv', 'w\_stop', 'dv\_max\_kmhps'] を順に走査し、各要素を key に入れて処理する。
\item   条件 key in cfg を判定し、真なら内部処理を行う。
\item     params.append(...) を実行する。
\item 条件 'dv\_max\_kmh\_per\_s' in cfg and 'dv\_max\_kmhps' not in cfg を判定し、真なら内部処理を行う。
\item   params.append(...) を実行する。
\item 条件 'horizon\_steps' in cfg を判定し、真なら内部処理を行う。
\item   params.append(...) を実行する。
\item 条件 'dt\_control' in cfg を判定し、真なら内部処理を行う。
\item   params.append(...) を実行する。
\item 条件 params を判定し、真なら内部処理を行う。
\item   self.set\_parameters(...) を実行する。
\item 条件 'horizon\_steps' in cfg を判定し、真なら内部処理を行う。
\item   self.Np に int(cfg['horizon\_steps']) の結果を代入する。
\item 条件 'dt\_control' in cfg を判定し、真なら内部処理を行う。
\item   self.dt に float(cfg['dt\_control']) の結果を代入する。
\item 条件 'model\_a0' in cfg を判定し、真なら内部処理を行う。
\item   self.model\_coeffs[0] に float(cfg['model\_a0']) の結果を代入する。
\item 条件 'model\_a1' in cfg を判定し、真なら内部処理を行う。
\item   self.model\_coeffs[1] に float(cfg['model\_a1']) の結果を代入する。
\item 条件 'model\_a2' in cfg を判定し、真なら内部処理を行う。
\item   self.model\_coeffs[2] に float(cfg['model\_a2']) の結果を代入する。
\item 条件 'model\_a3' in cfg を判定し、真なら内部処理を行う。
\item   self.model\_coeffs[3] に float(cfg['model\_a3']) の結果を代入する。
\item 条件 'model\_a4' in cfg を判定し、真なら内部処理を行う。
\item   self.model\_coeffs[4] に float(cfg['model\_a4']) の結果を代入する。
\item 条件 'online\_id\_enabled' in cfg を判定し、真なら内部処理を行う。
\item   self.online\_id\_enabled に bool(cfg['online\_id\_enabled']) の結果を代入する。
\item 条件 'stop\_points\_yaml' in cfg を判定し、真なら内部処理を行う。
\item   self.\_load\_stops(...) を実行する。
            \end{enumerate}
            \paragraph{代表コード断片}
            \begin{lstlisting}[language=Python]
    def _apply_config(self, cfg: dict):
        params = []
        for key in ['v_min_kmh', 'v_max_kmh', 'v_ref_kmh', 'w_fuel', 'w_speed', 'w_dv', 'w_stop',
                    'dv_max_kmhps']:
            if key in cfg:
                params.append(Parameter(key, value=cfg[key]))
        if 'dv_max_kmh_per_s' in cfg and 'dv_max_kmhps' not in cfg:
            params.append(Parameter('dv_max_kmhps', value=cfg['dv_max_kmh_per_s']))
        if 'horizon_steps' in cfg:
            params.append(Parameter('passo_horizon_steps', value=cfg['horizon_steps']))
        if 'dt_control' in cfg:
            params.append(Parameter('passo_dt', value=cfg['dt_control']))
        if params:
            self.set_parameters(params)
        if 'horizon_steps' in cfg:
            self.Np = int(cfg['horizon_steps'])
        if 'dt_control' in cfg:
            self.dt = float(cfg['dt_control'])
        if 'model_a0' in cfg:
            self.model_coeffs[0] = float(cfg['model_a0'])
        if 'model_a1' in cfg:
            self.model_coeffs[1] = float(cfg['model_a1'])
        if 'model_a2' in cfg:
            self.model_coeffs[2] = float(cfg['model_a2'])
        if 'model_a3' in cfg:
            self.model_coeffs[3] = float(cfg['model_a3'])
        if 'model_a4' in cfg:
            self.model_coeffs[4] = float(cfg['model_a4'])
        if 'online_id_enabled' in cfg:
            self.online_id_enabled = bool(cfg['online_id_enabled'])
        if 'stop_points_yaml' in cfg:
            self._load_stops(str(cfg['stop_points_yaml']))
\end{lstlisting}

\subsection{L2778 関数 MPCNode.\_stop\_penalty\_passo}
            \begin{itemize}[leftmargin=1.6em]
              \item 行範囲: L2778--L2790
              \item 所属: \path|MPCNode|
              \item 定義: \path|_stop_penalty_passo(self, s_km: float, v_kmh: float) -> float|
              \item docstring: 明示されていない。
              \item このブロックが直接呼ぶ主な関数・メソッド: abs, float, get, get\_parameter, max
              \item 戻り値の要点: pen
              \item 主なローカル代入名: dwell\_s, pen, s\_stop, st, v\_ms, vmax\_kmh, vmax\_ms, w\_stop, width\_km
              \item 読み取る主なインスタンス属性: self.get\_parameter, self.stops
              \item 更新する主なインスタンス属性: 特になし。
              \item 明示的に送出する例外: 明示的raiseはない。
              \item 制御構造の規模: 条件分岐 1、ループ 1、try 0
            \end{itemize}
            \paragraph{この定義を読むためのPython構文}
            \begin{itemize}[leftmargin=1.6em]
            \item 第1引数selfは、このメソッドを呼んだインスタンス自身を受け取る慣習名である。
\item `->`以後は戻り値の型注釈であり、実行時の値を自動変換する命令ではない。
            \end{itemize}
            \paragraph{上から順の処理}
            \begin{enumerate}[leftmargin=1.8em]
            \item v\_ms に v\_kmh / 3.6 の結果を代入する。
\item vmax\_kmh に float(self.get\_parameter('v\_max\_kmh').value) の結果を代入する。
\item vmax\_ms に vmax\_kmh / 3.6 の結果を代入する。
\item pen に 0.0 の結果を代入する。
\item w\_stop に float(self.get\_parameter('w\_stop').value) の結果を代入する。
\item self.stops を順に走査し、各要素を st に入れて処理する。
\item   s\_stop に float(st.get('s\_km', 0.0)) の結果を代入する。
\item   dwell\_s に float(st.get('dwell\_s', 0.0)) の結果を代入する。
\item   width\_km に max(0.05, dwell\_s * vmax\_ms / 1000.0 * 0.5) の結果を代入する。
\item   条件 abs(s\_km - s\_stop) <= width\_km を判定し、真なら内部処理を行う。
\item     pen を Add で更新する。
\item pen を返す。
            \end{enumerate}
            \paragraph{代表コード断片}
            \begin{lstlisting}[language=Python]
    def _stop_penalty_passo(self, s_km: float, v_kmh: float) -> float:
        v_ms = v_kmh / 3.6
        vmax_kmh = float(self.get_parameter('v_max_kmh').value)
        vmax_ms = vmax_kmh / 3.6
        pen = 0.0
        w_stop = float(self.get_parameter('w_stop').value)
        for st in self.stops:
            s_stop = float(st.get('s_km', 0.0))
            dwell_s = float(st.get('dwell_s', 0.0))
            width_km = max(0.05, (dwell_s * vmax_ms) / 1000.0 * 0.5)
            if abs(s_km - s_stop) <= width_km:
                pen += w_stop * (v_ms ** 2)
        return pen
\end{lstlisting}

\subsection{L2792 関数 MPCNode.\_fuel\_model\_lph}
            \begin{itemize}[leftmargin=1.6em]
              \item 行範囲: L2792--L2798
              \item 所属: \path|MPCNode|
              \item 定義: \path|_fuel_model_lph(self, v_kmh: float, acc_kmhps: float, grade: float) -> float|
              \item docstring: 明示されていない。
              \item このブロックが直接呼ぶ主な関数・メソッド: float, isfinite, max
              \item 戻り値の要点: max(0.0, float(fuel))
              \item 主なローカル代入名: a0, a0\_eff, a1, a2, a3, a4, fuel
              \item 読み取る主なインスタンス属性: self.idle\_fuel\_lph, self.model\_coeffs
              \item 更新する主なインスタンス属性: 特になし。
              \item 明示的に送出する例外: 明示的raiseはない。
              \item 制御構造の規模: 条件分岐 1、ループ 0、try 0
            \end{itemize}
            \paragraph{この定義を読むためのPython構文}
            \begin{itemize}[leftmargin=1.6em]
            \item 第1引数selfは、このメソッドを呼んだインスタンス自身を受け取る慣習名である。
\item `->`以後は戻り値の型注釈であり、実行時の値を自動変換する命令ではない。
            \end{itemize}
            \paragraph{上から順の処理}
            \begin{enumerate}[leftmargin=1.8em]
            \item (a0, a1, a2, a3, a4) に self.model\_coeffs の結果を代入する。
\item a0\_eff に a0 の結果を代入する。
\item 条件 np.isfinite(self.idle\_fuel\_lph) を判定し、真なら内部処理を行う。
\item   a0\_eff に float(self.idle\_fuel\_lph) の結果を代入する。
\item fuel に a0\_eff + a1 * v\_kmh + a2 * v\_kmh ** 2 + a3 * acc\_kmhps ** 2 + a4 * grade * v\_kmh の結果を代入する。
\item max(0.0, float(fuel)) を返す。
            \end{enumerate}
            \paragraph{代表コード断片}
            \begin{lstlisting}[language=Python]
    def _fuel_model_lph(self, v_kmh: float, acc_kmhps: float, grade: float) -> float:
        a0, a1, a2, a3, a4 = self.model_coeffs
        a0_eff = a0
        if np.isfinite(self.idle_fuel_lph):
            a0_eff = float(self.idle_fuel_lph)
        fuel = a0_eff + a1 * v_kmh + a2 * (v_kmh ** 2) + a3 * (acc_kmhps ** 2) + a4 * grade * v_kmh
        return max(0.0, float(fuel))
\end{lstlisting}

\subsection{L2800 関数 MPCNode.\_update\_identification}
            \begin{itemize}[leftmargin=1.6em]
              \item 行範囲: L2800--L2834
              \item 所属: \path|MPCNode|
              \item 定義: \path|_update_identification(self, now_sec: float, acc_kmhps: float)|
              \item docstring: 明示されていない。
              \item このブロックが直接呼ぶ主な関数・メソッド: append, array, column\_stack, float, isfinite, len, list, lstsq, maximum, mean, ones\_like, popleft
              \item 戻り値の要点: 戻り値が無いか、return 文が明示されていない。
              \item 主なローカル代入名: X, \_, a, alpha, coeffs, g, grade, r2, resid, rmse, s, samples, ss\_res, ss\_tot, v, y, y\_hat
              \item 読み取る主なインスタンス属性: self.fuel\_rate\_lph, self.grade, self.id\_ema\_alpha, self.id\_min\_samples, self.id\_samples, self.id\_window\_sec, self.idle\_fuel\_lph, self.model\_coeffs, self.obd\_ok, self.online\_id\_enabled, self.v\_now
              \item 更新する主なインスタンス属性: self.id\_r2, self.id\_rmse, self.model\_coeffs
              \item 明示的に送出する例外: 明示的raiseはない。
              \item 制御構造の規模: 条件分岐 6、ループ 1、try 0
            \end{itemize}
            \paragraph{この定義を読むためのPython構文}
            \begin{itemize}[leftmargin=1.6em]
            \item 第1引数selfは、このメソッドを呼んだインスタンス自身を受け取る慣習名である。
\item 内包表記は、反復・条件・値生成を一つの式で記述してcollectionを作る。
            \end{itemize}
            \paragraph{上から順の処理}
            \begin{enumerate}[leftmargin=1.8em]
            \item 条件 not self.online\_id\_enabled を判定し、真なら内部処理を行う。
\item    を返す。
\item 条件 self.obd\_ok < 0.5 を判定し、真なら内部処理を行う。
\item    を返す。
\item 条件 not np.isfinite(self.v\_now) or not np.isfinite(self.fuel\_rate\_lph) を判定し、真なら内部処理を行う。
\item    を返す。
\item 条件 self.v\_now <= 5.0 を判定し、真なら内部処理を行う。
\item    を返す。
\item grade に float(self.grade) if np.isfinite(self.grade) else 0.0 の結果を代入する。
\item self.id\_samples.append(...) を実行する。
\item 条件 self.id\_samples and now\_sec - self.id\_samples[0][0] > self.id\_window\_sec が成り立つ間くり返す。
\item   self.id\_samples.popleft(...) を実行する。
\item 条件 len(self.id\_samples) < self.id\_min\_samples を判定し、真なら内部処理を行う。
\item    を返す。
\item samples に list(self.id\_samples) の結果を代入する。
\item v に np.array([s[1] for s in samples], dtype=float) の結果を代入する。
\item a に np.array([s[2] for s in samples], dtype=float) の結果を代入する。
\item g に np.array([s[3] for s in samples], dtype=float) の結果を代入する。
\item y に np.array([s[4] for s in samples], dtype=float) の結果を代入する。
\item X に np.column\_stack([np.ones\_like(v), v, v ** 2, a ** 2, g * v]) の結果を代入する。
\item (coeffs, \_, \_, \_) に np.linalg.lstsq(X, y, rcond=None) の結果を代入する。
\item coeffs に np.maximum(coeffs, 0.0) の結果を代入する。
\item alpha に self.id\_ema\_alpha の結果を代入する。
\item self.model\_coeffs に (1.0 - alpha) * self.model\_coeffs + alpha * coeffs の結果を代入する。
\item 条件 np.isfinite(self.idle\_fuel\_lph) を判定し、真なら内部処理を行う。
\item   self.model\_coeffs[0] に float(self.idle\_fuel\_lph) の結果を代入する。
\item y\_hat に X @ coeffs の結果を代入する。
\item resid に y - y\_hat の結果を代入する。
\item rmse に float(np.sqrt(np.mean(resid ** 2))) if len(resid) else math.nan の結果を代入する。
\item ss\_tot に float(np.sum((y - np.mean(y)) ** 2)) if len(y) else 0.0 の結果を代入する。
\item ss\_res に float(np.sum(resid ** 2)) の結果を代入する。
\item r2 に 1.0 - ss\_res / ss\_tot if ss\_tot > 0.0 else math.nan の結果を代入する。
\item self.id\_rmse に rmse の結果を代入する。
\item self.id\_r2 に r2 の結果を代入する。
            \end{enumerate}
            \paragraph{代表コード断片}
            \begin{lstlisting}[language=Python]
    def _update_identification(self, now_sec: float, acc_kmhps: float):
        if not self.online_id_enabled:
            return
        if self.obd_ok < 0.5:
            return
        if not np.isfinite(self.v_now) or not np.isfinite(self.fuel_rate_lph):
            return
        if self.v_now <= 5.0:
            return
        grade = float(self.grade) if np.isfinite(self.grade) else 0.0
        self.id_samples.append((now_sec, float(self.v_now), float(acc_kmhps), grade, float(self.fuel_rate_lph)))
        while self.id_samples and (now_sec - self.id_samples[0][0]) > self.id_window_sec:
            self.id_samples.popleft()
        if len(self.id_samples) < self.id_min_samples:
            return
        samples = list(self.id_samples)
        v = np.array([s[1] for s in samples], dtype=float)
        a = np.array([s[2] for s in samples], dtype=float)
        g = np.array([s[3] for s in samples], dtype=float)
        y = np.array([s[4] for s in samples], dtype=float)
        X = np.column_stack([np.ones_like(v), v, v ** 2, a ** 2, g * v])
        coeffs, _, _, _ = np.linalg.lstsq(X, y, rcond=None)
        coeffs = np.maximum(coeffs, 0.0)
        alpha = self.id_ema_alpha
        self.model_coeffs = (1.0 - alpha) * self.model_coeffs + alpha * coeffs
        if np.isfinite(self.idle_fuel_lph):
            self.model_coeffs[0] = float(self.idle_fuel_lph)
        y_hat = X @ coeffs
        resid = y - y_hat
        rmse = float(np.sqrt(np.mean(resid ** 2))) if len(resid) else math.nan
        ss_tot = float(np.sum((y - np.mean(y)) ** 2)) if len(y) else 0.0
        ss_res = float(np.sum(resid ** 2))
        r2 = 1.0 - ss_res / ss_tot if ss_tot > 0.0 else math.nan
        self.id_rmse = rmse
        self.id_r2 = r2
\end{lstlisting}

\subsection{L2836 関数 MPCNode.\_solve\_passo\_mpc}
            \begin{itemize}[leftmargin=1.6em]
              \item 行範囲: L2836--L2877
              \item 所属: \path|MPCNode|
              \item 定義: \path|_solve_passo_mpc(self, w_fuel_override = None) -> np.ndarray|
              \item docstring: 明示されていない。
              \item このブロックが直接呼ぶ主な関数・メソッド: \_fuel\_model\_lph, \_stop\_penalty\_passo, abs, all, array, clip, dict, float, get\_parameter, isfinite, max, minimize
              \item 戻り値の要点: np.array(res.x, dtype=float) / J / x0
              \item 主なローカル代入名: J, Np, bounds, dv, dv\_max, fuel\_l, fuel\_lph, grade, k, res, s\_km, v0, v\_k, v\_max, v\_min, v\_prev, v\_ref, w\_dv, w\_dv\_limit, w\_fuel
              \item 読み取る主なインスタンス属性: self.Np, self.\_fuel\_model\_lph, self.\_stop\_penalty\_passo, self.dt, self.get\_parameter, self.grade, self.s\_km, self.v\_now
              \item 更新する主なインスタンス属性: 特になし。
              \item 明示的に送出する例外: 明示的raiseはない。
              \item 制御構造の規模: 条件分岐 3、ループ 1、try 0
            \end{itemize}
            \paragraph{この定義を読むためのPython構文}
            \begin{itemize}[leftmargin=1.6em]
            \item 第1引数selfは、このメソッドを呼んだインスタンス自身を受け取る慣習名である。
\item `->`以後は戻り値の型注釈であり、実行時の値を自動変換する命令ではない。
            \end{itemize}
            \paragraph{上から順の処理}
            \begin{enumerate}[leftmargin=1.8em]
            \item v\_min に float(self.get\_parameter('v\_min\_kmh').value) の結果を代入する。
\item v\_max に float(self.get\_parameter('v\_max\_kmh').value) の結果を代入する。
\item v\_ref に float(self.get\_parameter('v\_ref\_kmh').value) の結果を代入する。
\item w\_fuel に float(self.get\_parameter('w\_fuel').value) の結果を代入する。
\item 条件 w\_fuel\_override is not None を判定し、真なら内部処理を行う。
\item   w\_fuel に float(w\_fuel\_override) の結果を代入する。
\item w\_speed に float(self.get\_parameter('w\_speed').value) の結果を代入する。
\item w\_dv に float(self.get\_parameter('w\_dv').value) の結果を代入する。
\item w\_dv\_limit に float(self.get\_parameter('w\_dv\_limit').value) の結果を代入する。
\item dv\_max に float(self.get\_parameter('dv\_max\_kmhps').value) の結果を代入する。
\item Np に max(1, self.Np) の結果を代入する。
\item v0 に v\_ref if not np.isfinite(self.v\_now) else float(np.clip(self.v\_now, v\_min, v\_max)) の結果を代入する。
\item x0 に np.array([v0] * Np, dtype=float) の結果を代入する。
\item bounds に [(v\_min, v\_max)] * Np の結果を代入する。
\item grade に float(self.grade) if np.isfinite(self.grade) else 0.0 の結果を代入する。
\item 関数 cost を定義する。
\item res に minimize(cost, x0, method='L-BFGS-B', bounds=bounds, options=dict(maxiter=120)) の結果を代入する。
\item 条件 not np.all(np.isfinite(res.x)) を判定し、真なら内部処理を行う。
\item   x0 を返す。
\item np.array(res.x, dtype=float) を返す。
            \end{enumerate}
            \paragraph{代表コード断片}
            \begin{lstlisting}[language=Python]
    def _solve_passo_mpc(self, w_fuel_override=None) -> np.ndarray:
        v_min = float(self.get_parameter('v_min_kmh').value)
        v_max = float(self.get_parameter('v_max_kmh').value)
        v_ref = float(self.get_parameter('v_ref_kmh').value)
        w_fuel = float(self.get_parameter('w_fuel').value)
        if w_fuel_override is not None:
            w_fuel = float(w_fuel_override)
        w_speed = float(self.get_parameter('w_speed').value)
        w_dv = float(self.get_parameter('w_dv').value)
        w_dv_limit = float(self.get_parameter('w_dv_limit').value)
        dv_max = float(self.get_parameter('dv_max_kmhps').value)

        Np = max(1, self.Np)
        v0 = v_ref if not np.isfinite(self.v_now) else float(np.clip(self.v_now, v_min, v_max))
        x0 = np.array([v0] * Np, dtype=float)
        bounds = [(v_min, v_max)] * Np

        grade = float(self.grade) if np.isfinite(self.grade) else 0.0

        def cost(v_vec):
            s_km = float(self.s_km)
            v_prev = v0
            J = 0.0
            for k in range(Np):
                v_k = float(v_vec[k])
                dv = (v_k - v_prev) / max(self.dt, 1.0e-3)
                fuel_lph = self._fuel_model_lph(v_k, dv, grade)
                fuel_l = fuel_lph * (self.dt / 3600.0)
                J += w_fuel * fuel_l
                J += w_speed * (v_k - v_ref) ** 2
                J += w_dv * (dv ** 2)
                if dv_max > 0.0:
                    J += w_dv_limit * max(0.0, abs(dv) - dv_max) ** 2
                s_km += v_k * (self.dt / 3600.0)
                J += self._stop_penalty_passo(s_km, v_k)
...
\end{lstlisting}

\subsection{L2855 関数 MPCNode.\_solve\_passo\_mpc.cost}
            \begin{itemize}[leftmargin=1.6em]
              \item 行範囲: L2855--L2872
              \item 所属: \path|MPCNode._solve_passo_mpc|
              \item 定義: \path|cost(v_vec)|
              \item docstring: 明示されていない。
              \item このブロックが直接呼ぶ主な関数・メソッド: \_fuel\_model\_lph, \_stop\_penalty\_passo, abs, float, max, range
              \item 戻り値の要点: J
              \item 主なローカル代入名: J, dv, fuel\_l, fuel\_lph, k, s\_km, v\_k, v\_prev
              \item 読み取る主なインスタンス属性: self.\_fuel\_model\_lph, self.\_stop\_penalty\_passo, self.dt, self.s\_km
              \item 更新する主なインスタンス属性: 特になし。
              \item 明示的に送出する例外: 明示的raiseはない。
              \item 制御構造の規模: 条件分岐 1、ループ 1、try 0
            \end{itemize}
            \paragraph{この定義を読むためのPython構文}
            \begin{itemize}[leftmargin=1.6em]
            \item 特別な構文上の補足は少ない。
            \end{itemize}
            \paragraph{上から順の処理}
            \begin{enumerate}[leftmargin=1.8em]
            \item s\_km に float(self.s\_km) の結果を代入する。
\item v\_prev に v0 の結果を代入する。
\item J に 0.0 の結果を代入する。
\item range(Np) を順に走査し、各要素を k に入れて処理する。
\item   v\_k に float(v\_vec[k]) の結果を代入する。
\item   dv に (v\_k - v\_prev) / max(self.dt, 0.001) の結果を代入する。
\item   fuel\_lph に self.\_fuel\_model\_lph(v\_k, dv, grade) の結果を代入する。
\item   fuel\_l に fuel\_lph * (self.dt / 3600.0) の結果を代入する。
\item   J を Add で更新する。
\item   J を Add で更新する。
\item   J を Add で更新する。
\item   条件 dv\_max > 0.0 を判定し、真なら内部処理を行う。
\item     J を Add で更新する。
\item   s\_km を Add で更新する。
\item   J を Add で更新する。
\item   v\_prev に v\_k の結果を代入する。
\item J を返す。
            \end{enumerate}
            \paragraph{代表コード断片}
            \begin{lstlisting}[language=Python]
        def cost(v_vec):
            s_km = float(self.s_km)
            v_prev = v0
            J = 0.0
            for k in range(Np):
                v_k = float(v_vec[k])
                dv = (v_k - v_prev) / max(self.dt, 1.0e-3)
                fuel_lph = self._fuel_model_lph(v_k, dv, grade)
                fuel_l = fuel_lph * (self.dt / 3600.0)
                J += w_fuel * fuel_l
                J += w_speed * (v_k - v_ref) ** 2
                J += w_dv * (dv ** 2)
                if dv_max > 0.0:
                    J += w_dv_limit * max(0.0, abs(dv) - dv_max) ** 2
                s_km += v_k * (self.dt / 3600.0)
                J += self._stop_penalty_passo(s_km, v_k)
                v_prev = v_k
            return J
\end{lstlisting}

\subsection{L2879 関数 MPCNode.\_publish\_trajectory\_passo}
            \begin{itemize}[leftmargin=1.6em]
              \item 行範囲: L2879--L2892
              \item 所属: \path|MPCNode|
              \item 定義: \path|_publish_trajectory_passo(self, v_seq: np.ndarray)|
              \item docstring: 明示されていない。
              \item このブロックが直接呼ぶ主な関数・メソッド: Path, PoseStamped, append, enumerate, float, get\_clock, now, publish, to\_msg
              \item 戻り値の要点: 戻り値が無いか、return 文が明示されていない。
              \item 主なローカル代入名: k, path, pose, s\_tmp, v\_k
              \item 読み取る主なインスタンス属性: self.dt, self.get\_clock, self.pub\_path, self.s\_km
              \item 更新する主なインスタンス属性: 特になし。
              \item 明示的に送出する例外: 明示的raiseはない。
              \item 制御構造の規模: 条件分岐 0、ループ 1、try 0
            \end{itemize}
            \paragraph{この定義を読むためのPython構文}
            \begin{itemize}[leftmargin=1.6em]
            \item 第1引数selfは、このメソッドを呼んだインスタンス自身を受け取る慣習名である。
            \end{itemize}
            \paragraph{上から順の処理}
            \begin{enumerate}[leftmargin=1.8em]
            \item path に Path() の結果を代入する。
\item path.header.stamp に self.get\_clock().now().to\_msg() の結果を代入する。
\item path.header.frame\_id に 'map' の結果を代入する。
\item s\_tmp に float(self.s\_km) の結果を代入する。
\item enumerate(v\_seq) を順に走査し、各要素を (k, v\_k) に入れて処理する。
\item   s\_tmp を Add で更新する。
\item   pose に PoseStamped() の結果を代入する。
\item   pose.header に path.header の結果を代入する。
\item   pose.pose.position.x に float(k) の結果を代入する。
\item   pose.pose.position.y に float(v\_k) の結果を代入する。
\item   pose.pose.position.z に s\_tmp の結果を代入する。
\item   path.poses.append(...) を実行する。
\item self.pub\_path.publish(...) を実行する。
            \end{enumerate}
            \paragraph{代表コード断片}
            \begin{lstlisting}[language=Python]
    def _publish_trajectory_passo(self, v_seq: np.ndarray):
        path = Path()
        path.header.stamp = self.get_clock().now().to_msg()
        path.header.frame_id = 'map'
        s_tmp = float(self.s_km)
        for k, v_k in enumerate(v_seq):
            s_tmp += float(v_k) * (self.dt / 3600.0)
            pose = PoseStamped()
            pose.header = path.header
            pose.pose.position.x = float(k)
            pose.pose.position.y = float(v_k)
            pose.pose.position.z = s_tmp
            path.poses.append(pose)
        self.pub_path.publish(path)
\end{lstlisting}

\subsection{L2894 関数 MPCNode.\_step\_passo}
            \begin{itemize}[leftmargin=1.6em]
              \item 行範囲: L2894--L2962
              \item 所属: \path|MPCNode|
              \item 定義: \path|_step_passo(self)|
              \item docstring: 明示されていない。
              \item このブロックが直接呼ぶ主な関数・メソッド: Float32, Float32MultiArray, String, \_fuel\_model\_lph, \_publish\_trajectory\_passo, \_solve\_passo\_mpc, \_update\_identification, array, float, get\_clock, get\_parameter, isfinite
              \item 戻り値の要点: 戻り値が無いか、return 文が明示されていない。
              \item 主なローカル代入名: acc\_kmhps, dt, dt\_step, fuel\_pred\_lph, fuel\_valid, grade, msg, now\_sec, obd\_valid, run\_ready, run\_ready\_sec, speed\_valid, status, v\_ref, v\_seq
              \item 読み取る主なインスタンス属性: self.Np, self.\_fuel\_model\_lph, self.\_publish\_trajectory\_passo, self.\_solve\_passo\_mpc, self.\_update\_identification, self.config\_ready, self.fuel\_rate\_lph, self.get\_clock, self.get\_parameter, self.grade, self.id\_r2, self.id\_rmse, self.last\_step\_time, self.max\_acc\_dt\_sec, self.mpc\_state, self.obd\_ok, self.prev\_speed\_kmh, self.prev\_speed\_time, self.pub\_mpc\_state, self.pub\_speed, self.pub\_status, self.s\_km, self.v\_cmd, self.v\_now
              \item 更新する主なインスタンス属性: self.last\_step\_time, self.mpc\_state, self.prev\_speed\_kmh, self.prev\_speed\_time, self.v\_cmd, self.valid\_fuel\_sec, self.valid\_obd\_sec, self.valid\_speed\_sec
              \item 明示的に送出する例外: 明示的raiseはない。
              \item 制御構造の規模: 条件分岐 8、ループ 0、try 0
            \end{itemize}
            \paragraph{この定義を読むためのPython構文}
            \begin{itemize}[leftmargin=1.6em]
            \item 第1引数selfは、このメソッドを呼んだインスタンス自身を受け取る慣習名である。
            \end{itemize}
            \paragraph{上から順の処理}
            \begin{enumerate}[leftmargin=1.8em]
            \item now\_sec に self.get\_clock().now().nanoseconds / 1000000000.0 の結果を代入する。
\item 条件 self.last\_step\_time is None を判定し、真なら内部処理を行う。
\item   dt\_step に 0.0 の結果を代入する。
\item   上の条件が偽の場合:
\item   dt\_step に max(0.0, min(5.0, now\_sec - self.last\_step\_time)) の結果を代入する。
\item self.last\_step\_time に now\_sec の結果を代入する。
\item acc\_kmhps に 0.0 の結果を代入する。
\item 条件 np.isfinite(self.v\_now) を判定し、真なら内部処理を行う。
\item   条件 self.prev\_speed\_time is not None を判定し、真なら内部処理を行う。
\item     dt に now\_sec - self.prev\_speed\_time の結果を代入する。
\item     条件 0.0 < dt <= self.max\_acc\_dt\_sec and np.isfinite(self.prev\_speed\_kmh) を判定し、真なら内部処理を行う。
\item       acc\_kmhps に (float(self.v\_now) - float(self.prev\_speed\_kmh)) / dt の結果を代入する。
\item   self.prev\_speed\_time に now\_sec の結果を代入する。
\item   self.prev\_speed\_kmh に float(self.v\_now) の結果を代入する。
\item v\_ref に float(self.get\_parameter('v\_ref\_kmh').value) の結果を代入する。
\item speed\_valid に np.isfinite(self.v\_now) の結果を代入する。
\item fuel\_valid に np.isfinite(self.fuel\_rate\_lph) の結果を代入する。
\item obd\_valid に self.obd\_ok > 0.5 の結果を代入する。
\item 条件 dt\_step > 0.0 を判定し、真なら内部処理を行う。
\item   self.valid\_obd\_sec に self.valid\_obd\_sec + dt\_step if obd\_valid else 0.0 の結果を代入する。
\item   self.valid\_speed\_sec に self.valid\_speed\_sec + dt\_step if speed\_valid else 0.0 の結果を代入する。
\item   self.valid\_fuel\_sec に self.valid\_fuel\_sec + dt\_step if fuel\_valid else 0.0 の結果を代入する。
\item run\_ready\_sec に float(self.get\_parameter('run\_ready\_sec').value) の結果を代入する。
\item run\_ready に self.valid\_obd\_sec >= run\_ready\_sec and self.valid\_speed\_sec >= run\_ready\_sec and (self.valid\_fuel\_sec >= run\_ready\_sec) の結果を代入する。
\item 条件 not self.config\_ready を判定し、真なら内部処理を行う。
\item   self.mpc\_state に 'IDLE' の結果を代入する。
\item   self.v\_cmd に v\_ref の結果を代入する。
\item   v\_seq に np.array([self.v\_cmd] * max(1, self.Np), dtype=float) の結果を代入する。
\item   上の条件が偽の場合:
\item   条件 not run\_ready を判定し、真なら内部処理を行う。
\item     self.mpc\_state に 'DEGRADED\_RUN' の結果を代入する。
\item     self.v\_cmd に float(self.v\_now) if speed\_valid else v\_ref の結果を代入する。
\item     v\_seq に np.array([self.v\_cmd] * max(1, self.Np), dtype=float) の結果を代入する。
\item     上の条件が偽の場合:
\item     self.mpc\_state に 'RUN' の結果を代入する。
\item     self.\_update\_identification(...) を実行する。
\item     v\_seq に self.\_solve\_passo\_mpc() の結果を代入する。
\item     self.v\_cmd に float(v\_seq[0]) の結果を代入する。
\item msg に Float32() の結果を代入する。
\item msg.data に float(self.v\_cmd) の結果を代入する。
\item self.pub\_speed.publish(...) を実行する。
\item self.\_publish\_trajectory\_passo(...) を実行する。
\item fuel\_pred\_lph に math.nan の結果を代入する。
\item 条件 np.isfinite(self.v\_cmd) を判定し、真なら内部処理を行う。
\item   grade に float(self.grade) if np.isfinite(self.grade) else 0.0 の結果を代入する。
\item   fuel\_pred\_lph に self.\_fuel\_model\_lph(float(self.v\_cmd), 0.0, grade) の結果を代入する。
\item status に Float32MultiArray() の結果を代入する。
\item status.data に [float(self.fuel\_rate\_lph) if np.isfinite(self.fuel\_rate\_lph) else math.nan, float(fuel\_pred\_lph) if np.isfinite(fuel\_pred\_lph) else math.nan, float(self.v\_now) if np.isfinite(self.v\_now) else math.nan, float(self.v\_cmd), float(self.s\_km), float(self.id\_rmse) if np.isfinite(self.id\_rmse) else math.nan, float(self.id\_r2) if np.isfinite(self.id\_r2) else math.nan] の結果を代入する。
\item self.pub\_status.publish(...) を実行する。
\item self.pub\_mpc\_state.publish(...) を実行する。
            \end{enumerate}
            \paragraph{代表コード断片}
            \begin{lstlisting}[language=Python]
    def _step_passo(self):
        now_sec = self.get_clock().now().nanoseconds / 1e9
        if self.last_step_time is None:
            dt_step = 0.0
        else:
            dt_step = max(0.0, min(5.0, now_sec - self.last_step_time))
        self.last_step_time = now_sec

        acc_kmhps = 0.0
        if np.isfinite(self.v_now):
            if self.prev_speed_time is not None:
                dt = now_sec - self.prev_speed_time
                if 0.0 < dt <= self.max_acc_dt_sec and np.isfinite(self.prev_speed_kmh):
                    acc_kmhps = (float(self.v_now) - float(self.prev_speed_kmh)) / dt
            self.prev_speed_time = now_sec
            self.prev_speed_kmh = float(self.v_now)

        v_ref = float(self.get_parameter('v_ref_kmh').value)
        speed_valid = np.isfinite(self.v_now)
        fuel_valid = np.isfinite(self.fuel_rate_lph)
        obd_valid = self.obd_ok > 0.5

        if dt_step > 0.0:
            self.valid_obd_sec = self.valid_obd_sec + dt_step if obd_valid else 0.0
            self.valid_speed_sec = self.valid_speed_sec + dt_step if speed_valid else 0.0
            self.valid_fuel_sec = self.valid_fuel_sec + dt_step if fuel_valid else 0.0

        run_ready_sec = float(self.get_parameter('run_ready_sec').value)
        run_ready = (self.valid_obd_sec >= run_ready_sec and
                     self.valid_speed_sec >= run_ready_sec and
                     self.valid_fuel_sec >= run_ready_sec)

        if not self.config_ready:
            self.mpc_state = 'IDLE'
            self.v_cmd = v_ref
...
\end{lstlisting}

\subsection{L2965 関数 main}
            \begin{itemize}[leftmargin=1.6em]
              \item 行範囲: L2965--L2975
              \item 所属: module直下
              \item 定義: \path|main()|
              \item docstring: 明示されていない。
              \item このブロックが直接呼ぶ主な関数・メソッド: MPCNode, MultiThreadedExecutor, add\_node, destroy\_node, init, shutdown, spin
              \item 戻り値の要点: 戻り値が無いか、return 文が明示されていない。
              \item 主なローカル代入名: executor, node
              \item 読み取る主なインスタンス属性: 特になし。
              \item 更新する主なインスタンス属性: 特になし。
              \item 明示的に送出する例外: 明示的raiseはない。
              \item 制御構造の規模: 条件分岐 0、ループ 0、try 1
            \end{itemize}
            \paragraph{この定義を読むためのPython構文}
            \begin{itemize}[leftmargin=1.6em]
            \item try文は例外が起き得る処理と、異常時または終了時の経路を分ける。
            \end{itemize}
            \paragraph{上から順の処理}
            \begin{enumerate}[leftmargin=1.8em]
            \item rclpy.init(...) を実行する。
\item node に MPCNode() の結果を代入する。
\item executor に MultiThreadedExecutor(num\_threads=4) の結果を代入する。
\item executor.add\_node(...) を実行する。
\item 例外処理を伴う try ブロックを実行する。
\item   executor.spin(...) を実行する。
\item   成否にかかわらずfinallyで:
\item   executor.shutdown(...) を実行する。
\item   node.destroy\_node(...) を実行する。
\item   rclpy.shutdown(...) を実行する。
            \end{enumerate}
            \paragraph{代表コード断片}
            \begin{lstlisting}[language=Python]
def main():
    rclpy.init()
    node = MPCNode()
    executor = MultiThreadedExecutor(num_threads=4)
    executor.add_node(node)
    try:
        executor.spin()
    finally:
        executor.shutdown()
        node.destroy_node()
        rclpy.shutdown()
\end{lstlisting}

        \subsection{設定値と入力パラメータ}
            \begin{longtable}{p{0.07\linewidth} p{0.35\linewidth} p{0.45\linewidth}}
            \toprule
            行 & 名前 & 既定値・補足 \\
            \midrule
            58 & \path|passo_mode| & False \\
292 & \path|forecast_csv| & inputs/forecast\_10min.csv \\
293 & \path|maps_dir| & maps \\
294 & \path|drive_eff_map| &  \\
295 & \path|regen_eff_map| &  \\
296 & \path|rint_map| &  \\
297 & \path|drive_map_eco| &  \\
298 & \path|drive_map_power| &  \\
299 & \path|regen_map_eco| &  \\
300 & \path|regen_map_power| &  \\
301 & \path|panel_eff_map| &  \\
302 & \path|mppt_eff_map| &  \\
303 & \path|ocv_soc_map| &  \\
304 & \path|dt| & 600.0 \\
305 & \path|horizon_steps| & 9 \\
306 & \path|v_max_kmh| & 110.0 \\
307 & \path|terminal_soc_min| & 0.1 \\
308 & \path|stop_yaml| & inputs/stop\_points.yaml \\
309 & \path|forecast_time_mode| & auto \\
310 & \path|forecast_time_tz| & UTC \\
311 & \path|forecast_start_time_utc| &  \\
312 & \path|forecast_time_offset_sec| & 0.0 \\
313 & \path|forecast_reload_sec| & 60.0 \\
314 & \path|replan_on_forecast_reload| & False \\
315 & \path|params_yaml| &  \\
316 & \path|profile_runtime_mode| &  \\
317 & \path|route_profile_csv| &  \\
318 & \path|speed_profile_csv| &  \\
319 & \path|drive_schedule_yaml| &  \\
320 & \path|initial_upper_policy_csv| &  \\
321 & \path|use_measured_s| & True \\
322 & \path|use_measured_speed| & True \\
323 & \path|soc0| & 0.95 \\
324 & \path|Tb0| & 30.0 \\
325 & \path|s0_km| & 0.0 \\
326 & \path|speed_meas_timeout_sec| & 3.0 \\
327 & \path|distance_meas_timeout_sec| & 5.0 \\
328 & \path|battery_meas_timeout_sec| & 15.0 \\
329 & \path|speed_meas_filter_tau_sec| & 0.6 \\
330 & \path|speed_meas_max_accel_kmhps| & 12.0 \\
331 & \path|speed_meas_max_decel_kmhps| & 20.0 \\
332 & \path|distance_meas_max_rate_kmps| & 0.06 \\
333 & \path|distance_meas_max_backtrack_km| & 0.02 \\
334 & \path|battery_meas_filter_tau_sec| & 1.0 \\
335 & \path|w_dv| & 0.05 \\
336 & \path|w_dv_limit| & 2.0 \\
337 & \path|dv_max_kmhps| & 5.0 \\
338 & \path|w_T| & 5.0 \\
339 & \path|w_speed_limit| & 50.0 \\
340 & \path|w_drive_window| & 100000.0 \\
            \bottomrule
            \end{longtable}

        \subsection{ROS topic I/O}
            \begin{longtable}{p{0.07\linewidth} p{0.18\linewidth} p{0.34\linewidth} p{0.28\linewidth}}
            \toprule
            行 & 種別 & topic & callback / 補足 \\
            \midrule
            737 & publisher & \path|/planner/speed_cmd| & publisher \\
738 & publisher & \path|/planner/upper_speed_cmd| & publisher \\
739 & publisher & \path|/planner/throttle_cmd_pct| & publisher \\
740 & publisher & \path|/planner/drive_mode| & publisher \\
741 & publisher & \path|/planner/trajectory| & publisher \\
742 & publisher & \path|/planner/upper_plan| & publisher \\
743 & publisher & \path|/planner/lower_plan| & publisher \\
744 & publisher & \path|/planner/env| & publisher \\
745 & publisher & \path|/planner/metrics| & publisher \\
746 & publisher & \path|/planner/status| & publisher \\
2700 & publisher & \path|/planner/speed_cmd| & publisher \\
2701 & publisher & \path|/planner/trajectory| & publisher \\
2702 & publisher & \path|/planner/status| & publisher \\
2703 & publisher & \path|/system/mpc_state| & publisher \\
749 & subscription & \path|/vehicle/s_km| & self.\_on\_s\_km\_solar \\
750 & subscription & \path|/vehicle/speed_kmh| & self.\_on\_speed\_solar \\
751 & subscription & \path|/vehicle/batt_soc| & self.\_on\_soc\_solar \\
752 & subscription & \path|/vehicle/batt_temp_c| & self.\_on\_tb\_solar \\
753 & subscription & \path|/vehicle/batt_current_a| & self.\_on\_i\_solar \\
754 & subscription & \path|/vehicle/batt_voltage_v| & self.\_on\_v\_solar \\
755 & subscription & \path|/calib/solar_gain| & self.\_on\_calib\_solar\_gain \\
756 & subscription & \path|/calib/drive_power_gain| & self.\_on\_calib\_drive\_power\_gain \\
757 & subscription & \path|/calib/aux_power_w| & self.\_on\_calib\_aux\_power \\
2688 & subscription & \path|/vehicle/s_km| & self.\_on\_s\_km \\
2689 & subscription & \path|/vehicle/speed_kmh| & self.\_on\_speed \\
2690 & subscription & \path|/vehicle/fuel_rate_lph| & self.\_on\_fuel \\
2691 & subscription & \path|/vehicle/throttle_pct| & self.\_on\_throttle \\
2692 & subscription & \path|/vehicle/obd_ok| & self.\_on\_obd\_ok \\
2693 & subscription & \path|/vehicle/grade| & self.\_on\_grade \\
2694 & subscription & \path|/vehicle/idle_fuel_lph| & self.\_on\_idle\_fuel \\
2695 & subscription & \path|/system/config| & self.\_on\_config \\
2696 & subscription & \path|/system/config_ready| & self.\_on\_config\_ready \\
2697 & subscription & \path|/system/state| & self.\_on\_system\_state \\
            \bottomrule
            \end{longtable}







        \section{処理の流れ}
        \begin{enumerate}[leftmargin=1.7em]
        \item 初期化時に設定値や入力パスを読み込む。
\item publisher / subscription / timer を準備する。
\item timer callback 周期で主処理を進める。
        \end{enumerate}

        \section{このファイルの位置づけ}
        この資料群では、\path|mpc_solarcar/mpc_node.py| を
        \textbf{runtime core}の中核として扱う。
        したがって、ここで扱う処理は単独で完結せず、
        前段の profile / launch / telemetry と後段の planner / logger / report へ繋がる。

        \end{document}
